\documentclass[11pt,a4paper]{article}
\usepackage[T1]{fontenc}
\usepackage[utf8]{inputenc}
\usepackage{lmodern,amsmath,amssymb,amsthm,mathtools,microtype,mathrsfs}
\usepackage[margin=25mm,top=23mm,bottom=25mm]{geometry}
\usepackage{booktabs,array,longtable,enumitem,xcolor,fancyhdr}
\usepackage[colorlinks=true,linkcolor=blue!45!black,urlcolor=blue!45!black,citecolor=blue!45!black]{hyperref}
\hypersetup{pdftitle={Reading the Alpoge--Claude RH and Dirichlet-GRH proof: working version 0.43},pdfauthor={ChatGPT},pdfsubject={Partial human-readable reconstruction with explicit arithmetic inputs},pdfkeywords={Riemann hypothesis, Lean, theta lifts, prime sums, Arakelov slopes}}
\setlength{\parindent}{0pt}
\setlength{\parskip}{5pt plus 1pt}
\setlist{topsep=3pt,itemsep=3pt,leftmargin=*}
\setlength{\headheight}{14pt}
\pagestyle{fancy}\fancyhf{}
\fancyhead[L]{\small RH and Dirichlet GRH: a reconstruction}
\fancyhead[R]{\small Working version 0.43}
\fancyfoot[C]{\thepage}
\renewcommand{\headrulewidth}{0.3pt}
\numberwithin{equation}{section}
\newtheorem{theorem}{Theorem}[section]
\newtheorem{proposition}[theorem]{Proposition}
\newtheorem{lemma}[theorem]{Lemma}
\newtheorem{corollary}[theorem]{Corollary}
\theoremstyle{remark}\newtheorem{remark}[theorem]{Remark}
\newcommand{\Rea}{\operatorname{Re}}
\newcommand{\HS}{\mathrm{HS}}
\newcommand{\tr}{\operatorname{tr}}
\newcommand{\codim}{\operatorname{codim}}
\newcommand{\covol}{\operatorname{covol}}
\newcommand{\OK}{\mathcal O_K}
\newcommand{\src}[2]{\href{https://github.com/CocoSgt/riemann-hypothesis-lean/blob/7cbf7a0c32fca1badf05931879a37c3f8a2b03b4/src/#1}{\texttt{#2}}}
\newcommand{\status}[1]{\textbf{Status:} #1}
\begin{document}
\begin{center}
{\LARGE\bfseries RH and Dirichlet GRH: new reconstruction chapters}\\[7pt]
{\large Version 0.43 --- 23 November 2026}\\[5pt]
{\large Integral types, arithmetic fibers, and finite local weights}
\end{center}
\textbf{Scope.} This edition contains the current audit and the five new
chapters. It remains a partial human reconstruction. Earlier named
claims are in the cumulative manuscript and preserved component files.
The underlying formal proof is accepted as correct; the exact remaining
human reconstruction boundaries are stated below.
\setcounter{tocdepth}{1}
\tableofcontents
\clearpage
\begingroup\emergencystretch=3em
\section{Current dependency audit: integral types, arithmetic fibers, and finite local weights}
\label{sec:dependency-audit-v43}

\subsection{Current status and reading order}

Version 0.43 supplies the original integral type and Hecke-data
constructor, the actual compact-unitary automorphic decomposition,
the exact classical and ordinary part decomposition, and the dense
arithmetic-fiber and scalar-separation arguments. It also completes
the finite local estimates used by the selected retained-row weight
consumer. These results are proved on the original objects and at
their stated scopes; their still-required global suppliers are listed
below.

The pinned Alpoge--Claude formal proof is accepted as fully checked
and correct. This remains a \textbf{partial human reconstruction} of
RH and unrestricted Dirichlet GRH. The common arithmetic and geometric
families and the global quantitative producers have not all been
reconstructed. No formal correctness question is reopened, and no
Lean build or kernel replay is claimed.

There are 78 new and 2441 cumulative claims. All 2363 predecessor
claim records, their source arrays and dependency edges are preserved.
The selected sufficient-route graph still has 121 nodes and 161 edges.
No whole selected leaf is closed by this increment: its 31 remaining
whole obligations consist of 22 class-C and 9 class-B tasks. Of these,
24 are shared by RH and Dirichlet GRH, and 7 are additional for Dirichlet
GRH. These counts measure recorded scopes, not a percentage of proof
completion. Progress inside the two affected compound obligations is
substantial and recorded exactly.

Read ITA for the integral type, ADP then CFD for the two decompositions,
AFS for arithmetic points and separation, and FLC for the finite local
weight theorem. The earlier CMP1--CMP6 contain the complete conditional
terminal deduction; they are not repeated or promoted to an unconditional
human proof. The separate reader's edition contains this audit and only
the five new chapters. The cumulative edition preserves the earlier
chapters, whose conditional statements retain their printed hypotheses.

\subsection{The original integral type and Hecke data}

ITA1--ITA13 produce one representation over the original coefficient
DVR $O$, before the later choices of a Teichmuller lift family,
local correspondence or test representation. Uniqueness of the original
Teichmuller lifts handles the first universal quantifier. On the toral
zero-monodromy sector, the written reciprocity and singleton Langlands
quotient calculation handles every weaker source \texttt{Is\allowbreak Rec\allowbreak T} map;
it does not presume global uniqueness for those maps.

The explicitly matched crystalline inertial-type theorem supplies a
complex local detector. The original normalization makes its inducing
norm exponents integers, so arbitrary coefficient automorphisms fixing
$j(O)$ preserve the detected predicate. A stable lattice is constructed
over a finite extension and then restricted as a module to the original
$O$. Every conjugate summand still detects exactly the prescribed type.
The dual representation gives the source's precise inverse-transpose
column equation. Successive finite tensor--Hom constructions establish
the product theorem for every nonzero individually isotypic module,
without adding finite length or admissibility of that module.

This proves \texttt{Type\allowbreak For\allowbreak Alt} at every original allowed prime and
rank, including coefficient prime three, repeated inertia characters,
and an empty set of type places. ITA13 substitutes it into OTW13 and
CAP17, giving the original \texttt{Hd\allowbreak Ex}/\texttt{Unitary\allowbreak Hecke\allowbreak Data}.
Classical local Langlands and the source's \texttt{Rec\allowbreak TExists}
nonemptiness remain explicit inherited foundations; the type theorem
is a matched established input, not a new proof of the entire local
classification.

\subsection{Actual automorphic and classical decompositions}

ADP1--ADP12 work with the original identity-form definite unitary group.
The exact finite class representatives yield a compact fundamental set;
the original quotient measure supplies a positive definite invariant
form on actual continuous functions. Finite multiplicity and the
checked algebraic semisimplicity argument give finite globally
equivariant projections. The explicit scalarization
\[
 \mathcal T_F(\lambda^*)(g)
  =\lambda^*\bigl(R(g^{-1})E(F(p(g)))\bigr)
\]
and its bidual inverse keep the prescribed complex embedding, basis,
infinity representation and rational equivariance. Projecting all
scalarizations and inverting preserves the full matrix-valued
coefficient equations. This proves the original \texttt{form\allowbreak U} and
\texttt{decomp\allowbreak U}, on the actual quotient rather than a substituted
representation space.

CFD1--CFD12 show how a supplied decomposition into the original parts
can be refined without changing the level, depth or scalar characters.
Finite averaging uses left cosets and is valid for nonnormal
stabilizers. Commuting compact-group scalar profiles preserve the
original classical equations. Polynomial generalized-eigenspace
projectors then preserve the exact norm-one ordinary condition for
the original normalized $U_z$. This proves full \texttt{cl\allowbreak Parts\_\allowbreak alt},
including its useful arbitrary-predicate refinement.

The distinction between these endpoints is essential. ADP gives
actual automorphic constituents. The original \texttt{Part\allowbreak Alt\allowbreak Alt}
additionally includes a Galois representation, spherical readings,
local comparisons and arithmetic point properties. Its producer,
\texttt{part\allowbreak X\_\allowbreak alt}/\texttt{part\allowbreak U\_\allowbreak alt}, still has to be reconstructed.
CFD retains the exact \texttt{Automorphic\allowbreak Hecke\allowbreak Set} premise that this
producer will supply.

\subsection{Dense weights, nonflat fibers, scalar action and the same ideal}

AFS1--AFS8 prove the domain and full \texttt{Dense\allowbreak Wt} statements for
every supplied original universal weight ring and every original
dominant weight. Bounded integral power series are separated by the
written infinite Cartesian torsion sets in the open unit disk. The
torsion levels distinguish the finite-part characters, and the exact
stored Artin-map and cyclotomic identities give the literal arithmetic
and good-point predicates. The target remains
$\overline{\mathbf Q}_\ell$, not its completion.

AFS9--AFS19 supply the original fiber and scalar arguments. Every
specified character of the faithful finite operator algebra occurs
in its module fiber, even for a nonflat point map. Every nonzero
finite free ordinary block has nonzero fibers; this does not assert
that an arbitrary residual label gives a nonzero block. The compact
averaging description of the Jacquet kernel and the stated original
local literature clauses prove the distinct-unit principal-series
criterion with the exact scalar normalization. The determinant
argument turns the source's integral eigenvalues into units.

The tensor identification uses the split ordinary idempotent and
the already reconstructed specialization map. Consequently it
transports the actual part decomposition, every Hecke generator,
each specified fiber character and the unique scalar character on
each nonzero good part. This proves full \texttt{part\allowbreak Act0\_\allowbreak alt}.
AFS20 then proves the exact \texttt{TCore\allowbreak Sep}/\texttt{Out.\allowbreak sep}
assertion: a Hecke element vanishing at all good arithmetic points
is zero. The proof expands its action in the original ordinary
basis and applies \texttt{Dense\allowbreak Wt} to every coefficient.

AFS21--AFS22 apply this separation to the original map $f:B\to T$.
For the same ideal $J=\ker f$,
\[
 \bigl[x(b)=0\text{ for every original good point }x
                  \text{ killing }J\bigr]\quad\Longrightarrow\quad b\in J.
\]
No radical, larger support ideal or replacement quotient is introduced.
The zero-module case has $J=B$ and is treated separately. The complete
density implication is written with the exact remaining
\texttt{Automorphic\allowbreak Hecke\allowbreak Set} supplier exposed. Thus the bare original
\texttt{TJDense} installation, the point-output record and the other
same-ideal Hida fields remain part of the current frontier.

\subsection{The selected finite local weight theorem}

FLC1--FLC5 prove the remaining original root-injectivity collision
estimate with its constant $23$, for every finite field of cardinality
at least $2^{12}$ and every original degree. The same-factor estimate
uses the shortened Frobenius orbit after removing the prime-primary
part. In the cross-factor estimate, a common Frobenius translation
places every possible collision in one fixed fiber, avoiding a loss
by the degree of the fixed factor. Weighted divisor bounds and the
explicit numerical series estimates give the stated constant.

FLC6--FLC18 prove the exact finite mirabolic and root-exchange
calculations, the lower $e_J$ comparison and the squarefree $e_U$
weight formula. Green's cuspidal classification is the one explicit
published finite-group input, with its original simplex, sign, degree
and elliptic-value dictionary checked. The comparison to the supplied
normalized-type record preserves the actual series projectors and
factor degrees. Polynomial root predicates are still evaluated on
the original polynomial, and are never transported by an arbitrary
permutation of labels.

FLC19 assembles the actual \texttt{Weight\allowbreak All\allowbreak Rows} assertion for the
selected eventual family, with the unchanged $W_g$ and all five local
error terms. The last term uses RRT10's eventual rank threshold;
the original selected rank growth, not \texttt{Sizes\allowbreak At} alone,
absorbs it uniformly. The stronger fourth \texttt{adm\allowbreak Lq} upper
comparison is unused by this consumer and is not claimed.

The original normalized-type compatibility, adelic reduction and
inflation with its Haar normalization, row identities, common
\texttt{All\allowbreak Rows} object and arithmetic cylinder counts remain actual
inputs to be produced. A local weight theorem therefore does not close
the whole \texttt{GLOBAL\_\allowbreak RETAINED\_\allowbreak DEFECT} obligation.

\subsection{Every remaining selected obligation}

Each contract below is the full selected output. Completed nested
theorems remain available inputs; no compound contract is closed until
its other suppliers, shared objects and choice-order conditions have
been reconstructed. The current notes govern the historical frontier
descriptions retained in the cumulative chapters.
\paragraph{1. GLOBAL\_\allowbreak DENOMINATOR — C — Actual LinkN local table and binding producer}

Closed(Des $\mathcal C$ 3) LinkN: the actual etaSum is bounded with exact dZ0 and same-host vvTerm, subject only to the literal ShaFinOdd and summability guards.

\emph{Shared objects and choice order.} one FamP; one class $\mathcal A$ and its target-compatible subclass; one coefficient field per family datum; one x and i at each application; original selected rows, words and levels; thresholds chosen before the common family; actual D.nd; actual etaJ and Aj; sha\_\allowbreak eq

\emph{Source.} \href{https://github.com/CocoSgt/riemann-hypothesis-lean/blob/7cbf7a0c32fca1badf05931879a37c3f8a2b03b4/src/Formalism/SharpRestriction/GlobalTheoremInputs.lean#L45-L48}{LinkN}.

\emph{Source.} \href{https://github.com/CocoSgt/riemann-hypothesis-lean/blob/7cbf7a0c32fca1badf05931879a37c3f8a2b03b4/src/Formalism/SharpRestriction/GlobalTheoremInputs.lean#L40-L44}{LinkN\_\allowbreak at}.

\emph{Source.} \href{https://github.com/CocoSgt/riemann-hypothesis-lean/blob/7cbf7a0c32fca1badf05931879a37c3f8a2b03b4/src/Development/AutomorphicArithmetic/CharacterBoxes/DivinNOfHypV.lean#L117-L134}{divinN\_\allowbreak of\_\allowbreak hypV}.



\paragraph{2. GLOBAL\_\allowbreak RETAINED\_\allowbreak DEFECT — C — Original retained-row defect bound LinkIRR}

Rows.DelpNonneg and delp$\le$exp(-x\textasciicircum{}(ez/7)) uniformly for every late actual row datum.

\emph{Shared objects and choice order.} one FamP; one class $\mathcal A$ and its target-compatible subclass; one coefficient field per family datum; one x and i at each application; original selected rows, words and levels; thresholds chosen before the common family; AllRows A; LabelCountGoodInjective; AllJoinA; CountAllRows

\emph{New completed inputs.} FLC1–FLC19, reusing RRT1–RRT16 and GQT1–GQT5

\emph{Current remaining subfrontier.} Actual supplied normalized-type compatibility, original HidaSetting, reduction/inflation EtOne with the same Haar normalization, RowsAdd/RowsGood, one common AllRows object and the three arithmetic cylinder counts. All five selected Wg terms now have local producers, with the explicit eventual rank threshold for the last. The stronger unused fourth admLq upper comparison is not asserted.

\emph{Source.} \href{https://github.com/CocoSgt/riemann-hypothesis-lean/blob/7cbf7a0c32fca1badf05931879a37c3f8a2b03b4/src/Formalism/SharpRestriction/GlobalTheoremInputs.lean#L55-L58}{LinkIRR}.

\emph{Source.} \href{https://github.com/CocoSgt/riemann-hypothesis-lean/blob/7cbf7a0c32fca1badf05931879a37c3f8a2b03b4/src/Development/AbelianVarietyTheory/Proofs/GlobalInterpolation/LinkIRROfJoinEx.lean#L17-L38}{linkIRR\_\allowbreak of\_\allowbreak join\_\allowbreak ex}.

\emph{Source.} \href{https://github.com/CocoSgt/riemann-hypothesis-lean/blob/7cbf7a0c32fca1badf05931879a37c3f8a2b03b4/src/Development/AbelianVarietyTheory/Proofs/GlobalInterpolation/LinkIRROfJoinEx.lean#L143-L166}{link\_\allowbreak irr}.



\paragraph{3. GLOBAL\_\allowbreak PERIOD\_\allowbreak HEIGHT — C — Full actual period and height comparison LinkDIVt}

Closed conditional TuSize -> DGamSize -> (LinkEta -> DIVinN -> LinkDIV), with every original LinkDIV\_\allowbreak at clause.

\emph{Shared objects and choice order.} one FamP; one class $\mathcal A$ and its target-compatible subclass; one coefficient field per family datum; one x and i at each application; original selected rows, words and levels; thresholds chosen before the common family; same RR target; same Bu/besS; same Nu and Bk

\emph{Source.} \href{https://github.com/CocoSgt/riemann-hypothesis-lean/blob/7cbf7a0c32fca1badf05931879a37c3f8a2b03b4/src/Formalism/SharpRestriction/GlobalTheoremInputs.lean#L77-L90}{LinkDIV\_\allowbreak at}.

\emph{Source.} \href{https://github.com/CocoSgt/riemann-hypothesis-lean/blob/7cbf7a0c32fca1badf05931879a37c3f8a2b03b4/src/Formalism/SharpRestriction/GlobalTheoremInputs.lean#L94-L96}{LinkDIVt}.

\emph{Source.} \href{https://github.com/CocoSgt/riemann-hypothesis-lean/blob/7cbf7a0c32fca1badf05931879a37c3f8a2b03b4/src/Development/AbelianVarietyTheory/Proofs/GlobalInterpolation/LinkBrest.lean#L7-L25}{GlobalTheory.Top.link\_\allowbreak brest}.

\emph{Source.} \href{https://github.com/CocoSgt/riemann-hypothesis-lean/blob/7cbf7a0c32fca1badf05931879a37c3f8a2b03b4/src/Development/AbelianVarietyTheory/Proofs/GlobalInterpolation/WireSel.lean#L11-L41}{GlobalTheory.Top.wire\_\allowbreak sel}.



\paragraph{4. GLOBAL\_\allowbreak LOCAL\_\allowbreak GAIN — C — Original global LinkG installation}

One uniform LinkG with actual local lattice lengths, served-orbit coverage and MeanKappaBound; delta(x)->0 before varying actual indices.

\emph{Shared objects and choice order.} one FamP; one class $\mathcal A$ and its target-compatible subclass; one coefficient field per family datum; one x and i at each application; original selected rows, words and levels; thresholds chosen before the common family; actual G.Inst; actual GK; actual lenMN; same orbit representatives

\emph{Source.} \href{https://github.com/CocoSgt/riemann-hypothesis-lean/blob/7cbf7a0c32fca1badf05931879a37c3f8a2b03b4/src/Formalism/SharpRestriction/GlobalTheoremInputs.lean#L122-L137}{LinkG\_\allowbreak at}.

\emph{Source.} \href{https://github.com/CocoSgt/riemann-hypothesis-lean/blob/7cbf7a0c32fca1badf05931879a37c3f8a2b03b4/src/Formalism/SharpRestriction/GlobalTheoremInputs.lean#L138-L141}{LinkG}.

\emph{Source.} \href{https://github.com/CocoSgt/riemann-hypothesis-lean/blob/7cbf7a0c32fca1badf05931879a37c3f8a2b03b4/src/Development/AbelianVarietyTheory/Proofs/GlobalInterpolation/LinkG.lean#L14-L27}{GlobalTheory.Top.link\_\allowbreak g}.



\paragraph{5. GLOBAL\_\allowbreak SPECTRAL — C — Original normalized kernel/Gram estimates LinkS}

Closed TuSize -> LinkS with exact Gamma normalization, off-diagonal Hilbert-Schmidt bound, diagonal error and Dall/Dstar estimate, errors exp(-x\textasciicircum{}(epsilon/8)).

\emph{Shared objects and choice order.} one FamP; one class $\mathcal A$ and its target-compatible subclass; one coefficient field per family datum; one x and i at each application; original selected rows, words and levels; thresholds chosen before the common family; actual KD; same S.Cu spaces; GamOf; Dstar; Dall

\emph{Retained predecessor progress.} GQT6, GQT7. Actual normalized Hilbert–Schmidt and dimension estimates; sufficient weaker polynomial rate proved only for the selected final consumer.

\emph{Source.} \href{https://github.com/CocoSgt/riemann-hypothesis-lean/blob/7cbf7a0c32fca1badf05931879a37c3f8a2b03b4/src/Formalism/SharpRestriction/GlobalTheoremInputs.lean#L145-L149}{LinkS\_\allowbreak at}.

\emph{Source.} \href{https://github.com/CocoSgt/riemann-hypothesis-lean/blob/7cbf7a0c32fca1badf05931879a37c3f8a2b03b4/src/Formalism/SharpRestriction/GlobalTheoremInputs.lean#L152-L157}{LinkS}.

\emph{Source.} \href{https://github.com/CocoSgt/riemann-hypothesis-lean/blob/7cbf7a0c32fca1badf05931879a37c3f8a2b03b4/src/Development/AbelianVarietyTheory/Proofs/GlobalInterpolation/LinkAlt.lean#L16-L50}{GlobalTheory.Top.link\_\allowbreak alt}.



\paragraph{6. TASK\_\allowbreak PROP\_\allowbreak GK\_\allowbreak INSTALL — B — Install OneChar and the congruence integer into PropGK}

Complete the finite orbit/Frobenius/valuation argument producing one nonzero integer before every later coefficient prime, with the original e,N,kernel,p and fixed-K constants.

\emph{Source.} \href{https://github.com/CocoSgt/riemann-hypothesis-lean/blob/7cbf7a0c32fca1badf05931879a37c3f8a2b03b4/src/Development/AutomorphicArithmetic/AutomorphicVarieties/UnitaryThreeGroups/KatoUse.lean#L36-L42}{PropGK}.

\emph{Source.} \href{https://github.com/CocoSgt/riemann-hypothesis-lean/blob/7cbf7a0c32fca1badf05931879a37c3f8a2b03b4/src/Development/AutomorphicArithmetic/AutomorphicVarieties/UnitaryThreeGroups/OrbEquiv.lean#L111-L173}{UpperTheory.V.U3.propGK}.



\paragraph{7. TASK\_\allowbreak ONE\_\allowbreak CHAR\_\allowbreak PROOF — C — Reconstruct the original one-character argument}

From exact InpZ and AutomorphyUnramifiedCoefficients, prove the original OneChar trace/length estimate on the supplied coefficient datum and GD input.

\emph{Source.} \href{https://github.com/CocoSgt/riemann-hypothesis-lean/blob/7cbf7a0c32fca1badf05931879a37c3f8a2b03b4/src/Development/AutomorphicArithmetic/AutomorphicVarieties/UnitaryThreeGroups/KatoUse.lean#L54-L64}{OneChar}.

\emph{Source.} \href{https://github.com/CocoSgt/riemann-hypothesis-lean/blob/7cbf7a0c32fca1badf05931879a37c3f8a2b03b4/src/Development/AutomorphicArithmetic/AutomorphicVarieties/ModularityArgument/KatoQuadraticCharacters/OneChar.lean#L158-L178}{UpperTheory.V.Z.Kato.oneChar}.



\paragraph{8. TASK\_\allowbreak INP\_\allowbreak Z\_\allowbreak ASSEMBLY — B — Install the exact Kato input record}

Build the original CohInp/InpZ record using exactly lemM,lemMprime,lemP,RigII,AdVanishK and number on the same arithmetic objects. No rank-four replacement is permitted.

\emph{Source.} \href{https://github.com/CocoSgt/riemann-hypothesis-lean/blob/7cbf7a0c32fca1badf05931879a37c3f8a2b03b4/src/Development/AutomorphicArithmetic/AutomorphicVarieties/ModularityArgument/KatoQuadraticCharacters/KatoTwoOfUseN.lean#L55-L61}{InpZ}.



\paragraph{9. RIG\_\allowbreak II — C — Type-II polarized adjoint rigidity}

Marked original datum, actual type-II point and modularity of all source extensions imply vanishing for its exact adRep.

\emph{Source.} \href{https://github.com/CocoSgt/riemann-hypothesis-lean/blob/7cbf7a0c32fca1badf05931879a37c3f8a2b03b4/src/Development/AutomorphicArithmetic/AutomorphicVarieties/ModularityArgument/KatoQuadraticCharacters/KatoTwoOfUseN.lean#L49-L53}{RigII}.

\emph{Source.} \href{https://github.com/CocoSgt/riemann-hypothesis-lean/blob/7cbf7a0c32fca1badf05931879a37c3f8a2b03b4/src/Development/AutomorphicArithmetic/AutomorphicVarieties/ModularityArgument/KatoQuadraticCharacters/KatoTwoOfUseN.lean#L43-L47}{KatoTwo}.

\emph{Source.} \href{https://github.com/CocoSgt/riemann-hypothesis-lean/blob/7cbf7a0c32fca1badf05931879a37c3f8a2b03b4/src/Development/AutomorphicArithmetic/AutomorphicVarieties/ModularityArgument/KatoQuadraticCharacters/KatoTwoOfUseN.lean#L11-L41}{Used}.



\paragraph{10. AUT\_\allowbreak UNRAM\_\allowbreak COEF — C — Actual unramified-coefficient automorphy}

For each actual Dat with coefficient ramification one and fixed iota, choose one marked g1 supplying FinA, AutB and LemW together.

\emph{Source.} \href{https://github.com/CocoSgt/riemann-hypothesis-lean/blob/7cbf7a0c32fca1badf05931879a37c3f8a2b03b4/src/Development/AutomorphicArithmetic/AutomorphicVarieties/AutomorphicCases/AutomorphyUnramifiedCoefficients.lean#L24-L27}{AutomorphyUnramifiedCoefficients}.

\emph{Source.} \href{https://github.com/CocoSgt/riemann-hypothesis-lean/blob/7cbf7a0c32fca1badf05931879a37c3f8a2b03b4/src/Development/AutomorphicArithmetic/AutomorphicVarieties/AutomorphicCases/AutomorphyUnramifiedCoefficients.lean#L12-L13}{FinA}.

\emph{Source.} \href{https://github.com/CocoSgt/riemann-hypothesis-lean/blob/7cbf7a0c32fca1badf05931879a37c3f8a2b03b4/src/Development/AutomorphicArithmetic/AutomorphicVarieties/AutomorphicCases/AutomorphyUnramifiedCoefficients.lean#L15-L17}{AutB}.

\emph{Source.} \href{https://github.com/CocoSgt/riemann-hypothesis-lean/blob/7cbf7a0c32fca1badf05931879a37c3f8a2b03b4/src/Development/AutomorphicArithmetic/AutomorphicVarieties/AutomorphicCases/AutomorphyUnramifiedCoefficients.lean#L19-L22}{LemW}.



\paragraph{11. CNR3 — C — Uniform rank-three and rank-two congruence integer}

Fixed quadratic K -> c,A -> e,N,p,w,finite Galois-stable roots-of-unity orbit -> nonzero integer NN -> every later coefficient prime, closure, embedding, Frobenius and finite admissible collection.

\emph{Source.} \href{https://github.com/CocoSgt/riemann-hypothesis-lean/blob/7cbf7a0c32fca1badf05931879a37c3f8a2b03b4/src/Development/AutomorphicArithmetic/AutomorphicVarieties/AutomorphicForms/CongruenceNumberRankThree.lean#L125-L140}{CongruenceNumberRankThree}.



\paragraph{12. COUNT\_\allowbreak II\_\allowbreak SUPPORT — C — Two actual rank-four Links with common threshold and same-ideal Hida/P1}

For each of the TWO original QX rank-four Links, at fixed iota choose ONE common coefficient threshold before every later coefficient DVR and actual universal B,Lambda; produce one ideal J carrying ring,dense,pt,end1,end2,twist AND P1core. Inherited JSP/CIP/CAP/CCH/CFO supply the same-J consumer, every-map reciprocity/universal NiceI Norm, ArPl under unchanged CohP, supplied-type IsRes, and the full selected h=0,a=n=4,g=0 Choice with coupled B1/B2, D=4[L+:Q]+4T.q and B2 constant zero. The dualN bound retains one C before every m and an eventual N threshold uniform over each original admissible rep-prime under its finiteness premise, on actual TWPl. LIM constructs the full original S4. PRG now constructs the full original RRObj/RMaps, every finite-level map and actual ordered splitting target; PMS constructs the actual all-level PMod and full S5 with exact support, dimension, fiber and move fields, plus piKer/piSurj and the same-J annihilator equivalence. SCI constructs the actual quotient determinant, OoPair/ooMaps and point/OoPt on one splitting prime. The selected CRW full-block cotangent bound forces the support ideal to vanish, proving selected normalized Core. BCN supplies full Core=>CoreBC and its selected normalization/Unit consumer; PQL supplies full original TWQuot and matched pro-ell descent. Together SCI13 proves selected original Raw at h=0,a=n=4,k.g=0, with final B,J,K,f and every Selmer quantifier unchanged. NTL17 now proves full original neatU for arbitrary integral tau0 and every level. OTW12 proves one original theta before every K with complete-coefficient IsTheta/IsWt. ITA1–ITA13 now construct the full original integral TypeForAlt over the unchanged O, for every IsTeich family and every IsRecT map, and substitute it into OTW13/CAP17 to give the original UnitaryHeckeData/HdEx. The matched crystalline-type theorem and inherited local Langlands/RecTExists foundation are explicit. FNL, OST, WRP, OCT, OHA, SMP and ASP supply all original finite-level/HSet, ordinary/freeness/base-change, weight, control, Hecke-algebra and arithmetic-specialization fields. ADP constructs the actual formU/decompU on the original compact unitary quotient; CFD proves full clParts\_\allowbreak alt on its supplied original AutomorphicHeckeSet. AFS proves domain and DenseWt for every supplied LamObj and weight, TWtBridge, goodScalB\_\allowbreak alt, partAct0\_\allowbreak alt and TCoreSep on their original inputs. Its same-J density implication retains the exact AutomorphicHeckeSet installation input. The remaining arithmetic producers are the actual Galois-decorated partX\_\allowbreak alt/partU\_\allowbreak alt and AutomorphicHeckeSet supplier, the full point-output record, RingProps and the bare same-J density installation, the four exact Hida endpoints PointTheoremBaseObject, ArithmeticPointDecomposition, IrrArithPointsAutomorphicAway and TwistOk, and all-level same-ideal LevelAlt/TWAut/HasAut on both original Links. The selected Raw patching, normalization and descent are no longer open producers. Preserve full CohP, LocCFT, N4, exact determinant/pro-ell conditions, nontrivial residual cyclotomic character, ell=3, fixed iota and one common coefficient threshold. No arbitrary-Kap B1/B2 theorem, Sel1=0 assertion, arbitrary-chain Core or complete same-J arithmetic support producer is claimed.

\emph{New completed inputs.} ITA1–ITA13, ADP1–ADP12, CFD1–CFD12, AFS1–AFS22

\emph{Current remaining subfrontier.} Actual Galois-decorated partX\_\allowbreak alt/partU\_\allowbreak alt and AutomorphicHeckeSet; full Out.pt/RingProps and bare same-J density installation; all four exact Hida endpoints and all-level same-J automorphy on both Links. The complete generic density implication retains that exact supplier. Classical local Langlands/RecTExists and the explicitly named local literature/type inputs retain their stated foundation status.

\emph{Source.} \href{https://github.com/CocoSgt/riemann-hypothesis-lean/blob/7cbf7a0c32fca1badf05931879a37c3f8a2b03b4/src/Development/AutomorphicArithmetic/AutomorphicVarieties/WeakShimuraTheory/IsAlgClosureField.lean#L19-L52}{Link}.

\emph{Source.} \href{https://github.com/CocoSgt/riemann-hypothesis-lean/blob/7cbf7a0c32fca1badf05931879a37c3f8a2b03b4/src/Development/AutomorphicArithmetic/AutomorphicVarieties/ModularityArgument/ThresholdArguments/DenseThresholdCases/PointTheoremBaseObject.lean#L32-L38}{HidaPropsAlt}.

\emph{Source.} \href{https://github.com/CocoSgt/riemann-hypothesis-lean/blob/7cbf7a0c32fca1badf05931879a37c3f8a2b03b4/src/Development/AutomorphicArithmetic/AutomorphicVarieties/WeilGroupCharacters/UniversalOrdinaryDeterminantRing.lean#L284-L286}{P1core}.



\paragraph{13. TASK\_\allowbreak CASE3\_\allowbreak AUT\_\allowbreak JOIN — B — Complete the original Case-III automorphy/image join}

Install the actual Case-III ProModAlt, SeedIVAlt and shared PMExAlt with the inherited CohomologyHyp/Ferrero-Washington/HFit/SeedII facts into the original whole Case-III automorphy/image conclusion.

\emph{Source.} \href{https://github.com/CocoSgt/riemann-hypothesis-lean/blob/7cbf7a0c32fca1badf05931879a37c3f8a2b03b4/src/Development/AutomorphicArithmetic/AutomorphicVarieties/ModularityArgument/ThresholdArguments/DenseThresholdCases/PromodularityCaseThree.lean#L43-L44}{AutomorphyRankFourDense}.

\emph{Source.} \href{https://github.com/CocoSgt/riemann-hypothesis-lean/blob/7cbf7a0c32fca1badf05931879a37c3f8a2b03b4/src/Development/AutomorphicArithmetic/AutomorphicVarieties/ModularityArgument/ThresholdArguments/DenseThresholdCases/CuFinal.lean#L44-L44}{main4\_\allowbreak alt}.



\paragraph{14. TASK\_\allowbreak CASE3\_\allowbreak PROMOD\_\allowbreak JOIN — C — Reconstruct the actual Case-III chain-to-promodularity proof}

Use the original tensor-factor tower, ChainRankFour, ChainDesc4, seeded rank-four chain reduction, actual PMExAlt and SeedIVAlt to construct ProModAlt on every original Ctx3 star point at the same fixed iota, preserving all finite-coefficient, inertia and support maps.

\emph{Source.} \href{https://github.com/CocoSgt/riemann-hypothesis-lean/blob/7cbf7a0c32fca1badf05931879a37c3f8a2b03b4/src/Development/AutomorphicArithmetic/AutomorphicVarieties/ModularityArgument/ThresholdArguments/DenseThresholdCases/PromodularityCaseThreeOfLemmas.lean#L131-L153}{UpperTheory.V.Cu.promodularity\_\allowbreak case\_\allowbreak three\_\allowbreak of\_\allowbreak chain\_\allowbreak reduction}.

\emph{Source.} \href{https://github.com/CocoSgt/riemann-hypothesis-lean/blob/7cbf7a0c32fca1badf05931879a37c3f8a2b03b4/src/Development/AutomorphicArithmetic/AutomorphicVarieties/ModularityArgument/ThresholdArguments/DenseThresholdCases/PromodularityCaseThree.lean#L40-L42}{PromodularityCaseThreeDense}.



\paragraph{15. SEED\_\allowbreak IV\_\allowbreak ALT — C — Actual Case-III seed with prechosen residual datum}

Base Ctx3/Four tower, iota and Amb -> w,s before every later Dies tower -> compatible finite-coefficient seed with exact residual support/root entries.

\emph{Source.} \href{https://github.com/CocoSgt/riemann-hypothesis-lean/blob/7cbf7a0c32fca1badf05931879a37c3f8a2b03b4/src/Development/AutomorphicArithmetic/AutomorphicVarieties/ModularityArgument/ThresholdArguments/DenseThresholdCases/PromodularityCaseThree.lean#L19-L25}{SeedIVAlt}.



\paragraph{16. TENSOR\_\allowbreak FACTOR\_\allowbreak TOWER — B — Actual tensor-factor tower and characters}

Original finite-field, tame, factor and character constraints; source-level construction is sizeable and has not been reconstructed here.

\emph{Source.} \href{https://github.com/CocoSgt/riemann-hypothesis-lean/blob/7cbf7a0c32fca1badf05931879a37c3f8a2b03b4/src/Development/AutomorphicArithmetic/AutomorphicVarieties/ModularityArgument/ThresholdArguments/DenseThresholdCases/PromodularityCaseThreeOfLemmas.lean#L131-L153}{UpperTheory.V.Cu.promodularity\_\allowbreak case\_\allowbreak three\_\allowbreak of\_\allowbreak chain\_\allowbreak reduction}.



\paragraph{17. CHAIN\_\allowbreak RANK\_\allowbreak FOUR — C — Actual SomeIn-to-AllIn propagation}

Same source Space and PM ideal, residual point Good4, original bound 30(C+1)\textasciicircum{}2 <= d; geometry must preserve all identities.

\emph{Source.} \href{https://github.com/CocoSgt/riemann-hypothesis-lean/blob/7cbf7a0c32fca1badf05931879a37c3f8a2b03b4/src/Development/AutomorphicArithmetic/AutomorphicVarieties/ModularityArgument/ThresholdArguments/DenseThresholdCases/PromodularityCaseThree.lean#L26-L32}{ChainRankFour}.



\paragraph{18. CHAIN\_\allowbreak DESC4 — B — Original residual chain descent}

Finite residual descent retaining source Chain4 compatibility.

\emph{Source.} \href{https://github.com/CocoSgt/riemann-hypothesis-lean/blob/7cbf7a0c32fca1badf05931879a37c3f8a2b03b4/src/Development/AutomorphicArithmetic/AutomorphicVarieties/ModularityArgument/ThresholdArguments/DenseThresholdCases/PromodularityCaseThreeOfLemmas.lean#L131-L153}{UpperTheory.V.Cu.promodularity\_\allowbreak case\_\allowbreak three\_\allowbreak of\_\allowbreak chain\_\allowbreak reduction}.



\paragraph{19. RANK4\_\allowbreak CHAIN\_\allowbreak REDUCTION — C — Seeded reduction to the original residual chain}

Select T1 before every later Deep tower; retain one compatible seed and Chain4 for every original residual lift over Q.

\emph{Source.} \href{https://github.com/CocoSgt/riemann-hypothesis-lean/blob/7cbf7a0c32fca1badf05931879a37c3f8a2b03b4/src/Development/AutomorphicArithmetic/AutomorphicVarieties/ModularityArgument/ThresholdArguments/DenseThresholdCases/PromodularityCaseThree.lean#L33-L39}{RankFourChainReductionDense}.

\emph{Source.} \href{https://github.com/CocoSgt/riemann-hypothesis-lean/blob/7cbf7a0c32fca1badf05931879a37c3f8a2b03b4/src/Development/AutomorphicArithmetic/AutomorphicVarieties/ModularityArgument/ThresholdArguments/DenseThresholdCases/CuPChainIV.lean#L16-L82}{UpperTheory.V.Cu.chainIV\_\allowbreak alt}.



\paragraph{20. TASK\_\allowbreak COUNT\_\allowbreak V\_\allowbreak INSTALL — B — Complete the original CountV finite assembly}

From exactly the listed source host/cohomology inputs, complete countV\_\allowbreak of\_\allowbreak finiteness\_\allowbreak duality\_\allowbreak euler\_\allowbreak host, including original host/frame formation, normalized length, the CaseV guard and R1<=2\textasciicircum{}c2.

\emph{Source.} \href{https://github.com/CocoSgt/riemann-hypothesis-lean/blob/7cbf7a0c32fca1badf05931879a37c3f8a2b03b4/src/Development/AutomorphicArithmetic/AutomorphicVarieties/HeightAndVanishing/Counts.lean#L69-L69}{CountV}.

\emph{Source.} \href{https://github.com/CocoSgt/riemann-hypothesis-lean/blob/7cbf7a0c32fca1badf05931879a37c3f8a2b03b4/src/Development/AutomorphicArithmetic/AutomorphicVarieties/ModularityArgument/ThresholdArguments/DenseThresholdCases/CountVOf.lean#L33-L60}{UpperTheory.V.UnitaryThreeInertia.countV\_\allowbreak of\_\allowbreak finiteness\_\allowbreak duality\_\allowbreak euler\_\allowbreak host}.



\paragraph{21. THM\_\allowbreak S6\_\allowbreak S — C — Original unitary-three ordinary support and automorphy}

Original host/frame and fixed iota -> one marked g1 -> every universal R,u -> weight-zero finiteness and every Type-I BPrime\_\allowbreak s point.

\emph{Source.} \href{https://github.com/CocoSgt/riemann-hypothesis-lean/blob/7cbf7a0c32fca1badf05931879a37c3f8a2b03b4/src/Development/AutomorphicArithmetic/AutomorphicVarieties/ModularityArgument/ThresholdArguments/DenseThresholdCases/ThmS6CoreS.lean#L17-L26}{ThmS6\_\allowbreak s}.

\emph{Source.} \href{https://github.com/CocoSgt/riemann-hypothesis-lean/blob/7cbf7a0c32fca1badf05931879a37c3f8a2b03b4/src/Development/AutomorphicArithmetic/AutomorphicVarieties/ModularityArgument/ThresholdArguments/DenseThresholdCases/ThmS6CoreS.lean#L12-L16}{BPrime\_\allowbreak s}.



\paragraph{22. HOST\_\allowbreak FRAME\_\allowbreak LENGTH — B — Original Case-V host/frame and normalized-length installation}

Actual HostData, frame, local-condition equality, abstract count and normalized original Selmer length must all refer to the same objects.

\emph{Source.} \href{https://github.com/CocoSgt/riemann-hypothesis-lean/blob/7cbf7a0c32fca1badf05931879a37c3f8a2b03b4/src/Development/AutomorphicArithmetic/AutomorphicVarieties/ModularityArgument/ThresholdArguments/DenseThresholdCases/CountVOf.lean#L33-L60}{UpperTheory.V.UnitaryThreeInertia.countV\_\allowbreak of\_\allowbreak finiteness\_\allowbreak duality\_\allowbreak euler\_\allowbreak host}.



\paragraph{23. HOST\_\allowbreak NUMBER\_\allowbreak S — C — Uniform original-host rank-three catalogue}

One exponent before all closures, coefficient rings, levels, characters, frames, embeddings and marked points; one separating trace function controls all later finite pairwise nonisomorphic collections.

\emph{Source.} \href{https://github.com/CocoSgt/riemann-hypothesis-lean/blob/7cbf7a0c32fca1badf05931879a37c3f8a2b03b4/src/Development/AutomorphicArithmetic/AutomorphicVarieties/ModularityArgument/ThresholdArguments/DenseThresholdCases/HostNumberS.lean#L18-L28}{HostNumber\_\allowbreak s}.



\paragraph{24. TASK\_\allowbreak HOST\_\allowbreak VAN\_\allowbreak E\_\allowbreak INSTALL — B — Transport K/duality to literal HostVanE}

Use the already supplied LemmaK/local-global duality and finiteness on the actual CountV host/frame and exact original coefficient modules.

\emph{Source.} \href{https://github.com/CocoSgt/riemann-hypothesis-lean/blob/7cbf7a0c32fca1badf05931879a37c3f8a2b03b4/src/Development/AutomorphicArithmetic/AutomorphicVarieties/UnitaryThreeDeformations/HostNumber.lean#L9-L12}{HostVanE}.

\emph{Source.} \href{https://github.com/CocoSgt/riemann-hypothesis-lean/blob/7cbf7a0c32fca1badf05931879a37c3f8a2b03b4/src/Development/AutomorphicArithmetic/AutomorphicVarieties/LevelCohomology/SoHostVanE.lean#L11-L18}{UpperTheory.V.X.So.hostVanE}.



\paragraph{25. TASK\_\allowbreak HOST\_\allowbreak XPP\_\allowbreak INSTALL — B — Transport CorX to literal HostXpp}

Preserve Psi\textasciicircum{}-2 eta, the prime-to-ell quotient, actual local conditions and the original CountV host while installing BR7.

\emph{Source.} \href{https://github.com/CocoSgt/riemann-hypothesis-lean/blob/7cbf7a0c32fca1badf05931879a37c3f8a2b03b4/src/Development/AutomorphicArithmetic/AutomorphicVarieties/UnitaryThreeDeformations/HostNumber.lean#L14-L17}{HostXpp}.

\emph{Source.} \href{https://github.com/CocoSgt/riemann-hypothesis-lean/blob/7cbf7a0c32fca1badf05931879a37c3f8a2b03b4/src/Development/AutomorphicArithmetic/AutomorphicVarieties/LevelCohomology/Xpp.lean#L11-L84}{UpperTheory.V.X.xpp}.



\paragraph{26. DESIGN\_\allowbreak CHOICES\_\allowbreak COMMON — C — Construct the original common design family and bundle}

For the one target HostExW, construct the original compatible DesignChoices on Cdes=WithTw(WithSm Cprime), where Cprime=WithChN(A intersect target class), A is the one Built construction class, and cT(eps,y)=cTmin(y). Retain every original field, tuple, word, metric and common bundle duty. DSZ10 and RMP11 give size/rank properties of class members, not member existence. The original RecExP/RecAtG and common object/row route remains coupled. Construct ONE original common design family R and its recursive/geometric half with original Inputs/Pins/Holds/ExAt/InGeo and required CoreP. PKT now reconstructs full exists\_\allowbreak pkT, rowPk\_\allowbreak nonempty and exists\_\allowbreak moved with Copy/Coh on those original given inputs and a supplied fitting later stage: one catalogue, both ChT, complete LfT, actual integral lattices and all metric/measure/period constants are retained; one numerical record precedes every eta and V/S/N/G. RDR supplies one Side0 before all readers, ReadPk/pins and full ReadL with exact o.eta=q.eta. The initial actual geometry, Fits/division/central conditions, primitive rational tests and eta agreement must still be produced. Inherited CSR supplies finite columns, one ReadC before later twists, period/evaluation and scalar/error transports, RowSideDetermined and same-R Bundle/BTw/BWit/BunU consumers. RPI supplies exact RowPairInt for actual pairs and finite translates with original mu\_\allowbreak b/fd0; first literal DivDef3 at the same updated R,F,Y,Z,C,sEE follows from its retained original readings, Pins/Holds, stored nonzero IsItW chain, five real-group predicates and group tie. Remaining common fields include host/frame, MultEq nonzero-part support, that original nonzero chain/group tie, eigensystem/carrier/parts and all-embedding logLam, HD/d3b, NEb, actual rational tests and base rows. Retain every simultaneous family hypothesis and eventual threshold. Stronger second DivRows3 remains separate; universal quotient-height CuspRapidDecay and actual common member existence are not inferred from the comparison constructors. Version 0.41 supplies the finite phase-correct rational theta route: exact coefficient extension, split representatives, square-law/PhiIs, Pins from CoversAt, class-sum/point-identity transfer to OneMapOverL, actual-good-projector RatG and both-level LiftTestB. Its original higher covers, all-tuple class-sum uniqueness including phase rescaling, one PointIdentCore before O, base ClassSumsSpan and geometric LatFull remain explicit producers. MLat is available from the original LatFull input. This sharpens the rational-test work without claiming actual common geometry, Fits or eta agreement.



\paragraph{27. LOCAL\_\allowbreak NBUILT\_\allowbreak COMMON — C — Actual local specialization on the common selected family}

Construct, with the required closed-class property on the same Des(C,3), one eventual family of actual NBuilt certificates for the original Fm.D(x,i), retaining the certificate before forgetting to TSel. HLP now reconstructs full HeckeLift0, all-index commutation, HeckeStd0 and HeckeIntTw from the original MtInput and complex Hecke facts, in the full ambient GenAt range k<=nV. Every maximal ideal is an actual supported residual character class; literal MtSub is the O-span of projections of the original MtLat generators, is finite free, and localizes onto the exact NCov.part. Literal Psi has its exact range and off-space projection error, both original span assertions hold, and ActEx/ActHK are constructed on the same lattice under original semisimplicity. Original MjLat relation quotients, universal coefficient-map relation fields, selected tame generators, coefficient embeddings and all identity transports give the complete nonempty\_\allowbreak spec constructor on its full printed input list. A single compositum combines the four original coefficient requirements with big fixed beforehand. Whole ndSpec\_\allowbreak exists still needs actual common geometric/choice MtInput, complex HeckeStdPlain and the no-Jordan producer, original F3 bridges with valuation/GoodCoef/iq, original F4 common-overfield reading of the entire coefficient field, and exact arithmetic character/reciprocity/admissibility suppliers. Original AutSpec, Galois specialization and the remaining NBuilt fields on that same installation also remain. Do not replace the actual projected lattice or maps by abstract objects with the same dimensions.



\paragraph{28. SECOND\_\allowbreak LAYER\_\allowbreak COMPATIBLE\_\allowbreak SELECTION — C — Original second-layer selector and class duties}

For the same A, supply Duty on A and a Closed common original TSel family whose selected records are the original true-column rec7 records (or have the entire dependent tuple transported), with their real witness and original coefficient-field root data. Existence/compatibility remains arithmetic; merely abstractly isomorphic fields or independently selected columns do not suffice. CSR21–22 now construct the same-R Bundle/BTw/BWit/BunU outputs from the exact selected realization/row-side packet and NEb, with the original EvP/ParEv quantifier order. Actual common selection, the five stated realization packets and NEb remain upstream; none is replaced by an independently chosen compatible record. Version 0.38 closes the selected analytic extension and reconstructs the first updated DivDef3 consumer; it leaves the original common selection, geometric readings, rational tests, base rows, host/parts/NEb and exact retained DI premises upstream, with unchanged EvP/ParEv thresholds. Version 0.39 supplies the original given-data reader constructors and exact RowPairInt/first DivDef3 on original premises. It leaves one actual common geometric/comparison/test member, eta agreement, compatible selection, rational tests, base rows, host/parts/NEb and original EvP/ParEv thresholds upstream. Version 0.40 also supplies original comparison/RowPk/fitting-stage Moved algorithms and exact common data transport. The actual initial geometric member, fitting tests, eta agreement and compatible second-layer selector with original eventual thresholds remain upstream. Version 0.41 supplies the finite phase-correct rational theta route: exact coefficient extension, split representatives, square-law/PhiIs, Pins from CoversAt, class-sum/point-identity transfer to OneMapOverL, actual-good-projector RatG and both-level LiftTestB. Its original higher covers, all-tuple class-sum uniqueness including phase rescaling, one PointIdentCore before O, base ClassSumsSpan and geometric LatFull remain explicit producers. MLat is available from the original LatFull input. This sharpens the rational-test work without claiming actual common geometry, Fits or eta agreement.



\paragraph{29. BASE\_\allowbreak THETA\_\allowbreak READINGS — C — Actual base reading on the common bundle}

For every family in the original class, eventually on the SAME selected Bundle and for each i<m, supply the actual R,F,Read,ReadL, weak models/Weil data and phase-one pins. Identify their LiftTie.theta, mtRead, mjRead and ThDef3 basis with that Bundle.P.A3, P.Zs, P.S0 and its stored maps. Exact field, embeddings, metrics, measures, row family and weak choices are retained. This remains one part of the coupled BTH12 producer. PKT now supplies the complete original comparison packet, RowPk and fitting-stage Moved with Copy/Coh on original given geometry, Inputs/Pins/Holds/ExAt/InGeo and Fits. Its ChT/LfT/Reps, actual lattice carriers, original Haar and norm constants, period maps, representatives and all-place RrT are constructed together; its one numerical record precedes every eta/V/S/N/G. RDR gives the original level-zero side, moved-Read and ReadPk/pin constructors on those objects, with one Side0 before every reader and all twelve Read ties. Full ReadL retains the exact o.eta(j,v)=q.eta(j,v) guard for every 1<=j<=m and place v; arbitrary Over does not imply it. Inherited CSR fixes one ReadC before all TwArgs, transports both metrics and all thirteen reading maps, and gives actual period/evaluation with kappa=(c*nP)\textasciicircum{}(-1/2). RPI supplies exact row-pair integrability, so first updated DivDef3 follows from the retained original Pins/Holds, nonzero chain, five real-group predicates and group tie. Initial theta extension remains complete. Remaining work is actual initial common geometry and its original predicates, Fits/primitive rational tests, eta agreement and outer family/eventual/same-Bundle duties, reFace compatibility, MultEq support, original chain/group data, actual base rows and any stronger DivRows3 theta square. No new compatible reading, family threshold or Haar normalization is substituted. Version 0.41 supplies the finite phase-correct rational theta route: exact coefficient extension, split representatives, square-law/PhiIs, Pins from CoversAt, class-sum/point-identity transfer to OneMapOverL, actual-good-projector RatG and both-level LiftTestB. Its original higher covers, all-tuple class-sum uniqueness including phase rescaling, one PointIdentCore before O, base ClassSumsSpan and geometric LatFull remain explicit producers. MLat is available from the original LatFull input. This sharpens the rational-test work without claiming actual common geometry, Fits or eta agreement.



\paragraph{30. BASE\_\allowbreak THETA\_\allowbreak RATIONAL\_\allowbreak TESTS — C — Simultaneous original base rational tests}

On the SAME original Bundle/reading/theta map, for every i<m supply one finite separating complex-linear test family whose values lie in the fixed embedded coefficient field on theta of EVERY original rational source vector and on EVERY stored target-coordinate basis vector. BTH8/BTH11 use this same LiftTestB for rational adjoint image and matrix coordinates. RTG now reconstructs prescribed-field coefficient extension, original split class representatives/RepsOK, the self-dual split-lattice square law and finite residue PhiIs; every Pins field follows from the actual CoversAt. With the base coefficient record literally fixed, exact class-sum existence/uniqueness including phase-rescaled O, PointIdentCore with common lambda,z before every O, and base ClassSumsSpan imply all-tuple finite lift span and OneMapOverL. The phase is exactly cj(twAt embLL A3.A.v0)(R.phs i). RatG uses the actual good projector eOp(y), and full LiftTestB follows at BOTH original levels. These finite implication proofs are complete. Remaining original producers are the prescribed higher covers, those actual class sums/uniqueness, PointIdentCore, base ClassSumsSpan and the base/common geometric datum; exact MLat is supplied only from original LatFull, whose geometric production remains separate. Finite dimension or an unrelated rational theta model cannot replace these same-object identities. Original outer eventual-family/Bundle quantifiers remain.



\paragraph{31. SLR\_\allowbreak BASE\_\allowbreak ROWS\_\allowbreak PRODUCER — C — Actual base row identities on that same bundle}

For the SAME selected Bundle, common original fields and readings used by the other three BTH12 inputs, supply actual RowsHold0\_\allowbreak 3 uniformly eventually for every family. Its original primitive row realization, spectrum/type and basis identifications must be retained. BTH supplies a theta consumer only; it does not construct the base row identity. No separate witness may be chosen merely to satisfy this row requirement.



\clearpage
\section{Integral local types over the original coefficient ring}
\label{sec:integral-type-v43}

\status{This chapter supplies the original integral
\texttt{Type\allowbreak For\allowbreak Alt} input to OTW13.  The construction keeps the
given coefficient ring, residual characters, Galois quotient,
local groups and complex coefficient map.  Its only new substantial
published input is the crystalline inertial-type theorem, at the
precise scope stated in ITA6.  We prove the descent to the original
ring and the universal local-correspondence and Teichm\"uller
quantifiers.  Classical local Langlands and the original
nonemptiness theorem \texttt{Rec\allowbreak TExists} remain foundational inputs;
the latter is used only to prevent a universal quantifier over
correspondences from being vacuous.  No global automorphy,
nonvanishing, Hida endpoint, or RH/GRH conclusion is asserted here.}

Fix the original prime $\ell$, CM field $L$, algebraic closure
$\Omega$, setting $s:\texttt{Hida\allowbreak Setting}$, integer $n\geq1$,
level datum $V:s.\texttt{Lev}$, auxiliary datum $Q$, coefficient
datum $c$, residual datum $r$, and field isomorphism
\[
 \iota:\overline{\mathbb Q}_\ell\xrightarrow{\sim}\mathbb C.
\]
Write $O=c.O$, $K=\operatorname{Frac}(O)$,
$k=O/\mathfrak m_O$, and $q_O=|k|$.  The given coefficient
map is
\[
 j:O\lhook\joinrel\longrightarrow\overline{\mathbb Q}_\ell
              \xrightarrow{\iota}\mathbb C.
 \tag{ITA.1}
\]
We also write $j$ for its extension to $K$; neither $j$ nor
$\iota$ is assumed continuous for the complex topology.
The original residual characters are
$\bar\chi_i=r.\texttt{chib}_i:G_0\to k^\times$,
$0\leq i<n$.  Their kernels are open.  Their pairwise
distinctness, polarization and the other fields of $r$ will not
be needed for this construction.

The finite set $S=s.\texttt{St}$ is the source's set of type
places.  For $u\in S$, put
\[
 F_u=L_{s.w(u)},\quad G_u=\operatorname{GL}_n(F_u),\quad
 K_u=\operatorname{GL}_n(\mathcal O_{F_u}),\quad
 D_u=\operatorname{Dec}_L(s.A(u)),\quad
 I_u=\operatorname{Ine}_L(s.A(u)).
 \tag{ITA.2}
\]
These are the actual groups at the specified valuations.
The original map $V.\mathrm{pr}:G_L\to G_0$ restricts to
each $D_u$ and $I_u$.  Write
$G_S=\prod_{u\in S}G_u$ and $K_S=\prod_{u\in S}K_u$;
these identify with the source's \texttt{GV} and \texttt{KT}
by the coordinate inclusions.  The residue characteristic of
$F_u$ can differ from $\ell$; in this setting it does differ.
We will impose no bound comparing either prime to $n$.

Here is the exact conclusion.  There are an integer $m$ and one
homomorphism
\[
                    \tau_0:K_S\longrightarrow\operatorname{GL}_m(O)
 \tag{ITA.3}
\]
such that its complexification $\tau_{0,j}$ has open kernel and
the following equivalence holds.  It holds for every family
$t_i:G_0\to O^\times$ satisfying
\[
 \overline{t_i(g)}=\bar\chi_i(g),\qquad
                       t_i(g)^{q_O-1}=1,                    \tag{ITA.4}
\]
every nonzero complex vector space $M$, every representation
$\rho:G_S\to\operatorname{Aut}_{\mathbb C}(M)$, and every
family of smooth irreducible representations $\pi_u$ for which
$M$, restricted to $G_u$, is generated by the images of all
$G_u$-equivariant maps from $\pi_u$.  The last condition is the
literal \texttt{Iso\allowbreak Sum} condition; no finite length or
admissibility of $M$ is added.  Define
\[
 \chi^D_{u,i}(g)=j\bigl(t_i(V.\mathrm{pr}(g))\bigr),\qquad
                              \chi_{u,i}=\chi^D_{u,i}|_{I_u}.
 \tag{ITA.5}
\]
The asserted equivalence is
\begin{align}
 &\exists (F_a)_{0\leq a<m}\in M^m\setminus\{0\},\quad
   \sum_{b=0}^{m-1}j(\tau_0(k)_{ab})\,\rho(k)F_b=F_a
                 \quad(k\in K_S,\ 0\leq a<m)
                                                        \tag{ITA.6}\\
 &\hspace{25mm}\Longleftrightarrow\quad
   \forall u\in S\ \forall\mathrm{rc},\quad
   \texttt{Is\allowbreak Rec\allowbreak T}(s.hw(u),\mathrm{rc})\Longrightarrow
   \texttt{Is\allowbreak Type\allowbreak WD}\bigl(\chi_u,
                       \mathrm{rc}(\operatorname{algNorm}\pi_u)\bigr).
                                                        \notag
\end{align}
Here
$\operatorname{algNorm}\pi=\pi\otimes|\det|^{(1-n)/2}$.
The source predicate \texttt{Is\allowbreak Type\allowbreak WD} asks for a representative
of its Weil--Deligne isomorphism class on $\mathbb C^n$ with
$N=0$ and inertia matrices exactly
$\operatorname{diag}(\chi_0,\ldots,\chi_{n-1})$.  Repeated
inertia characters are allowed.  Formula (ITA.6), including the
order of all its quantifiers, is \texttt{Type\allowbreak For\allowbreak Alt}.

\subsection{The actual coefficient ring and every Teichm\"uller lift}

% BEGIN ITA1
\paragraph{ITA1. Completeness and the algebraic coefficient closure.}
\label{claim:ITA1}
The ring $O$ is the complete valuation ring of the finite extension
$K/\mathbb Q_\ell$, its residue field is finite, and
$\mathbb C$ is an algebraic closure of $j(K)$.
\begin{proof}
GA1 identifies the actual subring $O\subset\overline{\mathbb Q}_\ell$
with the elements of $K$ of norm at most one, and GA2 proves
that $K$ is finite over $\mathbb Q_\ell$.  The equivalence of
the coordinate and induced norms in GA2 proves that $K$ is
complete.  Its closed unit ball $O$ is therefore complete.
The norm topology and the maximal-ideal topology agree by
GA1's uniformizer formula.  Since $O$ is a finite
$\mathbb Z_\ell$-module, $O/\ell O$ is finite; $k$ is its
quotient because $\ell\in\mathfrak m_O$.  Finally every
element of $\overline{\mathbb Q}_\ell$ is algebraic over
$\mathbb Q_\ell\subset K$, and $\iota$ is onto the
algebraically closed field $\mathbb C$.  Thus every complex
number is algebraic over this particular embedded copy $j(K)$.
This is an algebraic statement and supplies no complex
continuity assertion.
\end{proof}
% END ITA1

% BEGIN ITA2
\paragraph{ITA2. One and only one family satisfies the original lift predicate.}
\label{claim:ITA2}
There is a unique multiplicative lift $[\ ]:k^\times\to O^\times$
with $[a]^{q_O-1}=1$ and residue $a$.  Hence (ITA.4) has the
unique solution $t_i=[\ ]\circ\bar\chi_i$, and
\[
                         \ker t_i=\ker\bar\chi_i.       \tag{ITA.7}
\]
\begin{proof}
For $a\in k^\times$, choose any unit lift $a_0\in O$.
The polynomial $f(X)=X^{q_O-1}-1$ vanishes at $a_0$
modulo $\mathfrak m_O$, and
$f'(a_0)=(q_O-1)a_0^{q_O-2}$ is a unit: the residue of
$q_O-1$ is $-1$.  Newton iteration
$a_{r+1}=a_r-f(a_r)/f'(a_r)$ stays in the same residue
class, keeps the derivative invertible, and at least doubles
the valuation of $f(a_r)$.  Its successive differences tend
to zero geometrically in the maximal-ideal filtration.
Completeness from ITA1 gives a root $[a]$ in that residue
class.  This also proves existence directly at $\ell=3$.

For uniqueness let $x,y\in O^\times$ have that residue and
$(q_O-1)$st power one.  Put $z=xy^{-1}$, so $\bar z=1$.
Then
\[
          (z-1)(1+z+\cdots+z^{q_O-2})=z^{q_O-1}-1=0.
\]
The second factor is a unit, since its residue is $-1$.
Consequently $z=1$.  Products of the chosen lifts satisfy
the same equation and have the product residue, proving
$[ab]=[a][b]$; also $[1]=1$.  Applying uniqueness to each
value proves uniqueness of the entire family $t_i$ among
all homomorphisms in the source predicate.  A value $t_i(g)$
equals one exactly when its residue equals one, which gives
(ITA.7).  All kernels are therefore open, and all values
belong to the finite group of $(q_O-1)$st roots of unity.
\end{proof}
% END ITA2

% BEGIN ITA3
\paragraph{ITA3. The precise local inertia types and their field of values.}
\label{claim:ITA3}
For every $u$, the characters (ITA.5) on the original $D_u$
have finite image and open kernel.  They have values in
$j(O^\times)$, and
$\vartheta_u=\bigoplus_{i=0}^{n-1}\chi_{u,i}$ is an
inertial type for $F_u$.
\begin{proof}
The map $V.\mathrm{pr}$ is the original continuous restriction
and quotient map.  By (ITA.7), the kernel of its composite
with $t_i$ is the inverse image of the open subgroup
$\ker\bar\chi_i$.  The embedding $j$ is injective, so
pushing coefficients through $j$ does not change this
kernel.  The image is finite by ITA2, proving all the stated
topological assertions without any continuity of $j$.
The value assertion is the displayed formula (ITA.5) itself.

Complete at the fixed valuation $s.A(u)$.  Restriction and
completion of finite extensions identify $D_u,I_u$ and its
Weil subgroup with $G_{F_u},I_{F_u}$ and $W_{F_u}$, carrying
the specified subgroup inclusions to the usual ones.  This
is the same completion identification used in KX3 and is
valid at every finite place.  The direct sum of the
$\chi^D_{u,i}|_{W_{F_u}}$ extends $\vartheta_u$ to the
Weil group and has finite image and open kernel there.
Thus it meets the published definition of inertial type.
Neither distinctness of the $\chi_{u,i}$ nor a tame-image
condition is required.
\end{proof}
% END ITA3

\subsection{The local-correspondence dictionary, including its quantifiers}

For one place write $F=F_u$, $D=D_u$, $I=I_u$, $W=W_{F_u}$,
and $\chi_i^D=\chi^D_{u,i}$.  Let $\mathrm{rec}$ denote
the usual complex local Langlands correspondence with geometric
reciprocity.  Retain its rank-one formula and its compatibility
with ordered Langlands quotients.  For an ordered list of smooth
characters $\mu_0,\ldots,\mu_{n-1}$ satisfying
\[
                  i<j\ \Longrightarrow\ \mu_j\ne\mu_i|\ |,
                                                               \tag{ITA.8}
\]
normalized induction has a unique irreducible quotient; its
parameter is the sum of the corresponding Weil characters with
$N=0$.  These are the rank-one part of the usual Langlands
classification, not a claim that every principal series is
irreducible.  They are also precisely the \texttt{one},
\texttt{langlands} and singleton-segment laws used below.
One primary reference for these conventions and the quotient
statement is Caraiani--Emerton--Gee--Geraghty--Pa\v{s}k\={u}nas--Shin,
\emph{Patching and the $p$-adic local Langlands correspondence},
\S\,1.8, \S\,3.1 and Proposition 3.9, with its proof via
Bernstein--Zelevinsky classification.
\footnote{Author's full paper:
\url{https://www.ma.imperial.ac.uk/~tsg/Index_files/Gpatching.pdf}.
The local statement is independent of the global patching
restriction in that paper.  For the unrestricted inertial-type
theorem itself we use the separate statement in ITA6.}

The source has a larger class of maps, called \texttt{Is\allowbreak Rec\allowbreak T}.
We use exactly the following fields of such a map: a map
$a:W\to F^\times$ satisfying \texttt{Is\allowbreak Local\allowbreak Artin}; the
rank-one formula; the ordered-segment formula; compatibility
with real unramified twists; and injectivity against a
\texttt{Good} representation.  Representations with a
cuspidal support consisting of rank-one characters are
\texttt{Good} by the first disjunct of its definition.
The source's theorem \texttt{Rec\allowbreak TExists} supplies at least
one such map at the given place.  Its exact foundational
endpoint is
\path{Foundations/LocalLanglands/LocalCentralUnits/RecTExists.lean},
and the caller is \texttt{Win\allowbreak Conv.\allowbreak ex\allowbreak Rec'\_\allowbreak alt}.  This
nonemptiness is retained as part of the original local
Langlands foundation.  Our argument does not replace
\texttt{Is\allowbreak Rec\allowbreak T} by a preferred map or assume its maps agree
on arbitrary representations.

% BEGIN ITA4
\paragraph{ITA4. Every allowed local Artin map at a type place is the geometric one.}
\label{claim:ITA4}
Every map $a:W\to F^\times$ satisfying the source's
\texttt{Is\allowbreak Local\allowbreak Artin} equals the inverse of geometric local
reciprocity on the Weil abelianization.
\begin{proof}
The $u\in s.\mathrm{St}$ branch of $s.\mathrm{split}$
says that the original place below $w=s.w(u)$ splits in
$L/L^+$.  The same decomposition-group argument as in
CIP3 therefore gives $w\ne\bar w$.  In particular CIP3's
local-unit injectivity applies to this place, although
it lies away from $\ell$.

Put $a_0=(\operatorname{rec}^{\mathrm{geo}}_F)^{-1}$ on
$W^{\mathrm{ab}}$, using the actual completion in ITA3.
Local class field theory gives its surjectivity, continuity
on inertia, and
$\operatorname{ord}_F(a_0(\sigma))=-\deg(\sigma)$,
where degree one denotes arithmetic Frobenius.  A smooth
Weil character factors through this abelianization and
therefore through $a_0$, giving the source's \texttt{factor}
field.  The resulting character of $F^\times$ has open
kernel, because the reciprocity quotient is open.
Local--global compatibility gives
\[
       \alpha_L(j_w a_0(\sigma))=\operatorname{toAb}(\sigma),
                                                               \tag{ITA.9}
\]
where $\alpha_L$ is the inverse-arithmetic global map of
CIP1.  CIP2 identifies every source global Artin map with
$\alpha_L$, so this proves the universally quantified
\texttt{glob} field as well.

For any other source map $a$, the quotient
$a(\sigma)/a_0(\sigma)$ is a unit by \texttt{ord\_\allowbreak eq}.
Its local idele is killed by $\alpha_L$, by the
\texttt{glob} field and (ITA.9).  CIP3 forces that unit
to be one.  This holds for every $\sigma\in W$ and
proves uniqueness.  It also checks the reciprocity sign:
a geometric Frobenius maps to a uniformizer.  No
cyclotomic-prime hypothesis has entered this proof.
\end{proof}
% END ITA4

% BEGIN ITA5
\paragraph{ITA5. Exact equivalence with the every-\texttt{Is\allowbreak Rec\allowbreak T} condition.}
\label{claim:ITA5}
For every smooth irreducible $\pi$ of $\operatorname{GL}_n(F)$,
the following three conditions are equivalent:
\begin{align*}
 P_\chi(\pi):\quad &N=0\text{ and }
     \mathrm{rec}(\operatorname{algNorm}\pi)|_I
                                  \simeq\bigoplus_i\chi_i;\\
 U_\chi(\pi):\quad &\texttt{Is\allowbreak Type\allowbreak WD}
       (\chi,\mathrm{rc}(\operatorname{algNorm}\pi))
                     \text{ for every }\mathrm{rc}\text{ satisfying }
                     \texttt{Is\allowbreak Rec\allowbreak T};\\
 E_\chi(\pi):\quad &\text{the same assertion for at least one such }
                     \mathrm{rc}.
\end{align*}
No global uniqueness theorem for the weaker source correspondence
is needed for this equivalence.
\begin{proof}
First consider a Frobenius-semisimple Weil--Deligne
representation $(R,N)$ with $N=0$ and
$R|_I\simeq\bigoplus_i\chi_i$, where every $\chi_i$
extends to the character $\chi_i^D$ of $D$.
Each occurring inertia character is fixed by conjugation
by $W$, since its extension is a character.  Consequently
every inertia-isotypic subspace is preserved by a
Frobenius lift $\Phi$.  On that subspace inertia acts
scalarly and $R(\Phi)$ is semisimple.  Choose an eigenbasis
there.  Since $W/I\simeq\mathbb Z$ is generated by
$\Phi$, the resulting basis splits the entire Weil
representation into smooth characters
\[
                            R=\bigoplus_{i=0}^{n-1}\xi_i.
                                                               \tag{ITA.10}
\]
Their inertia restrictions, with multiplicities, are exactly
the original $\chi_i$.  The inertia image is finite, so
every $\xi_i$ is smooth in the source sense.  Reordering
the basis to the original inertia list gives the literal
diagonal matrices demanded by \texttt{Is\allowbreak Type\allowbreak WD}; this
also proves its equivalence with the stated $N=0$
and inertia-isomorphism conditions.

By local reciprocity, or the \texttt{factor} field,
$\xi_i=\mu_i\circ a_0$ for smooth $\mu_i:F^\times\to
\mathbb C^\times$.  Reorder them by nondecreasing
$|\mu_i(\varpi_F)|$.  Since $|\varpi_F|<1$, this order
satisfies (ITA.8).  Let $\Pi$ be the unique irreducible
quotient of their normalized induction.  It has rank-one
cuspidal support and is therefore \texttt{Good}.  Its
parameter under the usual correspondence is (ITA.10),
with zero monodromy.  More significantly, for \emph{every}
source map $\mathrm{rc}'$, ITA4 identifies its Artin map
with $a_0$, and its \texttt{one} and \texttt{langlands}
fields give the same direct sum (ITA.10).

If $E_\chi(\pi)$ holds, apply this construction to
$\mathrm{rc}(\operatorname{algNorm}\pi)$ for its witnessing
map.  That map sends $\Pi$ to the same parameter.  Its
\texttt{inj} field, whose second argument $\Pi$ is
\texttt{Good}, gives
$\operatorname{algNorm}\pi=\Pi$ as irreducible classes.
Every other source map and the usual correspondence
therefore give the required type.  Thus
$E_\chi\Rightarrow U_\chi$ and $E_\chi\Rightarrow P_\chi$.

If $P_\chi(\pi)$ holds, construct $\Pi$ from its usual
parameter.  The bijectivity of classical local Langlands
gives $\operatorname{algNorm}\pi=\Pi$; the preceding
rank-one calculation gives $U_\chi(\pi)$ and, using
\texttt{Rec\allowbreak TExists}, $E_\chi(\pi)$.  Finally
$U_\chi\Rightarrow E_\chi$ uses exactly that same
nonemptiness theorem.  All three conditions are equivalent.

An unramified real twist of a Weil--Deligne representation
multiplies $R(w)$ by an unramified nonzero scalar and leaves
$N$ unchanged.  Hence $P_\chi(\pi)$ is also equivalent to
$N=0$ and the same inertia type for $\mathrm{rec}(\pi)$.
This checks explicitly the source's $(1-n)/2$ normalization
against the unnormalized form of the inertial-type theorem.
\end{proof}
% END ITA5

% BEGIN ITA6
\paragraph{ITA6. The published local type applies at the complete original scope.}
\label{claim:ITA6}
There is a nonzero finite-dimensional smooth irreducible complex
representation $A$ of $K_F=\operatorname{GL}_n(\mathcal O_F)$
such that for every smooth irreducible $\pi$ of
$\operatorname{GL}_n(F)$,
\[
             \operatorname{Hom}_{K_F}(A,\pi)\ne0
                         \quad\Longleftrightarrow\quad P_\chi(\pi).
                                                               \tag{ITA.11}
\]
Its kernel is open and it factors through a finite quotient of
$K_F$.
\begin{proof}[Published theorem and application]
Use the crystalline part of Matthew Emerton and Toby Gee,
\emph{Moduli stacks of \textup{\'etale} $(\varphi,\Gamma)$-modules
and the existence of crystalline lifts}, Annals of Mathematics
Studies 215, Theorem 8.2.1, with its inertial-type definition
in \S\,1.12.
\footnote{Author's text, printed p.~204 (PDF page 214),
and pp.~15--16 for conventions:
\url{https://www.math.uchicago.edu/~emerton/pdffiles/moduli.pdf}.}
For any finite extension $F/\mathbb Q_p$ and any inertial
type of rank $n$, it supplies a smooth irreducible
finite-dimensional representation $\sigma^{\mathrm{crys}}$
of $K_F$.  Its multiplicity in any irreducible smooth
representation is nonzero exactly for that inertia type
and zero monodromy.  There is no genericity hypothesis on
$\pi$, no restriction $p\nmid n$, and no exclusion of
$p=3$.  We need only existence and detection, not the
additional multiplicity bound.  This statement cites
CEGGPS, Theorem 3.7; using the book's unrestricted local
statement avoids importing that article's separate
global patching hypotheses.

Apply it to $\vartheta=\bigoplus_i\chi_i$, which is an
inertial type by ITA3.  The theorem is written over
$\overline{\mathbb Q}_p$; transport it to $\mathbb C$
through a field isomorphism.  This is an algebraic
transport of smooth representations, with no comparison
between complex and $p$-adic topologies.  Its usual
reciprocity convention is geometric, as in ITA4.
The possible unramified normalization is harmless by the
last paragraph of ITA5.  Thus its conclusion is exactly
(ITA.11), at the original $F$ and $n$.

Choose a finite basis of $A$.  The intersection of its
open stabilizers is contained in the kernel, so that
kernel is open.  It is normal, and an open subgroup
of compact $K_F$ has finite index: its cosets are an
open cover and compactness gives a finite subcover.
Thus $A$ factors through the finite quotient
$H=K_F/\ker A$.
\end{proof}
% END ITA6

% BEGIN ITA7
\paragraph{ITA7. Invariance under all algebraic automorphisms fixing the coefficient ring.}
\label{claim:ITA7}
For every $\gamma\in\operatorname{Aut}(\mathbb C/j(K))$
and every smooth irreducible $\pi$, one has
\[
                        P_\chi({}^{\gamma}\pi)
                              \Longleftrightarrow P_\chi(\pi).
                                                               \tag{ITA.12}
\]
Here scalar conjugation gives the original representation an
equivalent $\gamma$-semilinear model.  Arbitrary algebraic
automorphisms are allowed.
\begin{proof}
Suppose $P_\chi(\pi)$.  ITA5 writes
$\operatorname{algNorm}\pi$ as the ordered singleton
Langlands quotient attached to smooth characters
$\mu_0,\ldots,\mu_{n-1}$.  Their Weil restrictions
give (ITA.10), with the prescribed inertia.
Undoing the common twist, $\pi$ is the unique irreducible
quotient of normalized induction of
$\mu_i|\ |^{(n-1)/2}$.  On the upper triangular Borel,
the inducing transformation character, including the
normalizing modulus, is
\[
       \prod_{i=0}^{n-1}\mu_i(a_i)|a_i|^{n-1-i}.
                                                               \tag{ITA.13}
\]
Indeed the exponent of $a_i$ in the half-modulus is
$(n-1-2i)/2$; adding $(n-1)/2$ gives $n-1-i$.
These are integer exponents, so every value of the norm
factor in (ITA.13) is rational and fixed by $\gamma$.

Apply $\gamma$ pointwise to the functions in this induced
representation.  It is a semilinear bijection commuting
with right translation; it preserves smoothness because
stabilizers are unchanged.  Formula (ITA.13) shows that
the resulting induction is the normalized induction of
$(\gamma\mu_i)|\ |^{(n-1)/2}$.  The relation (ITA.8)
is preserved: $|\ |$ is rational-valued, so an equality
after applying $\gamma$ would give the same equality
before applying it.  Semilinear transport preserves
submodules, quotients and the unique irreducible cosocle.
Thus $\operatorname{algNorm}({}^{\gamma}\pi)$ is the
singleton quotient for $\gamma\mu_i$.

The rank-one and Langlands formulas give its parameter
$\bigoplus_i\gamma\xi_i$, with $N=0$.  Its inertia
characters are $\gamma\chi_i$.  Every value of $\chi_i$
is in $j(O)$ by ITA3, hence is fixed by $\gamma$.
This proves the forward implication in (ITA.12).
Apply the same argument to $\gamma^{-1}$ for the reverse
implication.  The cancellation in (ITA.13) is necessary:
one cannot assume that an arbitrary $\gamma$ fixes a
chosen positive square root of the residue cardinality.
\end{proof}
% END ITA7

\subsection{Descent of the detecting representation to the original DVR}

% BEGIN ITA8
\paragraph{ITA8. A stable lattice followed by restriction of scalars.}
\label{claim:ITA8}
There is a finite free, nonzero $O$-module $B$ and an action
$R:K_F\to\operatorname{Aut}_O(B)$ with open kernel such
that, with $B_j=\mathbb C\otimes_{O,j}B$,
\[
              \operatorname{Hom}_{K_F}(B_j,\pi)\ne0
                             \quad\Longleftrightarrow P_\chi(\pi)
                                                               \tag{ITA.14}
\]
for every smooth irreducible $\pi$.  The output coefficient
ring is exactly $O$.
\begin{proof}
Choose a complex basis of the finite representation $A$ of
ITA6 and its finite quotient $H$.  By ITA1 the finitely
many entries of its matrices, pulled back by $\iota^{-1}$,
belong to a finite extension $E/K$ inside the original
$\overline{\mathbb Q}_\ell$.  They define a representation
$A_E:H\to\operatorname{GL}_d(E)$ whose extension by
$\iota$ is the given $A$.  Its group identities hold over
$E$ because that field embeds in $\mathbb C$.

Let $O_E=\{x\in E:|x|_\ell\leq1\}$.  The norm-ball
argument of WRP2 proves that this is a DVR finite over
$\mathbb Z_\ell$, containing $O$.  In particular it is
finite over $O$: any finite set of $\mathbb Z_\ell$-module
generators also generates it over $O$.  Set
\[
                      \mathcal L=\sum_{h\in H}A_E(h)O_E^d
                                      \ \subset E^d.          \tag{ITA.15}
\]
The sum is finite and is preserved by every $A_E(h)$,
with equality because its inverse also belongs to $H$.
It contains $O_E^d$ and spans $E^d$ over $E$.
It is a finite torsion-free $O_E$-module, hence free of
rank $d$ over that DVR.  Regard the \emph{same} module as
an $O$-module $B=\operatorname{Res}_{O_E/O}\mathcal L$.
It is finite and torsion-free over the original DVR,
so is finite free.  Inverting $\ell$ gives $E^d$:
every element of $E$ becomes integral after multiplication
by a sufficiently high power of $\ell$.  Thus
\[
              \operatorname{rank}_O B=d[E:K]>0,
       \qquad K\otimes_O B\simeq\operatorname{Res}_{E/K}E^d.
                                                               \tag{ITA.16}
\]
The action $R$ on $B$ still factors through $H$.

Since $E/K$ is separable, scalar extension decomposes as
\[
 B_j\simeq
 \bigoplus_{\eta:E\hookrightarrow\mathbb C,\ \eta|_K=j}
                  \bigl(\mathbb C\otimes_{E,\eta}E^d\bigr).
                                                               \tag{ITA.17}
\]
For clarity, choose a primitive element of $E/K$.
Its minimal polynomial has distinct roots in $\mathbb C$;
the Chinese remainder theorem gives
$\mathbb C\otimes_K E\simeq\prod_\eta\mathbb C$
by $z\otimes x\mapsto(z\eta(x))_\eta$.  Applying this
to $E^d$ gives (ITA.17) and shows directly that each
summand's action is obtained by applying $\eta$ to the
original matrices.

Each $\eta$ extends, relative to the initial embedding
$\iota|_E$, to some $\gamma\in
\operatorname{Aut}(\mathbb C/j(K))$.  The usual extension
of an algebraic-field embedding to an algebraic closure
gives an embedding of $\mathbb C$ into itself over
$j(K)$.  It is onto: its image is algebraically closed,
and $\mathbb C$ is algebraic over $j(K)$.  Thus the
summand in (ITA.17) is ${}^{\gamma}A$.

Semilinearly transporting a map gives
\[
 \operatorname{Hom}_{K_F}({}^{\gamma}A,\pi)\ne0
 \ \Longleftrightarrow\
 \operatorname{Hom}_{K_F}(A,{}^{\gamma^{-1}}\pi)\ne0.
\]
By ITA6 and ITA7 this is equivalent to $P_\chi(\pi)$,
for \emph{every} summand of (ITA.17).  The sum is finite
and nonempty.  Therefore its Hom space is nonzero exactly
when $P_\chi(\pi)$, proving (ITA.14).
The extension $O_E$ was used only to construct the lattice;
restriction of scalars makes every final matrix an
$O$-matrix.  No assertion that $A$ itself descends
irreducibly to $O$, and no division by $|H|$ inside $O$,
is required.
\end{proof}
% END ITA8

% BEGIN ITA9
\paragraph{ITA9. The contragredient gives exactly the source's column equation.}
\label{claim:ITA9}
Choose an $O$-basis $e_0,\ldots,e_{d_B-1}$ of $B$ and let
\[
                \tau(k)=R(k^{-1})^{\mathsf T}
                       \in\operatorname{GL}_{d_B}(O).          \tag{ITA.18}
\]
This is an integral representation with open kernel.  For
every complex representation $\rho$ of $K_F$ on $M$,
the vector families satisfying the source's column equation
are naturally identified with
$\operatorname{Hom}_{K_F}(B_j,M)$.  In particular their
nonvanishing for a smooth irreducible $\pi$ is equivalent
to $U_\chi(\pi)$ of ITA5.
\begin{proof}
The module underlying (ITA.18) is $B^*=\operatorname{Hom}_O(B,O)$,
with the dual basis.  It is preserved because $R(k)$ and
$R(k^{-1})$ preserve $B$.  Matrix multiplication gives a
homomorphism, and the kernel equals that of $R$.
To a family $F_a\in M$ associate the unique linear map
$f:B_j\to M$ with $f(e_a)=F_a$.  For every $k,a$,
\[
 \sum_b j(\tau(k)_{ab})\rho(k)F_b
       =\rho(k)f\bigl(R(k^{-1})e_a\bigr).                   \tag{ITA.19}
\]
This equals $F_a$ for all $a$ exactly when
$fR(k)=\rho(k)f$.  Thus the family's equations are
precisely equivariance, and the family is nonzero exactly
when $f$ is nonzero.  Equivalently the family represents
an invariant tensor in $B_j^*\otimes_{\mathbb C}M$.
Apply ITA8 and ITA5 for the last assertion.  This explicitly
checks the inverse and transpose; using $R$ itself in
the column equation would detect its dual instead.
\end{proof}
% END ITA9

\subsection{The product of local types for arbitrary isotypic modules}

% BEGIN ITA10
\paragraph{ITA10. What the literal \texttt{Iso\allowbreak Sum} hypothesis supplies.}
\label{claim:ITA10}
Let $G$ be a group, $\pi$ an irreducible complex representation,
and $N$ a representation generated by images of equivariant
maps $\pi\to N$.  Then $N$ is a direct sum of copies of
$\pi$.  Every subrepresentation of such an $N$ is again
a sum of copies of $\pi$.  If $N\ne0$ and $C$ is a
finite-dimensional representation of a subgroup $K\subset G$,
then
\[
                 \operatorname{Hom}_K(C,N)\ne0
                   \quad\Longleftrightarrow\quad
                 \operatorname{Hom}_K(C,\pi)\ne0.            \tag{ITA.20}
\]
\begin{proof}
A nonzero map from $\pi$ is injective, and its image is
simple.  Choose a maximal family of these simple
submodules whose sum is direct.  If its sum $N_0$ is not
all of $N$, some generating simple submodule $P$ is not
contained in $N_0$.  Its intersection with $N_0$ is a
submodule of the simple $P$, hence is zero.  Adding $P$
contradicts maximality.  This proves the first statement.

For completeness the same argument is the usual
semisimple-module argument for submodules.  A sum of
simple modules is semisimple: given a submodule $Y$, take
a maximal direct sum of simple summands disjoint from
$Y$; its sum with $Y$ must be all of $N$, because a
simple generating summand not contained in that sum has
zero intersection with it and could be added.  Hence
every submodule has a complement.  Project the original
simple generating summands to $Y$ along that complement.
Their nonzero images are simple and isomorphic to
$\pi$, and generate $Y$.  This proves the asserted
isotypy of submodules, without a finite length assumption.

Finally a linear map from finite-dimensional $C$ into a
direct sum has image in a finite subsum: apply the
definition of direct sum to the images of a basis.
If a nonzero equivariant map $C\to N$ exists, one
coordinate projection to a copy of $\pi$ is nonzero
and equivariant.  Conversely, since $N\ne0$, there is
a copy of $\pi$ in $N$ into which any nonzero map
$C\to\pi$ can be composed.  This proves (ITA.20).
\end{proof}
% END ITA10

% BEGIN ITA11
\paragraph{ITA11. Simultaneous detection in the original product representation.}
\label{claim:ITA11}
For each $u\in S$ take the $O$-module $B_u$ of ITA8,
its action $R_u$, and its dual action $\tau_u$ from ITA9.
Put
\[
 B_S=\bigotimes_{u\in S,O}B_u,\qquad
 R_S(k)=\bigotimes_{u\in S}R_u(k_u),\qquad
 \tau_0=R_S^*,\qquad m=\prod_{u\in S}\operatorname{rank}_O B_u.
                                                               \tag{ITA.21}
\]
Choose the tensor basis and any fixed identification of its
finite index set with $\{0,\ldots,m-1\}$.  For every
$M,\rho,(\pi_u)$ satisfying the exact hypotheses before
(ITA.6), its left side holds if and only if
$P_{\chi_u}(\pi_u)$ holds for every $u$.
\begin{proof}
By ITA9 the left side is the nonvanishing of
\[
 \bigl((B_S)_j^*\otimes M\bigr)^{K_S}
                \simeq\operatorname{Hom}_{K_S}((B_S)_j,M).
                                                               \tag{ITA.22}
\]
Enumerate the finite set $S$.  Starting with $N_0=M$,
apply successively the functors
\[
                    N\longmapsto
                    \operatorname{Hom}_{K_u}((B_u)_j,N)
                      \simeq((B_u)_j^*\otimes N)^{K_u}.
                                                               \tag{ITA.23}
\]
The remaining groups act on the $N$ factor.  These actions
are well defined and commute with taking $K_u$-invariants,
since the original factors of $G_S$ commute.  For each
unprocessed factor $G_v$, the previous module is
$\pi_v$-isotypic: this is true at the start by the
literal \texttt{Iso\allowbreak Sum} hypothesis and ITA10.  The next
module is a $G_v$-submodule of the finite direct sum
$(B_u)_j^*\otimes N$, and so is again $\pi_v$-isotypic
by ITA10.

Whenever the previous module $N$ is nonzero, (ITA.20)
and (ITA.14) say that (ITA.23) is nonzero exactly when
$P_{\chi_u}(\pi_u)$ holds.  A zero module stays zero
under all subsequent steps.  Induction therefore gives
nonzero output exactly when every one of these local
conditions holds.  Successive finite tensor--Hom
identifications identify that output with (ITA.22):
equivariance in each coordinate is precisely equivariance
under their finite product.  This proves the assertion.
No decomposition of $M$ as an external tensor product,
no scalar-endomorphism version of Schur's lemma, and no
admissibility of $M$ have been used.

If $S$ is empty, (ITA.21) means $B_S=O$, $m=1$ and
the trivial action.  Then (ITA.22) is $M\ne0$, and
the right side is the empty universal condition, so the
same conclusion holds in this case as well.
\end{proof}
% END ITA11

% BEGIN ITA12
\paragraph{ITA12. The complete original integral \texttt{Type\allowbreak For\allowbreak Alt}.}
\label{claim:ITA12}
The homomorphism $\tau_0$ in (ITA.21), over the unchanged
ring $O=c.O$, satisfies the entire original
$\texttt{Type\allowbreak For\allowbreak Alt}(r,\iota,\tau_0)$.
\begin{proof}
First the intersection of the coordinate inverse images
of $\ker R_u$ is open in $K_S$.  Every element of that
intersection acts trivially on the tensor module in
(ITA.21); hence it belongs to $\ker\tau_0$.
A subgroup containing an open subgroup is open, since
it is a union of its open cosets.  Its image under the
injective scalar map $j$ has exactly the same kernel:
a matrix becomes the identity over $\mathbb C$ if and
only if it was the identity over $O$.  Thus the
complexified kernel in the source's \texttt{Is\allowbreak Type\allowbreak Alt}
field is open.

Fix an arbitrary family $t'$ satisfying the source's
$r.\texttt{Is\allowbreak Teich}$, not necessarily the family used
when choosing the local types.  ITA2 gives equality
$t'_i=t_i$ as homomorphisms for every $i$.  Consequently
the local characters (ITA.5) coincide as actual maps
on the original inertia subgroups.  For every allowed
$M,\rho,(\pi_u)$, ITA11 proves the left side of
(ITA.6) equivalent to all the local conditions
$P_{\chi_u}(\pi_u)$.  ITA5 turns each of them into
the universally quantified \texttt{Is\allowbreak Rec\allowbreak T} condition.
The coordinate formula in ITA9 is exactly the sum in
the source definition.  This gives every field of
\texttt{Is\allowbreak Type\allowbreak Alt} for the fixed $t'$, and $t'$ was
arbitrary.  Therefore it proves \texttt{Type\allowbreak For\allowbreak Alt}.

The choice order is the required one.  The data
$s,n,V,Q,c,r,\iota$ are fixed first; then the unique
Teichm\"uller lifts, finitely many local types and
lattices, their $O$-bases, and their tensor product
are chosen once.  The resulting $m,\tau_0$ precede
every later lift family, local correspondence, module,
and representation in (ITA.6).  The only coefficient
extension was internal to (ITA.15), followed by
restriction of scalars.  All output matrices lie in
the original $O$ and are complexified through the
specified $\iota$.

This proves the existential endpoint of
\path{TTypeEx.exists_integral_type_for_residual_rep}
under its original foundational inputs.  Its rank is
every original $n\geq1$.  The argument includes the
allowed original coefficient prime $\ell=3$, every
local residue prime, repeated inertia characters,
and the empty type-place set.  It imposes no
genericity, regularity, tame-ramification, large-prime
or prime-to-rank hypothesis.
\end{proof}
% END ITA12

% BEGIN ITA13
\paragraph{ITA13. The precise input supplied to the existing Hecke-data assembly.}
\label{claim:ITA13}
For the same original $V,Q,c,r,\iota$ and every actual
universal character ring $\Lambda$, ITA12 supplies the
formerly missing type input in OTW13.  Applying OTW13
with this $m,\tau_0$ gives the original
$\texttt{Unitary\allowbreak Hecke\allowbreak Data}(Q,\iota,\Lambda)$, and hence
its \texttt{Hd\allowbreak Ex} existential conclusion.
\begin{proof}
The integral representation and all its maps have the
types required by OTW13; ITA12 gives precisely its
\texttt{type} hypothesis.  OTW12 in that assembly
provides the one character $\theta$ together with
\texttt{theta} and \texttt{is\allowbreak Wt}, using this same
representation.  CAP17 provides the residual
Hecke-eigenvalue function $a$ and its \texttt{res}
field on the same original operators.  Thus
\[
        (m,\tau_0,\texttt{type},\theta,
                   \texttt{theta},\texttt{is\allowbreak Wt},a,\texttt{res})
\]
has exactly the eight fields of
\texttt{Unitary\allowbreak Hecke\allowbreak Data}.  This is direct substitution
into the already reconstructed OTW13 proof, with no
new coefficient extension or later alteration of
$\tau_0$, $r$ or $\Lambda$.  The conclusion is the
Hecke-data supplier only; the later common-ideal
and automorphy conclusions require their own
separate arguments.
\end{proof}
% END ITA13

\subsection{Proof boundary and exact source dictionary}

The following locators are in the pinned source commit
\begin{center}\texttt{7cbf7a0c32fca1badf05931879a37c3f8a2b03b4}.\end{center}
The companion metadata records their full file and range hashes.
\begin{center}\small
\begin{tabular}{p{0.26\linewidth}p{0.38\linewidth}p{0.26\linewidth}}
\hline
Original object or endpoint & Source file and declaration & Written proof \\
\hline
Actual coefficient subring & \texttt{Dirichlet\allowbreak Prime\allowbreak To\allowbreak LPart}\allowbreak\texttt{Mul\allowbreak LPower\allowbreak Part.\allowbreak Coef} & ITA1, ITA8 \\
Every lift family & \path{ArithmeticPointClassicality.Resid.IsTeich} & ITA2--ITA3 \\
Actual split type places & \path{ContinuousResP.HidaSetting}, \texttt{Lev.\allowbreak pr} & ITA3--ITA4 \\
Zero-monodromy diagonal inertia & \texttt{Distinct.\allowbreak Is\allowbreak Type\allowbreak WD} & ITA5 \\
Weaker local correspondences & \texttt{Is\allowbreak Rec\allowbreak TOf\allowbreak Is\allowbreak Rec.\allowbreak Laws\allowbreak T}, \texttt{Is\allowbreak Rec\allowbreak T} & ITA4--ITA5 \\
Conjugation and integral descent & \texttt{Type\allowbreak Conj}, \texttt{Type\allowbreak Loc}, \texttt{PType\allowbreak Ex.\allowbreak te\_\allowbreak loc} & ITA7--ITA9 \\
Arbitrary product modules & \texttt{Coefficient\allowbreak Rep}\allowbreak\texttt{Trivial\allowbreak On\allowbreak Unramified.\allowbreak Iso\allowbreak Sum} & ITA10--ITA11 \\
Full original type & \path{UnitaryHeckeData.TypeForAlt}, \texttt{TType\allowbreak Ex} & ITA12 \\
Original Hecke-data record & \texttt{Unitary\allowbreak Hecke\allowbreak Data}, \texttt{Hd\allowbreak Ex} & ITA13, OTW13 \\
\hline
\end{tabular}
\end{center}

The new proof simplifies the source's \texttt{Nonzero\allowbreak Types}
and generic integral-annihilator descent at the existential
endpoint actually required.  It does not claim that the chosen
lattice basis equals the source's auxiliary Hilbert-choice
basis, and no such equality is required by that endpoint.
The local type theorem is an explicit matched classical
input, not a new proof of all of Bushnell--Kutzko or
Schneider--Zink theory.  Likewise \texttt{Rec\allowbreak TExists}
is retained at its original foundation status.  The
mathematical progress here is the full coefficient,
normalization, universal-quantifier, duality and product
argument yielding ITA12 and its exact OTW13 consequence.


\section{The actual compact-unitary automorphic decomposition}
\label{sec:automorphic-decomposition-v43}

The preceding classical-parts theorem preserves exact level and ordinary
conditions inside a supplied decomposition.  Here we construct the
underlying automorphic decomposition on the original definite unitary
quotient.  The endpoints are the full original \texttt{form\allowbreak U} and
\texttt{decomp\allowbreak U}.  These endpoints are independent of the arithmetic
construction of a Galois representation for an automorphic constituent.
The latter construction remains necessary to obtain a full
\texttt{Part\allowbreak Alt\allowbreak Alt} and \texttt{Automorphic\allowbreak Hecke\allowbreak Set}.

Fix the original \(\ell,L,s,n,V,Q,c,\tau_0,\iota\), with
\(n\ge1\).  Write \(G=U_{\mathbb A}\),
\(\Gamma=U(L^+)\), \(G_f=U_f\), and \(G_\infty=U_\infty\),
with the literal identity Hermitian form and embeddings of FNL2.  The
source \(\mathcal A=\mathrm{autU}(L,n)\) consists of continuous,
left \(\Gamma\)-invariant complex functions on \(G\), smooth under
\(G_f\), and finite under \(G_\infty\).  Its action is right
translation.  An original automorphic representation is a nonzero
irreducible invariant algebraic subspace of this very function space.

For the specified weight and basis, let
\[
 D=\bigotimes_{a\in\operatorname{Emb}_\ell}
             \nabla_{\lambda(a)}(\mathbf C),\qquad
 E:\mathbf C^d\xrightarrow{\sim}D,
\]
and write \(R=\mathrm{infRep}(\iota,\lambda)\).  The letter \(E\)
is the exact source basis equivalence, not a newly chosen realization.
Let
\[
 p:G\longrightarrow G^a\times G^\ell,
 \qquad
 p(g)=\bigl(\mathrm{away}(g),(\mathrm{cLoc}_{w(v)}(g))_{v\mid\ell}\bigr).
\]
For an algebraic coefficient function \(F\), the exact source scalarization is
\begin{equation}\label{eq:adp-toform}
 \mathcal T_F(\lambda^*)(g)
       =\lambda^*\bigl(R(g^{-1})E(F(p(g)))\bigr),\qquad
                  \lambda^*\in D^*.
\end{equation}
The symbols \(\lambda\) and \(\lambda^*\) respectively denote the fixed
weight and an arbitrary linear functional; there is no dual-weight change
implicit in this notation.

\subsection{Algebraic semisimplicity on the original quotient}

% BEGIN ADP1
\paragraph{ADP1. A compact fundamental set and the actual positive invariant form.}
\label{claim:ADP1}
The quotient \(\Gamma\backslash G\) is compact and has a positive,
finite, right-invariant measure.  On \(\mathcal A\) the formula
\[
       \langle f,h\rangle
       =\int_{\Gamma\backslash G}\overline{f(g)}h(g)\,d\mu(g)
\]
is a positive definite, right-invariant Hermitian form.

Indeed choose an open compact subgroup \(K_f\le G_f\).  Such subgroups
are obtained by imposing integral matrices at almost all places and open
compact congruence conditions at the finitely many remaining places; the
unitary equations cut out closed subgroups.  FNL3 gives finitely many
\(t_i\) such that every \(g\) is \(\gamma t_i z k\), with
\(z\in G_\infty\), \(k\in K_f\).  The group \(G_\infty\) is a
finite product of compact identity-form unitary groups: at every real
place of \(L^+\), the two complex embeddings of \(L\) are conjugate,
and the equation is \(g^*g=1\).  Thus
\(\bigcup_it_iG_\infty K_f\) is a compact fundamental set.

The rational subgroup is discrete in these same adelic coordinates,
and FV5--FV8 give unimodularity and the invariant quotient measure for
this unitary group.  Compactness makes its measure finite; a nonempty
open subset has positive measure by the local Haar quotient formula.
Continuous functions on the compact quotient are bounded, so the
integral is defined.  Right invariance proves invariance of the form.
If the squared norm is zero and \(f(g_0)\ne0\), continuity supplies
a neighborhood of \(g_0\) on which \(|f|\) is bounded below by a
positive constant.  Its image in the quotient has positive measure,
contradicting the zero integral.  Hence the form is definite on the
actual functions, without passing to almost-everywhere equivalence classes.
% END ADP1

% BEGIN ADP2
\paragraph{ADP2. Finite multiplicity from the exact finite class representatives.}
\label{claim:ADP2}
Let \(S\) be a finite-dimensional complex representation of
\(G_\infty\), and let
\[
 \mathcal E_S=\operatorname{Hom}_{G_\infty}(S,\mathcal A).
\]
For every open \(K_f\le G_f\), the space \(\mathcal E_S^{K_f}\)
is finite-dimensional.  If \(t_1,\ldots,t_r\) are the representatives
from FNL3, the evaluation map
\[
 \mathcal E_S^{K_f}\longrightarrow (S^*)^r,
       \quad T\longmapsto\bigl(v\mapsto T(v)(t_i)\bigr)_i
\]
is injective.

If all evaluations vanish, write \(g=\gamma t_i z k\) in the exact
FNL3 order.  Left rational invariance removes \(\gamma\); right
\(K_f\)-invariance removes \(k\); equivariance gives
\(T(v)(t_i z)=T(zv)(t_i)=0\).  Thus \(T=0\), proving the bound
\(\dim\mathcal E_S^{K_f}\le r\dim S\).  Neither independence of the
representatives nor triviality of their stabilizers is required.

Each \(T\in\mathcal E_S\) is fixed by an open compact subgroup of
\(G_f\): take a finite basis of \(S\), intersect the open finite
stabilizers of its images, and choose an open compact subgroup inside
that intersection.  Every compact subgroup has finite orbits on each
element of \(\mathcal E_S\), since an open stabilizer has finite
index in a compact group.  These are the exact admissibility and
finite-orbit assertions needed below.
% END ADP2

% BEGIN ADP3
\paragraph{ADP3. Complete reducibility in each actual multiplicity space.}
\label{claim:ADP3}
Each \(\mathcal E_S\) is an algebraic sum of irreducible
\(G_f\)-subrepresentations.

Choose any basis \(s_1,\ldots,s_d\) of \(S\) and give
\(\mathcal E_S\) the Hermitian form
\[
       (T,T')_S=\sum_{j=1}^d
                         \langle T(s_j),T'(s_j)\rangle.
\]
ADP1 makes it positive definite and \(G_f\)-invariant.  For an arbitrary
\(T\), ADP2 gives an open compact \(K_f\) fixing it, with finite
dimensional fixed space and finite orbits on all other vectors.  The
abstract argument in PKT5--PKT6 now applies with precisely this
\(G_f,\mathcal E_S,K_f\) and this Hermitian form.  It expresses
\(T\) as a finite sum of vectors in irreducible invariant subspaces.
All hypotheses of that argument have just been produced here; none of
the separate geometric multiplicity inputs used in PKT4 are imported.
% END ADP3

% BEGIN ADP4
\paragraph{ADP4. Algebraic decomposition of the complete automorphic function space.}
\label{claim:ADP4}
Every \(f\in\mathcal A\) belongs to the algebraic sum of the original
automorphic representations in \(\mathcal A\).

Its \(G_\infty\)-orbit spans a finite-dimensional invariant space by
the definition of \(\mathcal A\).  The form of ADP1 makes this space
unitary, so successive invariant orthogonal complements decompose it
into finite-dimensional irreducibles \(S\).  For each such \(S\),
ADP3 decomposes its multiplicity space \(\mathcal E_S\) under \(G_f\).
It remains to check that, for an irreducible subspace
\(U\subset\mathcal E_S\), its evaluation image
\[
       \operatorname{ev}(S\otimes U)\subset\mathcal A,
          \qquad s\otimes T\longmapsto T(s),
\]
is an irreducible \(G_\infty\times G_f=G\)-representation.

The algebra spanned by the irreducible unitary action on \(S\) is
\(\operatorname{End}_{\mathbf C}(S)\), by the finite-dimensional
matrix-algebra argument written in PKT10.  Therefore an invariant
subspace of \(S\otimes U\) is of the form \(S\otimes U'\): in a
basis of \(S\), the matrix units extract each coefficient of a tensor
and insert it in every coordinate.  Stability under \(G_f\) makes
\(U'\) invariant, so irreducibility of \(U\) gives either zero or the
whole tensor product.  The evaluation map is equivariant and nonzero
because a nonzero map \(T\in U\) has some \(T(s)\ne0\).  Its kernel
is therefore zero, and its image is irreducible.  Each element of that
image is still an actual automorphic function.  This proves all fields
of the original \(\mathrm{IsAutRepU}\) for the image, and proves the
claimed algebraic decomposition.  In particular no Hilbert-space closure
is substituted for the required finite sum.
% END ADP4

% BEGIN ADP5
\paragraph{ADP5. Finite globally equivariant projections on a prescribed finite space.}
\label{claim:ADP5}
For every finite-dimensional subspace \(M\subset\mathcal A\), there
are finitely many original automorphic representations \(\sigma_i\)
and \(G\)-equivariant linear maps \(P_i:\mathcal A\to\mathcal A\)
such that \(P_i\mathcal A\subset\sigma_i\) and
\(\sum_iP_i f=f\) for all \(f\in M\).

ADP4 says that \(\mathcal A\), as a module over \(\mathbf C[G]\),
is a sum of simple submodules.  Choose a maximal direct family of simple
submodules by the usual union-of-chains argument.  If its sum were
proper, a simple submodule from the spanning family would not be
contained in it; its intersection with the sum is then zero, contradicting
maximality.  Hence this family is a direct-sum decomposition of
\(\mathcal A\).  Coordinate projections are \(G\)-equivariant.
Only finitely many coordinates occur in a finite basis of \(M\).
Select precisely those coordinates.  Their projections have the required
properties on every vector of \(M\).  Multiplicities cause no difficulty:
distinct summands may be isomorphic representations, and are kept as
distinct subspaces of the original function space.
% END ADP5

\subsection{The exact coefficient-to-function dictionary}

% BEGIN ADP6
\paragraph{ADP6. The finite projection and the prescribed infinity representation.}
\label{claim:ADP6}
The map \(p\) is a homomorphism, is surjective already on \(G_f\),
and has kernel \(G_\infty\).  The representation \(R\) is trivial on
\(G_f\), has continuous complex matrix coefficients, and satisfies
\[
 R(\gamma)E(v)=E\bigl(\xi(p(\gamma)_\ell)v\bigr)
                       \quad(\gamma\in\Gamma).
\]

The splitting of all selected places above \(\ell\), the exact local
sections, and the independent finite-coordinate adjustment are FNL4--FNL6.
They give surjectivity without changing the away coordinates.  If all
displayed coordinates of \(p(g)\) are one, every other finite coordinate
is one as well: at a selected split pair the opposite coordinate is the
inverse conjugate transpose of the displayed one.  Thus the kernel is
exactly \(G_\infty\).

Each tensor factor of \(R\) is the specified polynomial representation
evaluated on the complex infinite coordinate, so it is trivial on
finite-only elements.  Its matrix coefficients are polynomials in matrix
entries times an integral power of the nonzero determinant; they are
continuous, including negative determinant powers.  Finite tensor products
preserve this assertion.  The complex embedding of \(L\) indexing a
coordinate is fixed; this argument never regards \(\iota\) as a
continuous map on \(\overline{\mathbf Q}_\ell\).

For a rational matrix \(\gamma\), applying its prescribed embedding
to each entry of the selected local component is exactly its infinite
component at that embedding.  Apply the same polynomial representation
and tensor the equalities.  The original \(\mathrm{IsBasis}\) identity
for \(E,\xi\) gives the displayed rational compatibility.  This is the
literal dictionary for \(R,p,E\).
% END ADP6

% BEGIN ADP7
\paragraph{ADP7. Scalarization, its exact translation rule, and automorphy.}
\label{claim:ADP7}
For every coefficient function \(F\), the map \(\mathcal T_F\) of
\eqref{eq:adp-toform} is linear in \(F\) and in \(\lambda^*\), is
injective as a map of coefficient functions, and obeys
\[
 R_x\mathcal T_F(\lambda^*)
  =\mathcal T_{F(\,\cdot\,p(x))}
                    (\lambda^*\circ R(x^{-1})).
\]
If \(F\) has the original rational equivariance and an open finite
stabilizer, then every \(\mathcal T_F(\lambda^*)\) belongs to
\(\mathcal A\).

Linearity follows directly from the formula.  If every scalarization
vanishes, surjectivity of \(p\), invertibility of \(R(g)\), and
separation by linear functionals on finite-dimensional \(D\) give
\(E(F(y))=0\) for every \(y\), hence \(F=0\).  The translation
identity is the equation
\(R((gx)^{-1})=R(x^{-1})R(g^{-1})\), with the order as displayed.

By ADP6, rational equivariance says
\(E(F(p(\gamma)y))=R(\gamma)E(F(y))\).  Substitution proves left
\(\Gamma\)-invariance of the scalarization.  Right translation by
an open finite stabilizer fixes it, since \(R\) is trivial on \(G_f\).
Its \(G_\infty\)-translates lie in the image of the finite-dimensional
space \(D^*\), proving infinity finiteness.

For continuity, use the continuous finite projection
\(G=G_f\times G_\infty\to G_f\) from FNL2.  An open finite
stabilizer fixes \(F(p(g))\) locally in \(g\), while the dependence
on the infinite component is a continuous matrix coefficient of \(R\)
by ADP6.  Locally the scalarization is therefore a continuous matrix
coefficient with a fixed vector.  These are precisely all the original
\(\mathrm{IsAutFormU}\) conditions.
% END ADP7

% BEGIN ADP8
\paragraph{ADP8. Inverse scalarization on the original functions.}
\label{claim:ADP8}
Let \(\Phi:D^*\to\mathcal A\) be linear and suppose that
\[
 R_z\Phi(\lambda^*)
      =\Phi(\lambda^*\circ R(z^{-1}))\quad(z\in G_\infty).
\]
There exists a unique coefficient function \(F\) with the original
rational equivariance and \(\mathcal T_F=\Phi\).

Finite-dimensional biduality gives a unique \(\Psi(g)\in D\) with
\(\lambda^*(\Psi(g))=\Phi(\lambda^*)(g)\) for every functional.
The infinity relation is exactly
\(\Psi(gz)=R(z^{-1})\Psi(g)\).  Hence \(R(g)\Psi(g)\) is constant
on each fiber of \(p\), by ADP6.  Define
\[
                    F(p(g))=E^{-1}(R(g)\Psi(g)).
\]
This is independent of the lift \(g\).  Left rational invariance of
\(\Phi\) gives \(\Psi(\gamma g)=\Psi(g)\), hence
\(E(F(p(\gamma)y))=R(\gamma)E(F(y))\).  This is the exact rational
coefficient condition by ADP6.  Substitution gives
\(\mathcal T_F(\lambda^*)(g)=\Phi(\lambda^*)(g)\).  Uniqueness
follows from ADP7.  No coefficient field, basis equivalence, or
algebraic representation was changed in this inversion.
% END ADP8

% BEGIN ADP9
\paragraph{ADP9. Full original \texttt{form\allowbreak U}.}
\label{claim:ADP9}
The five clauses of the original \texttt{form\allowbreak U} all hold for the given
\(E,\xi,\iota\) and original projection \(p\): injectivity on the
rationally equivariant functions, automorphy under the original open
finite-stabilizer hypothesis, the exact right-translation identity,
surjectivity of the finite projection, and inverse scalarization for every
equivariant \(\Phi\).

ADP6 gives the finite projection and the exact rational dictionary.
ADP7 gives the first three function assertions, and ADP8 gives the final
surjectivity assertion, including the original function space in which
the inverse lies.  The argument works for any supplied finite-dimensional
basis equivalence, including a zero-dimensional coefficient space, and
uses no openness condition on \(\tau_0\) for this endpoint.  That
condition enters only when the level laws of a vector are used to produce
an open scalar-function stabilizer in ADP10.
% END ADP9

\subsection{Transport of finite automorphic projections}

% BEGIN ADP10
\paragraph{ADP10. The actual finite-level laws give simultaneous smoothness.}
\label{claim:ADP10}
Assume \(\ker\tau_0\) is open, and let \(F=(F_i)\in H.\mathrm{Big}\)
be fixed at one original away level \(K\) and one original Iwahori
depth \(b\).  Every scalarization of every \(F_i\) is an element of
\(\mathcal A\).

Choose a sufficiently small finite adelic neighborhood of one satisfying
the following finitely many conditions: its selected \(\ell\)-part lies
in \(I_1(b)\); its away part has the original level expression; and
its type coordinates lie in \(\ker\tau_0\).  Such a neighborhood
exists at every stated depth: the local principal congruence subgroup
\[
 \ker\bigl(\operatorname{GL}_n(O_{F_v})\longrightarrow
              \operatorname{GL}_n(O_{F_v}/\mathfrak m_{F_v}^{b+1})\bigr)
                         \ \subset I_1(b)
\]
is open.  To check the inclusion, eliminate a matrix congruent to the
identity modulo \(\mathfrak m_{F_v}^{b+1}\) into upper-unipotent,
diagonal and lower-unipotent factors.  Each pivot is congruent to one
and hence is a unit; each subtraction and division preserves the same
congruence.  The three factors therefore lie in the original
\(N\), \(T_c(b)\), and \(N^-(b)\), respectively.  The defining
Iwahori factorization puts their product in \(I_1(b)\).  The finite
product over the selected places is open as well.  Combine this with
FNL8--FNL9, continuity of the finite projections, openness of
the given \(K_q\)'s, and the assumed open kernel.  Outside these
finitely many places impose the usual integral conditions.  Thus every
element \(k\) of this neighborhood has
\(p(k)=(u,b')\), with \(u\in H.\mathrm{lev}(K)\),
\(\tau_C(u)=1\), and \(b'\in I_1(b)\).

The vector-valued level equation
\(\tau_C(u)F(y(u,1))=F(y)\) now says that every coordinate \(F_i\)
is fixed by this away translation.  Its \(I_1(b)\)-fixedness supplies
the local translation as well, so \(F_i(yp(k))=F_i(y)\).
Each scalar coordinate therefore has an open finite stabilizer.  Its
rational coefficient condition comes from \(H.\mathrm{Big}\), and
ADP7 applies.  This establishes a common neighborhood before the
coordinate or dual functional is selected; it does not assume separate
coordinate smoothness from vector-valued level invariance alone.
% END ADP10

% BEGIN ADP11
\paragraph{ADP11. Equivariant projections preserve the exact coefficient equations.}
\label{claim:ADP11}
For \(F\) as in ADP10, there are finitely many original automorphic
representations \(\sigma_t\) and vectors \(G_t\in H.\mathrm{Big}\)
such that
\[
 F=\sum_tG_t,\qquad
            \mathcal T_{(G_t)_i}(\lambda^*)\in\sigma_t
                         \quad\text{for every }i,\lambda^*.
\]

Apply ADP5 to the finite-dimensional span of all
\(\mathcal T_{F_i}(\lambda^*)\), as \(i\) ranges over its finite
index set and \(\lambda^*\) over a basis of \(D^*\).  Obtain the
global equivariant projections \(P_t\).  Since \(P_t\) commutes with
the infinity action, \(P_t\mathcal T_{F_i}\) satisfies the exact
equivariance required in ADP8.  Let \((G_t)_i\) be its inverse
scalarization.  These functions have rational coefficient equivariance
and have all scalarizations in \(\sigma_t\).

It remains to check the base away coefficient equation defining
\(H.\mathrm{Big}\).  For \(u\) in its original base subgroup, choose
\(k_u\in G_f\) with \(p(k_u)=(u,1)\), using ADP6.  ADP7 gives
\[
 \mathcal T_{F_i(\,\cdot\,(u,1))}(\lambda^*)
      =R_{k_u}\mathcal T_{F_i}(\lambda^*).
\]
The original matrix equation is a finite linear combination of these
identities with coefficients \((\tau_C(u))_{ij}\).  Apply \(P_t\),
commute it past \(R_{k_u}\), and use its complex linearity.  ADP7
injectivity recovers the identical matrix equation for \(G_t\).
Thus \(G_t\in H.\mathrm{Big}\).  Finally
\(\sum_tP_t\mathcal T_{F_i}=\mathcal T_{F_i}\); inverse
scalarization gives \(\sum_tG_t=F\).  This checks the actual
coefficient action, including all matrix mixing, rather than replacing
it by a scalar invariance assumption.
% END ADP11

% BEGIN ADP12
\paragraph{ADP12. Full original \texttt{decomp\allowbreak U} and its exact remaining boundary.}
\label{claim:ADP12}
For every datum and vector in the original \texttt{decomp\allowbreak U} statement,
\[
 F\in\sum_{\sigma:\,\mathrm{IsAutRepU}(\sigma)}
       Y_{\sigma},\qquad
 Y_{\sigma}=H.\mathrm{Big}\cap
       \operatorname{span}_{\mathbf C}
       \{G:\mathcal T_{G_i}(\lambda^*)\in\sigma
                                   \text{ for all }i,\lambda^*\}.
\]
This is the literal source \(YOf\), with the specified \(E,\xi\).

ADP10 proves automorphy for all scalarizations from the exact hypotheses
of \texttt{decomp\allowbreak U}.  ADP11 gives a finite sum of vectors in the
displayed spaces.  Each vector lies even in the generating set before
taking its span.  This proves the endpoint with the same functions,
coefficient representation, embeddings, and rational group.

The theorem makes no arithmetic assertion about \(\sigma\).  To turn
\(Y_\sigma\ne0\) into a \(\mathrm{PartAltAlt}\) one must still
construct its Galois representation, spherical polynomial readings,
local ordinary comparison, auxiliary local correspondence, and the
original point-property record.  The source's
\texttt{part\allowbreak U\_\allowbreak alt}/\texttt{part\allowbreak X\_\allowbreak alt} is the separate producer of
those data.  Once it is supplied, discard the zero \(Y_\sigma\)'s,
decorate each remaining one without changing its subspace, and apply
CFD12.  No assertion of arithmetic nonvanishing or completion of the
same-ideal Hida theorem follows solely from this decomposition.
% END ADP12


\section{Decomposing the actual classical and ordinary spaces}
\label{sec:classical-parts-v43}

This chapter proves the original \texttt{TCl\allowbreak Parts.\allowbreak cl\allowbreak Parts\_\allowbreak alt} and its
version for an arbitrary predicate on the parts.  The input is the original
\texttt{Automorphic\allowbreak Hecke\allowbreak Set}: vectors in the big automorphic space, with
some finite level and some Iwahori depth, are algebraic sums of the stated
parts.  The conclusion preserves the \emph{given} level, finite-part
character, depth, and ordinary generalized-eigenvalue condition.  In
particular, no passage to a closure or to a different automorphic space is
made.  The proof of the big automorphic decomposition, and the arithmetic
construction of the Galois data decorating each part, are distinct inputs.

All source correspondences use commit
\texttt{7cbf7a0c32fca1badf05931879a37c3f8a2b03b4}.
The notation is that of FNL, OST, and ASP.  Fix the original datum
\(H\), satisfying \(H.\mathrm{Ax}\), a complete Noetherian local
coefficient ring \(A\), its original weight character \(\theta\),
an arithmetic point \(x:A\to\overline{\mathbf Q}_\ell\),
the fixed field isomorphism \(\iota\), an admitted weight \(\lambda\),
and the original algebraic representation \(\xi\) on \(\mathbf C^d\).
Write
\[
 \mathcal M=\{F:G^a\times G^\ell\longrightarrow
                  \mathbf C^{m}\otimes\mathbf C^d\}.
\]
This is the same ambient function space as the source's iterated function
type.  Let \(\rho_a\) denote the away action with coefficient factor
\(\tau_C\), and \(\rho_p\) the product local right action.  For an
auxiliary level \(K=(K_q)\), depth \(e\), and the original finite-part
character \(\phi=\mathrm{finPart}(\theta,x,\lambda)\), put
\(C_{K,e}=H.\mathrm{ClAt}(\theta,\iota,x,\lambda,\xi,K,e)\).
The ordinary space is denoted \(O_{K,e}\).

An original \emph{part} \(Y\) is the actual subspace in a
\texttt{Part\allowbreak Alt\allowbreak Alt} record.  Among its fields are \(Y\subset H.\mathrm{Big}\),
nonvanishing, and local isotypic presentations at every auxiliary place
and every place above \(\ell\).  Its other fields retain the spherical
eigenvalues, Galois representation, and local comparison conditions.  This
chapter never modifies any of these fields.  A family may be selected by
\(P(\lambda,r_Y)\), as in \texttt{cl\allowbreak Parts\_\allowbreak alt}, or by a predicate on the
whole record, as needed in \texttt{Ord\allowbreak Part\allowbreak Pt}.

\subsection{Averaging while preserving all scalar equations}

% BEGIN CFD1
\paragraph{CFD1. The scalar profile of one vector.}
\label{claim:CFD1}
Let \(\rho:G\to\operatorname{GL}(M)\) be a representation over
\(\mathbf C\), \(K\le G\), and \(F\in M\).  Define
\[
 E_K(F)=\{y\in M: \rho(k)F=aF\Longrightarrow \rho(k)y=ay
                       \text{ for every }k\in K,\ a\in\mathbf C\}.
\]
This is a linear subspace.  Every endomorphism commuting with all
\(\rho(k)\), \(k\in K\), preserves it.  If \(F\ne0\), then
\[
 S=\{k\in K:\rho(k)F\in\mathbf C F\}
\]
is a subgroup and there is a unique character
\(\chi:S\to\mathbf C^\times\) with \(\rho(s)F=\chi(s)F\).
Moreover
\[
 E_K(F)=\{y:\rho(s)y=\chi(s)y\ (s\in S)\}.
\]

For the first assertions, apply linearity to each displayed scalar
equation, and move a commuting endomorphism across \(\rho(k)\).
When \(F\ne0\), the scalar in such an equation is unique and nonzero,
since \(\rho(k)\) is invertible.  Applying \(\rho(s)\rho(t)\) to
\(F\) proves multiplicativity of the scalar; applying \(\rho(s^{-1})\)
proves the inverse assertion.  These calculations prove both the subgroup
claim and the description of \(E_K(F)\).  The nonzero hypothesis is used
only for this description; the definition itself also covers \(F=0\).
% END CFD1

% BEGIN CFD2
\paragraph{CFD2. Finite-index averaging without a normality assumption.}
\label{claim:CFD2}
Let \(\pi:S\to\operatorname{GL}(M)\), and let \(U\le S\) have
finite index.  Choose representatives \(r_q\) of the left cosets
\(S/U\), and define the endomorphism
\[
 P=\frac1{[S:U]}\sum_{q\in S/U}\pi(r_q).
\]
If \(y\) is \(U\)-fixed, then \(Py\) is \(S\)-fixed and does not
depend on the chosen representatives.  If \(y\) is already \(S\)-fixed,
then \(Py=y\).  Finally \(P\) preserves every \(S\)-stable subspace.

Replacing \(r_q\) by \(r_qu\), \(u\in U\), does not change its
action on a \(U\)-fixed vector.  For \(s\in S\), the representatives
\(sr_q\) enumerate the same left cosets, so \(\pi(s)Py=Py\).
The other assertions follow term by term.  No assertion that \(P\)
projects all of \(M\) is needed: invariance of \(Py\) is asserted on
the \(U\)-fixed domain.  In particular the argument does not incorrectly
treat a nonnormal \(U\) as a quotient group.
% END CFD2

% BEGIN CFD3
\paragraph{CFD3. Refinement of a finite sum to the scalar profile.}
\label{claim:CFD3}
Let \((Y_i)_{i\in I}\) be subspaces of \(M\), stable under \(K\).
Suppose every \(y\in Y_i\) is fixed by a subgroup \(U_y\le G\) for
which \([K:K\cap U_y]<\infty\).  If
\(F\in\sum_iY_i\), then
\[
                         F\in\sum_i(Y_i\cap E_K(F)).
\]
The sum here is algebraic, so each membership involves only finitely many
summands.

There is nothing to prove for \(F=0\).  Otherwise write
\(F=\sum_{i\in I_0}y_i\), with \(I_0\) finite.  Let \(S,\chi\)
be as in CFD1, and put \(U=S\cap\bigcap_{i\in I_0}U_{y_i}\).
Its index in \(S\) is finite: restriction embeds each coset set into the
corresponding finite coset set for \(K\), and intersections inject into
products of these coset sets.  Every \(u\in U\) fixes all \(y_i\), hence
fixes their nonzero sum, so \(\chi(u)=1\).  Apply CFD2 to the twisted
representation \(\pi(s)=\chi(s)^{-1}\rho(s)\) and this \(U\).
Each \(y_i\) is \(\pi(U)\)-fixed; therefore \(Py_i\) belongs to
\(Y_i\cap E_K(F)\).  Since \(F\) is \(\pi(S)\)-fixed,
\(F=PF=\sum_iPy_i\).  This proves the assertion without assuming that
the original sum of the \(Y_i\)'s is direct.
% END CFD3

% BEGIN CFD4
\paragraph{CFD4. Simultaneous refinement for commuting local factors.}
\label{claim:CFD4}
Suppose a finite family of groups \(G_j\) acts on \(M\), the actions
of distinct factors commute, and every \(Y_i\) satisfies the hypotheses
of CFD3 for each \(K_j\le G_j\).  Then
\[
 F\in\sum_iY_i\quad\Longrightarrow\quad
 F\in\sum_i\left(Y_i\cap\bigcap_j E_{K_j}(F)\right).
\]

Apply CFD3 to the first factor.  By CFD1 and commutation, every later
factor preserves the scalar-profile subspace already imposed.  Its
intersection with a \(Y_i\) is still stable, and its vectors still have
the required finite-index stabilizers, since they are vectors of \(Y_i\).
Induction on the finite set of factors proves the displayed formula.
An empty family gives the initial algebraic sum itself.
% END CFD4

% BEGIN CFD5
\paragraph{CFD5. Local isotypic sums supply the required stabilizers.}
\label{claim:CFD5}
Let \(G\) be a topological group, \(K\le G\) compact, and let
\(Y\) be an algebraic sum of images of equivariant maps from one smooth
representation \(\sigma\).  Then \(Y\) is \(G\)-stable and every
\(y\in Y\) is fixed by a subgroup \(U\) with
\([K:K\cap U]<\infty\).

Write \(y\) as a finite sum of images of vectors of \(\sigma\).
Smoothness gives an open stabilizer for each vector; their finite
intersection \(U\) fixes \(y\).  The open cosets of \(K\cap U\)
cover compact \(K\), so finitely many suffice.  Equivariance proves
stability term by term.  This is exactly the content of the original
\texttt{Iso\allowbreak Sum} fields of a part at each auxiliary place and each place
above \(\ell\).  Irreducibility of \(\sigma\) is not required for
this particular implication, although it remains part of the original
part record.
% END CFD5

\subsection{The original level and ordinary operators}

% BEGIN CFD6
\paragraph{CFD6. Commutation on the original ambient function space.}
\label{claim:CFD6}
The actions at different auxiliary places commute, the actions at
different places above \(\ell\) commute, and every auxiliary action
commutes with every action above \(\ell\).  Every original part is
stable under the whole product \(G^\ell\).

For the first assertion, the \(Q\)-\(Q\) field of \(H.\mathrm{Ax}\)
identifies the original sections as commuting elements in \(U_b\).
Apply the coefficient action \(\rho_a\) to that identity.  For the
second assertion, the local elements lie in distinct coordinates of
\(G^\ell\), hence commute.  For the mixed assertion, on each coordinate
of a function the two expressions are
\[
  \tau_C(u)F\bigl(g(u,1)(1,b)\bigr),\qquad
  \tau_C(u)F\bigl(g(1,b)(u,1)\bigr),
\]
which agree.  These are the original actions; no abstract conjugacy is
used.  FNL12 verifies the relevant section identities for the actual
unitary datum.  Local stability of a part follows from CFD5, and finite
coordinate factorization gives stability under all of \(G^\ell\).
% END CFD6

% BEGIN CFD7
\paragraph{CFD7. Exact classical equations imply the required Iwahori fixedness.}
\label{claim:CFD7}
The original classical space is exactly the space of \(F\in H.\mathrm{Big}\)
satisfying
\begin{align*}
 \rho_a(u)F&=F &&(u\in H.\mathrm{lev}(K)),\\
 \rho_p(b)F&=F &&(b\in N\text{ or }b\in N^-(e)),\\
 \rho_p(e_V(v,i,y))F&=\iota(\phi(v,i,y))^{-1}F
                  &&(y\in\mathcal O_{L_{w(v)}}^\times).
\end{align*}
Suppose the specialization datum has depth \(c_0\), so that
\(\phi(v,i,y)=1\) on the original \(U_c(c_0+1)\).  For every
\(e\ge c_0\), a vector of \(C_{K,e}\) is fixed by \(I_1(e)\).

The first assertion is expansion of the original definition, as already
identified in ASP5.  The subgroup \(I_1(e)\) is generated by
\(N,T_c(e),N^-(e)\).  The first and last are fixed by the displayed
equations.  Each torus coordinate of \(T_c(e)\) is a product of the
original elementary diagonal elements \(e_V(v,i,y_i)\), with
\(y_i\in U_c(e+1)\subset U_c(c_0+1)\).  The finite-part character
is therefore one on each factor.  This fixes the middle generator too.
The factorization and the shift by one are the exact GL and product
Iwahori identities in OST26--OST27 and ASP4, not a replacement convention
for depth.
% END CFD7

% BEGIN CFD8
\paragraph{CFD8. The actual distinguished operator preserves all required spaces.}
\label{claim:CFD8}
Let \(T=H.\mathrm{uzC}(\iota,\lambda,d)\), with its original
\(\mathrm{lowVal}\) normalization.  If \(e\ge c_0\), then \(T\)
preserves \(C_{K,e}\).  It also preserves every original part \(Y\).

On the product action, \(T\) is its stated nonzero scalar times the
finite \(N\)-coset sum defining \(U_z\).  CFD6 shows that each part is
stable under every term.  The big space and its away level condition are
also preserved term by term: finite right translation preserves rational
equivariance, and the away action commutes with it.

For the remaining classical conditions, CFD7 gives \(I_1(e)\)-fixedness.
The original Iwahori level calculation, OST5--OST6, shows that \(U_z\)
preserves this fixedness at depth \(e\).  The same double-coset calculation
commutes \(U_z\) with each elementary integral torus element on
\(N\)-fixed vectors.  Thus its prescribed scalar equation is preserved.
Together these give every equation in CFD7.  Multiplication by the
original normalization scalar does not change stability.  This also
agrees with the positive-operator identification in ASP6.
% END CFD8

% BEGIN CFD9
\paragraph{CFD9. A decomposition in big parts preserves the given classical space.}
\label{claim:CFD9}
Let \((Y_i)\) be any family of original parts and let
\(F\in C_{K,e}\cap\sum_iY_i\).  Then
\[
                         F\in\sum_i(Y_i\cap C_{K,e}).
\]
This implication does not require \(e\ge c_0\).

Apply CFD4 first to the auxiliary groups with their compact subgroups
\(K_q\), and then to the factors above \(\ell\) with
\(\mathrm{GL}_n(\mathcal O_{L_{w(v)}})\).  CFD5 gives finite-index
stabilizers and CFD6 gives all commutations.  We obtain a finite sum
\(F=\sum_iF_i\), where \(F_i\in Y_i\) and every \(F_i\) satisfies
every scalar equation of \(F\) under those compact subgroups.

Since \(Y_i\subset H.\mathrm{Big}\), the base away subgroup is already
fixed on \(F_i\).  The extra \(K_q\)-conditions are inherited from
the scalar profile, with scalar one.  They generate
\(H.\mathrm{lev}(K)\) together with that base subgroup.  The original
\(N\), \(N^-(e)\), and integral torus all lie in the above maximal
compact product, by the GL Iwahori identities of OST26.  Their fixedness
or prescribed finite-part scalar equations are inherited from \(F\).
Thus all equations of CFD7 hold for each \(F_i\), proving the claim.
Neither the level \(K\) nor the depth \(e\) has changed.
% END CFD9

% BEGIN CFD10
\paragraph{CFD10. Polynomial projection preserves a decomposition into invariant pieces.}
\label{claim:CFD10}
Let \(W\) be a finite-dimensional complex vector space, \(T\in
\operatorname{End}(W)\), and let \(W=\sum_iW_i\), with every
\(W_i\) stable under \(T\).  Then, for each \(\mu\in\mathbf C\),
\[
 W[\mu]_{\mathrm{gen}}
       =\sum_i\bigl(W_i\cap W[\mu]_{\mathrm{gen}}\bigr).
\]

If \(\mu\) is not an eigenvalue, both sides vanish.  Otherwise factor
the minimal polynomial as \(\prod_\nu(X-\nu)^{a_\nu}\).  The elementary
Chinese remainder theorem supplies \(p_\mu\in\mathbf C[X]\) that is one
modulo \((X-\mu)^{a_\mu}\) and zero modulo every other factor.
Consequently \(p_\mu(T)\) is the identity on the \(\mu\)-generalized
eigenspace and zero on the others.  Given a finite expression
\(w=\sum_iw_i\), apply \(p_\mu(T)\): stability keeps each projected
term in \(W_i\).  This proves the nontrivial inclusion; the other is
immediate.  No diagonalizability and no directness of \(\sum_iW_i\)
has been assumed.
% END CFD10

% BEGIN CFD11
\paragraph{CFD11. Refinement to the exact ordinary generalized eigenspaces.}
\label{claim:CFD11}
Assume \(e\ge c_0\), the original specialization assertions, and
\(C_{K,e}=\sum_i(Y_i\cap C_{K,e})\).  Then
\[
                         O_{K,e}=\sum_i(Y_i\cap O_{K,e}).
\]

The original \(O_{K,e}\) is the sum, inside \(C_{K,e}\), of the
generalized eigenspaces of \(T\) with
\(\lvert\iota^{-1}(\mu)\rvert_\ell=1\).  This identification is ASP6;
it involves no continuity of \(\iota\).  The original specialization
finiteness clause, proved in ASP10, makes \(C_{K,e}\) finite-dimensional.
CFD8 proves stability of both \(C_{K,e}\) and \(Y_i\cap C_{K,e}\).
Apply CFD10 to every eigenvalue with the indicated norm and add the
equalities.  Each resulting summand lies in \(Y_i\cap O_{K,e}\).
The reverse inclusion is automatic.  The argument includes nilpotent
Jordan parts and uses the source's unit condition, without replacing it
by a condition on the complex absolute value.
% END CFD11

% BEGIN CFD12
\paragraph{CFD12. Full original classical-parts endpoint and its predicate refinement.}
\label{claim:CFD12}
Under every hypothesis of the original \texttt{cl\allowbreak Parts\_\allowbreak alt}, for every
original \(K\) and \(e\ge c_0\),
\begin{align*}
 C_{K,e}&=\sum_{Y:\,P(\lambda,r_Y)}(Y\cap C_{K,e}),\\
 O_{K,e}&=\sum_{Y:\,P(\lambda,r_Y)}(Y\cap O_{K,e}).
\end{align*}
The same conclusion holds with a predicate on the whole part record,
provided the corresponding big-space decomposition is supplied.

For \(F\in C_{K,e}\), membership in \(H.\mathrm{Big}\) and the
given away level are among its defining equations.  CFD7 gives the
required Iwahori-fixedness.  The algebraic-representation assertion of
the original specialization datum supplies the exact algebraic-weight
hypothesis.  Apply \(H.\mathrm{AutomorphicHeckeSet}\) at this
\(\lambda,\xi,F\), obtaining membership in the stated sum of parts.
CFD9 preserves the original level and depth in that sum, proving the
first equality.  CFD11 proves the second.  The proof used the selecting
predicate only to specify the family being summed, so it proves the
predicate refinement as well.

This closes the original generic structural theorem, including every
level and the original map and depth assertions.  For the actual unitary
construction, FNL16 supplies \(H.\mathrm{Ax}\), and ASP13 supplies
the specialization datum; the automorphic part producer remains a
separate arithmetic input until it is explicitly constructed.  In
particular no nonvanishing of a selected residual Hecke block is inferred
from either equality when the block could be zero.
% END CFD12


\section{Arithmetic fibers, scalar Hecke actions, and separation}
\label{sec:arithmetic-fibers-v43}

\status{This chapter continues the completed ordinary, control,
weight-representation, and specialization chapters.  It constructs
arithmetic points of the original weight ring, proves their density,
and develops the fiber and scalar-action arguments needed to separate
the original ordinary Hecke algebra.  The original automorphic-part
hypothesis is kept explicit wherever it is used.  Source predicates
are unpacked into the mathematical statements actually needed below;
no formal theorem is accepted in place of a written argument.}

Write $\overline{\mathbb Q}_\ell$ for the original algebraic closure,
with its original $\ell$-adic norm.  It is not replaced by its
completion.  A coefficient ring $O=c.O$ is the original finite
extension of $\mathbb Z_\ell$ inside this algebraic closure.  The
specified isomorphism $\iota:\overline{\mathbb Q}_\ell\simeq\mathbb C$
is always retained as an algebraic isomorphism.  No continuity of
$\iota$ is used.  All assertions involving an $\mathrm{HSet}$ retain
its complete original $\mathrm{Ax}$ hypothesis; in particular its
odd-prime condition includes $\ell=3$.  The weight-density argument
itself works for every prime, exactly as its source statement does.

\subsection{Dense arithmetic points in the actual inertial weight algebra}

% BEGIN AFS1
\paragraph{AFS1. Integral power series are separated by infinite Cartesian sets.}
\label{claim:AFS1}
Let $O$ be a coefficient DVR and let $S$ be a finite set.  For each
$s\in S$, let $D_s$ be an infinite subset of
$\{z\in\overline{\mathbb Q}_\ell:|z|_\ell<1\}$.  If
$f\in O[[X_s:s\in S]]$ evaluates to zero at every point of
$\prod_sD_s$, then $f=0$.  Each evaluation here is an actual
$O$-algebra map to $\overline{\mathbb Q}_\ell$.
\begin{proof}
For any finite tuple of arguments choose a finite extension $E$ of
$\mathbb Q_\ell$ containing them and the fraction field of $O$.
All arguments lie in $\mathfrak m_E$.  A series with integral
coefficients converges in $E$ at this tuple: the terms of total
degree at least $r$ have norm at most $\rho^r$, where
$\rho<1$ is the largest argument norm.  The same bound justifies
addition, multiplication, and regrouping.  Thus evaluation is the
claimed algebra map, with values in the algebraic closure itself.

We prove a stronger one-variable assertion.  If
$g(T)=\sum_{k\geq0}a_kT^k\in O_E[[T]]$ is nonzero, let
$C=\max_k|a_k|>0$ and let $N$ be the first index attaining this
maximum.  The maximum exists because the nonzero norms in a DVR
form a discrete sequence.  Then $g$ has at most $N$ distinct zeros
in the open unit disk of $\overline{\mathbb Q}_\ell$.
If $\alpha$ is one such zero, work in the finite extension
$E(\alpha)$ and put
\[
 b_j=\sum_{k>j}a_k\alpha^{k-j-1},\qquad
 h(T)=\sum_{j\geq0}b_jT^j.
\]
The sums converge, $|b_j|\leq C$, and the coefficient identities
$a_0=-\alpha b_0$, $a_j=b_{j-1}-\alpha b_j$ give
$g=(T-\alpha)h$.  If $N=0$, the term $a_0$ strictly dominates
every other term at $\alpha$, which contradicts $g(\alpha)=0$.
For $N>0$ the first coefficient of $h$ of norm $C$ is
$b_{N-1}$: its first summand is $a_N$ and all subsequent summands
have smaller norm; for $j<N-1$, the finitely many earlier terms
have norm less than $C$ and the remaining tail has norm at most
$C|\alpha|^{N-j-1}<C$.  Division therefore reduces the first
maximum index by one.  Every other zero remains a zero of $h$,
since it differs from $\alpha$.  Repeating over the finite
compositum of any proposed finite list of roots proves the bound.
This argument does not assume that an infinite root set has an
accumulation point inside the open disk.

Induct on $|S|$.  Fix all arguments but one.  Evaluation in those
arguments produces a one-variable series with coefficients in the
integer ring of one finite extension.  Its infinitely many zeros
force all its coefficients to vanish by the preceding bound.
Each coefficient is the evaluation of a series in the remaining
variables.  The induction hypothesis kills all those series.
For $S=\varnothing$ the sole evaluation is the inclusion of $O$,
so the assertion also holds.
\end{proof}
% END AFS1

% BEGIN AFS2
\paragraph{AFS2. Infinite disjoint torsion boxes with the original index condition.}
\label{claim:AFS2}
Let $n\geq1$, let $S$ be finite, and let
$\mathrm{idx}:S\to\{0,\ldots,n-1\}$.  Suppose for each $s$ that
$u_s\in\overline{\mathbb Q}_\ell$ has a positive power
satisfying $|u_s^{M_s}-1|<1$.  There are infinite sets $R_s$ of
roots of unity such that
\[
 |u_s\eta-1|<1\quad(\eta\in R_s),\qquad
 R_s\cap R_t=\varnothing\quad\hbox{if }\mathrm{idx}(s)\ne\mathrm{idx}(t).
\]
Consequently $f\in O[[X_s]]$ is zero if every algebra map $y$
with
\[
 1+y(X_s)=u_s\eta_s,
 \quad\eta_s\text{ of finite order},\qquad
 \eta_s\ne\eta_t\text{ whenever }\mathrm{idx}(s)\ne\mathrm{idx}(t)
\]
kills $f$.
\begin{proof}
The hypothesis implies $|u_s|=1$.  Put all the finitely many $u_s$
in one finite extension $E$.  If its residue field has $q$ elements,
$M=q-1$ is prime to $\ell$ and $|u_s^M-1|<1$ for every $s$.
Factoring $T^M-1$ in the algebraic closure gives a root $\xi_s$
with $|u_s-\xi_s|<1$: otherwise the product of all the distances
would be at least one.  Its norm is one, so
$\omega_s=\xi_s^{-1}$ satisfies
$\omega_s^M=1$ and $|u_s\omega_s-1|<1$.

For each $k\geq1$ choose a primitive $\ell^k$th root $\zeta_k$.
It satisfies $|\zeta_k-1|<1$: after reduction in the integer ring
of its finite extension, $T^{\ell^k}-1=(T-1)^{\ell^k}$, so its
residue is one.  Set
\[
 R_s=\{\omega_s\zeta_{\mathrm{idx}(s)+1+nt}:t\geq0\}.
\]
The $\ell$-primary part of the order of
$\omega_s\zeta_k$ is exactly $\ell^k$, because the order of
$\omega_s$ divides $M$ and $(M,\ell)=1$.  Equality of two such
elements forces equality of their $k$'s.  The displayed levels
are distinct for different $t$, and equality between levels for
$s,t$ forces their indices modulo $n$ to coincide.  This proves
infinitude and the required disjointness.  The norm inequality
follows from
$ab-1=a(b-1)+(a-1)$ for two units congruent to one.
The sets $D_s=\{u_s\eta-1:\eta\in R_s\}$ are infinite subsets
of the open unit disk.  Every evaluation on their Cartesian
product satisfies the prescribed conditions, so AFS1 applies.
\end{proof}
% END AFS2

% BEGIN AFS3
\paragraph{AFS3. Torsion on inertia coordinates gives the actual finite-part depth.}
\label{claim:AFS3}
Use the original local decomposition groups $G_v$, inertia groups
$I_v$, and stored Artin maps $a_v=\mathrm{artOf}(s.hw(v))$.
Choose the exact local abelian coordinates furnished by KX3,
\[
 a_v^{\rm ab}:G_v\longrightarrow
            \mathbb Z_\ell^{d_v}\times\mathbb Z_\ell,
 \qquad d_v=[L_{s.w(v)}:\mathbb Q_\ell]>0,
\]
and choose $\gamma_{v,j}\in I_v$ mapping to $(e_j,0)$.
If a continuous character
$\Phi:G_v\to\overline{\mathbb Q}_\ell^{\times}$ has finite-order
values on every $\gamma_{v,j}$, then $\Phi|_{I_v}$ has finite
image.  Some integer $e$ has the following property for all
$e'\geq e$: if $t$ is in the original Weil group and
$a_v(t)\in U_c(L_{s.w(v)},e')$, then $\Phi(t)=1$.
\begin{proof}
KX3 identifies the image of inertia in the topological
abelianization of the actual $G_v$ with
\[
  O_{L_{s.w(v)}}^{\times}\simeq
       T_v\times\mathbb Z_\ell^{d_v},
\]
where $T_v$ is finite, including both its prime-to-$\ell$ and
$\ell$-primary torsion.  Under this identification the selected
map $a_v^{\rm ab}$ is projection to the free coordinates on
inertia.  A continuous character into a Hausdorff abelian group
kills the closure of the commutator subgroup.  It thus factors
through this same abelianization.  Choose a common positive
multiple $N$ of the orders of its finitely many torsion values
and of the values on the $\gamma_{v,j}$.  The character $\Phi^N$
is trivial on $T_v$ and on all the coordinate generators.
Their integer span is dense in $\mathbb Z_\ell^{d_v}$, so
continuity makes $\Phi^N$ trivial on all inertia.  Its image
therefore lies in the finite group of $N$th roots of unity.
Its kernel is open, since this finite subset of a Hausdorff
field is discrete.

CIP4 identifies the stored Artin map with the same original
local reciprocity quotient, with
$\operatorname{ord}(a_v(t))=-\deg(t)$.  The character on inertia
therefore descends continuously to a finite character of the
local unit group.  An open subgroup of that group contains
$U_c(F_v,e)$ for some $e$, and then contains every
$U_c(F_v,e')$ with $e'\geq e$.  The valuation identity shows
that $a_v(t)$ is a unit precisely when $t$ belongs to inertia.
The descended character evaluated at $a_v(t)$ is $\Phi(t)$;
this proves the last assertion.  This uses the full local
abelianization already proved in KX3, rather than assuming
the private formal lemma about an arbitrary abstract
$\mathrm{LocAb}$ record.
\end{proof}
% END AFS3

% BEGIN AFS4
\paragraph{AFS4. The original algebraic character and finite-part identity.}
\label{claim:AFS4}
Fix all the literal data of
\texttt{TDense\allowbreak U.\allowbreak is\allowbreak Domain\_\allowbreak and\_\allowbreak dense\allowbreak Wt\_\allowbreak of\_\allowbreak is\allowbreak Theta}, including
the supplied $\mathrm{LamObj}\ r$, its characters $\kappa_{v,i}$,
the character $\theta$ satisfying $\mathrm{IsTheta}$, and any
dominant weight $\lambda$.  For every $v,i$ there is a
continuous character $\delta_{v,i}:G_v\to
\overline{\mathbb Q}_\ell^{\times}$ such that, for inertia
elements $t$ of the original Weil group,
\[
 \delta_{v,i}(t)=
  \operatorname{algPow}(s.w(v),\operatorname{lowW}(\lambda,i),a_v(t)).
\]
Choose any full witnessing character $\psi_{v,i}$ extending
$\kappa_{v,i}$.  For each original algebra point
$x:\Lambda\to\overline{\mathbb Q}_\ell$ put
\[
 \Phi_{x,v,i}=(x\circ\psi_{v,i})\,
                     \varepsilon_\ell^{\,i}\delta_{v,i}.
\]
These are continuous, and if $a_v(t)=y$ is a local unit then
\begin{equation}\label{eq:afs-finpart}
 \operatorname{finPart}(\theta,x,\lambda,v,i,y)
                       =\Phi_{x,v,i}(t)^{-1}.
\end{equation}
\begin{proof}
The finite set of continuous embeddings
$j:F_v\to\overline{\mathbb Q}_\ell$ defines on local units
the continuous character
\[
 \alpha(y)=\prod_jj(y)^{
          \operatorname{lowW}(\lambda,i)(j|_L)}.
\]
All values lie in the unit group of one finite extension:
each embedding is an isometric $\mathbb Q_\ell$-embedding
(WRP1), and finitely many such finite image fields have finite
compositum.  Extend $\alpha$ to $F_v^{\times}$ by making it
one on a chosen uniformizer.  It extends continuously to
$\widehat{F_v^{\times}}=\widehat{\mathbb Z}\times O_{F_v}^{\times}$.
Composing with the actual reciprocity identification of KX3
and CIP4 gives $\delta_{v,i}$.  Only its restriction to inertia
is required, so its unramified choice changes none of the
following identities.

ASP2 gives, for the given $x$, a number $0<\eta<1$ with
$|x(\mathfrak m_\Lambda^r)|\leq\eta^r$.  Thus composition
with $x$ carries any adically continuous character to a
continuous $\ell$-adic character: choose $r$ with
$\eta^r<\epsilon$ and use equality modulo
$\mathfrak m_\Lambda^r$ on a sufficiently small neighborhood.
The cyclotomic character and $\delta_{v,i}$ are continuous,
which proves the asserted continuity of $\Phi$.

The literal $\mathrm{IsTheta}$ identity, evaluated at $x$, is
\[
 x(\theta(e_{v,i}(y)))\,x(\kappa_{v,i}(t))\,
                     \varepsilon_\ell(t)^i=1.
\]
Here $t$ is in inertia because its Artin image is a unit,
and $\psi_{v,i}(t)=\kappa_{v,i}(t)$.  Divide this identity by
the original $\operatorname{algPow}$ factor.  The defining
formula for $\delta$ proves \eqref{eq:afs-finpart}, with the
inverse and the exponent $i$ exactly as in the source.
\end{proof}
% END AFS4

% BEGIN AFS5
\paragraph{AFS5. Every supplied original weight ring is a domain and has dense good points.}
\label{claim:AFS5}
Under precisely the quantified hypotheses of
\texttt{TDense\allowbreak U.\allowbreak is\allowbreak Domain\_\allowbreak and\_\allowbreak dense\allowbreak Wt\_\allowbreak of\_\allowbreak is\allowbreak Theta}, the given
$\Lambda.\Lambda$ is a domain, and for every $\lambda$ it
satisfies the original $H.\mathrm{DenseWt}(\theta,\lambda)$,
where $H=\mathrm{hset}(Q,\tau_0)$.  In expanded form,
\[
 \bigl[\forall x:\Lambda\to_O\overline{\mathbb Q}_\ell,
       \ H.\mathrm{Good}(\theta,x,\lambda)\Rightarrow x(f)=0\bigr]
                    \quad\Longrightarrow\quad f=0.
\]
\begin{proof}
KX4--KX5 give, for this supplied universal object, the unique
isomorphism
\[
 e:O[[X_{v,i,j}:v\mid\ell,\ 0\leq i<n,\ 1\leq j\leq d_v]]
                         \overset\sim\longrightarrow\Lambda
\]
and the actual finite coefficient characters $\psi^0_{v,i}$.
At the selected inertia generators its formula is
\begin{equation}\label{eq:afs-coordinate}
 \kappa_{v,i}(\gamma_{v,j})
       =\psi^0_{v,i}(\gamma_{v,j})\bigl(1+e(X_{v,i,j})\bigr).
\end{equation}
The map $e$ respects every $O$-algebra map in the original
universal property; it does not replace the prescribed $\Lambda$
by an unrelated model.  A power-series ring over a domain is a
domain: choose the least nonzero monomial in the well-ordered
graded lexicographic order; the least monomial in a product has
the product of those coefficients.  Hence $\Lambda$ is a domain.

For the fixed $\lambda$ take AFS4's $\delta$ and set
\[
 u_{v,i,j}=
 \bigl(\psi^0_{v,i}(\gamma_{v,j})
       \varepsilon_\ell(\gamma_{v,j})^i
       \delta_{v,i}(\gamma_{v,j})\bigr)^{-1}.
\]
Every factor is a unit in the integer ring of a finite extension,
so each $u_{v,i,j}$ has a positive power congruent to one.
AFS2 supplies a Cartesian family of evaluations $y$ such that
\[
 1+y(X_{v,i,j})=u_{v,i,j}\eta_{v,i,j},
\]
where all $\eta$ have finite order and their values are
different whenever their $i$ indices differ.  Put
$x=y\circ e^{-1}$.  Applying $x$ to
\eqref{eq:afs-coordinate} and canceling its unit factors gives
\begin{equation}\label{eq:afs-torsionvalues}
             \Phi_{x,v,i}(\gamma_{v,j})=\eta_{v,i,j}.
\end{equation}
AFS3 gives, for each $v,i$, a depth on which this character
is trivial on Artin preimages.  There are finitely many
$v,i$, so take one common upper bound.  Equation
\eqref{eq:afs-finpart} proves the exact $H.\mathrm{Arith}$
condition at this depth.

For $i\ne i'$ at a given $v$, choose $j=1$, which exists
because $d_v>0$ in KX3.  The original Artin image
$y=a_v(\gamma_{v,1})$ is a local unit.  By
\eqref{eq:afs-finpart} and \eqref{eq:afs-torsionvalues}, the
two finite parts at $y$ are the inverses of the distinct
$\eta_{v,i,1}$ and $\eta_{v,i',1}$.  They are distinct.
For the original $\mathrm{hset}(Q,\tau_0)$ every dominant
$\lambda$ is admitted.  This verifies all three clauses of
$H.\mathrm{Good}$, without changing the original characters.

If $f$ vanishes at every such good point, then $e^{-1}(f)$
vanishes on the entire torsion box just constructed.  AFS2
forces it to be zero, whence $f=0$.  All choices were made
after fixing $\lambda$ and before choosing the tested element
$f$; no point is selected by a coefficient of $f$.
\end{proof}
% END AFS5

% BEGIN AFS6
\paragraph{AFS6. The original increasing weight has a consistent dominant realization.}
\label{claim:AFS6}
For the original Hida setting $s$, every strictly increasing
weight $h:s.\mathrm{Wt}(n)$ has a dominant weight $\lambda$
with the literal relation $\mathrm{WtOf}(h,\lambda)$.
\begin{proof}
At a continuous embedding $\tau:F_v\to\overline{\mathbb Q}_\ell$
define
\[
 \lambda(\tau|_L)_j=h(v,n-1-j,\tau)-(n-1-j),
                      \qquad 0\leq j<n.
\]
Strict increase of the integer-valued $h$ implies
$h(v,b,\tau)-h(v,a,\tau)\geq b-a$ whenever $a\leq b$.
The displayed sequence is therefore dominant in the source's
order.  Its reversed coordinate gives
\[
 \operatorname{lowW}(\lambda,i)(\tau|_L)+i=h(v,i,\tau),
\]
which is exactly $\mathrm{WtOf}$.

This definition is consistent.  A continuous field embedding
of a local completion preserves its valuation, so its
restriction to $L$ determines the associated finite prime.
More explicitly, if distinct primes gave the same restriction,
choose integral $a,b$ with $a$ in the first prime, $b$ in the
second and $a+b=1$.  The common embedding would give both
$|a|<1$ and $|b|<1$, contradicting the ultrametric inequality
at $1$.  At the same prime, density of $L$ in its completion
and continuity show that the two embeddings coincide.  The
chosen primes $s.w(v)$ are themselves distinct because their
restrictions to the maximal real subfield are the original
$v$.  Thus each prescribed global embedding receives only
one value.  Set $\lambda=0$ at any remaining global embeddings.
The resulting dominant weight has all the required identities.
\end{proof}
% END AFS6

% BEGIN AFS7
\paragraph{AFS7. Local algebraicity and finite parts agree on the stored Artin map.}
\label{claim:AFS7}
Let $\kappa:G_v\to\overline{\mathbb Q}_\ell^{\times}$ be a
continuous character, let $w$ be an integer weight on the
continuous embeddings of $F_v$, and suppose a function $f$
on local units satisfies, for every inertia element $\sigma$,
\begin{equation}\label{eq:afs-isfin}
 \kappa(\sigma)f(a_v(\sigma))
                =\prod_\tau\tau(a_v(\sigma))^{-w(\tau)}.
\end{equation}
Then the original $\mathrm{IsLocAlg}(s.hw(v),\kappa,w)$ is
equivalent to $f$ being one on some $U_c(F_v,e)$.
For every witnessing finite-part character $\varphi$ in
$\mathrm{IsWtPart}$ one has
$\varphi(\sigma)=f(a_v(\sigma))$.
\begin{proof}
The last identity follows by cancellation of the nonzero
$\kappa(\sigma)$ from the defining $\mathrm{IsWtPart}$
identity and \eqref{eq:afs-isfin}.  The stored Artin map is
surjective from inertia to local units: its full Weil
surjectivity and $\operatorname{ord}\circ a_v=-\deg$ give
this immediately.  It is a quotient map on inertia, because
it is a continuous surjection from a compact group to the
Hausdorff compact local unit group.

If $\varphi$ has open kernel, the identity just proved shows
that it is constant on Artin fibers and descends to $f$.
Its kernel consequently gives an open subgroup of the local
units, which contains some $U_c(F_v,e)$.  Conversely, define
on inertia
\[
 \varphi(\sigma)=\kappa(\sigma)^{-1}
                       \prod_\tau\tau(a_v(\sigma))^{-w(\tau)}.
\]
It is a character, equals $f\circ a_v$, and its kernel
contains the open inverse image of $U_c(F_v,e)$.  It has
the required open kernel and the required product identity.
The original definition quantifies over every map satisfying
$\mathrm{IsLocalArtin}$.  CIP4's uniqueness identifies each
of those with the stored $a_v$, so the displayed identity
satisfies that full quantifier as well.
\end{proof}
% END AFS7

% BEGIN AFS8
\paragraph{AFS8. Exact arithmetic and good-point bridge for the original determinant ring.}
\label{claim:AFS8}
Retain all data and hypotheses of
\texttt{TWt\allowbreak Bridge.\allowbreak arith\_\allowbreak good\_\allowbreak iff\_\allowbreak hset\_\allowbreak of\_\allowbreak theta\_\allowbreak incr\_\allowbreak weight}.
For every increasing $h$ choose $\lambda$ as in AFS6.  For
every $\lambda$ satisfying $\mathrm{WtOf}(h,\lambda)$ and
every original point $x:B.\mathrm{Pt}$, put
$x_\Lambda=x\circ\Lambda.\mathrm{toPair}(B.u)$.  Then
\begin{align*}
 x.\mathrm{Arith}(h)&\ \Longleftrightarrow\
       H.\mathrm{Arith}(\theta,x_\Lambda,\lambda),\\
 x.\mathrm{Good}(h)&\ \Longleftrightarrow\
       H.\mathrm{Good}(\theta,x_\Lambda,\lambda).
\end{align*}
The maps, inertia characters, and finite-part distinctions are
the original ones.
\begin{proof}
The universal property of $\Lambda.\mathrm{toPair}$ identifies
the image of $\kappa_{v,i}$ with the restriction of the
original $x.\psi_{v,i}$ to inertia.  The literal
$\mathrm{IsTheta}$ identity therefore gives
\[
 x.\psi_{v,i}(\sigma)\,
 x_\Lambda(\theta(e_{v,i}(a_v(\sigma))))\,
                         \varepsilon_\ell(\sigma)^i=1.
\]
Under the same reciprocity normalization, the local cyclotomic
identity is
$\varepsilon_\ell(\sigma)=\prod_\tau\tau(a_v(\sigma))$.
This is the local cyclotomic reciprocity identity on the
stored Artin map proved in CIP5.  Dividing by the
algebraic factor and using
$h(v,i,\tau)=i+\operatorname{lowW}(\lambda,i)(\tau|_L)$
gives \eqref{eq:afs-isfin} for
$\kappa=x.\psi_{v,i}$ and
$f=H.\operatorname{finPart}(\theta,x_\Lambda,\lambda,v,i)$.
AFS7 identifies local algebraicity with triviality at some
local unit depth.  A common upper bound for the finitely
many depths is exactly the single depth in $H.\mathrm{Arith}$.
The increasing-weight requirement on the determinant side
is the assumed $h.\mathrm{Incr}$.

Every finite-part inertia character is, by AFS7, exactly
$f\circ a_v$.  Since $a_v:I_v\to O_{F_v}^{\times}$ is
surjective, two such characters are equal if and only if
their functions $f$ are equal.  Thus their injective indexing
by $i$ is equivalent to the source condition
$\forall i\ne j\ \exists y,\ f_i(y)\ne f_j(y)$.
All dominant weights are admitted for the literal $H$ here.
These observations prove both equivalences, in both
directions, with the original universal-point map.
\end{proof}
% END AFS8

\subsection{Finite fibers and local scalar actions}

% BEGIN AFS9
\paragraph{AFS9. Every point of a faithful finite algebra occurs in its module fiber.}
\label{claim:AFS9}
Let $A$ be a commutative ring, let $N$ be a finite $A$-module,
and let $T\subset\operatorname{End}_A(N)$ be a commutative
$A$-subalgebra.  For a field $K$, an algebra map $x:A\to K$,
and a ring map $y:T\to K$ extending $x$, the finite-dimensional
space $V=K\otimes_A N$ contains a nonzero vector $v$ with
\[
             t v=y(t)v\qquad(t\in T).
\]
More generally, if $V=\sum_jV_j$ with every $V_j$ stable
under $T$, such a vector can be chosen in one of the $V_j$.
No flatness of $A\to K$ is assumed.
\begin{proof}
Put $\mathfrak p=\ker y$.  The module $N$ is finite over $T$
because an $A$-generating set also generates it over $T$.
It is faithful over $T$ by the specified inclusion into its
endomorphisms.  Thus $N_{\mathfrak p}\ne0$: if the localization
vanished, a product of finitely many elements outside
$\mathfrak p$ would annihilate all the chosen generators,
contradicting faithfulness.  Nakayama's lemma over the local
ring $T_{\mathfrak p}$ gives
$N_{\mathfrak p}/\mathfrak pN_{\mathfrak p}\ne0$.
The map $T/\mathfrak p\to K$ extends to its fraction field,
so extension of this nonzero vector space gives
\[
       K\otimes_{T,y}N\ne0.
\]
This is a quotient of $V$ on which $T$ acts by $y$.

Here is the finite algebra argument that turns this quotient
into an eigenvector, also keeping a stable summand.  Let
$D\subset\operatorname{End}_K(V)$ be the $K$-algebra generated
by the $T$-action.  It is commutative and finite-dimensional.
The nonzero quotient just constructed makes its scalar rule
a well-defined character $\chi:D\to K$.  A finite-dimensional
commutative algebra is a product of Artinian local algebras.
For clarity, the descending powers of its radical stabilize;
Nakayama after localization at each of its finitely many
maximal ideals makes the stable power zero.  The Chinese
remainder theorem for sufficiently large powers of the
maximal ideals then gives this product and its idempotents.
Let $e_\chi$ be the idempotent of the factor with character
$\chi$.  The space $e_\chi V$ is nonzero because it maps
onto the nonzero quotient.  Its local radical $\mathfrak r$
is nilpotent.  In a nonzero last power
$\mathfrak r^q(e_\chi V)$, every nonzero vector is killed by
$\mathfrak r$ and therefore has the required eigencharacter.

If $V=\sum_jV_j$, the scalar and $T$ stability of each $V_j$
implies $D$ stability.  Hence
$e_\chi V=\sum_je_\chi V_j$, so one $e_\chi V_j$ is nonzero.
Apply the same last-power argument within that space.
This proves both assertions, including an infinite indexing
set $j$, since membership in a sum uses finitely many terms.
The argument also shows the converse elementary fact used
below: an eigencharacter in a finite-dimensional commuting
representation occurs in a nonzero one-dimensional quotient
of its corresponding local block.
\end{proof}
% END AFS9

% BEGIN AFS10
\paragraph{AFS10. Nonzero ordinary blocks have nonzero fibers at every given point.}
\label{claim:AFS10}
For the original $H,A,\theta,a,K$ satisfying $H.\mathrm{Ax}$,
$\mathrm{IsCO}$ and $H.\mathrm{IsWt}$, set
\[
 N_0=H.X.\mathrm{Ord}(H.\tau_A,\theta,H.\mathrm{lev}(K)),
 \qquad N=H.MK(\theta,a,K),\qquad T=H.TK(\theta,a,K).
\]
If $N\ne0$, every $O$-algebra point $x:A\to
\overline{\mathbb Q}_\ell$ has a nonzero fiber
$\overline{\mathbb Q}_\ell\otimes_{A,x}N$ and extends to at
least one $O$-algebra point of $T$.  Every specified such
extension occurs as a nonzero simultaneous eigenvector in
that fiber.  At an arithmetic point, ASP's same specialization
map identifies its complex scalar extension injectively with
$\operatorname{span}_{\mathbb C}\operatorname{sp}(N)$.
\begin{proof}
OHA6--OHA7 identify $N$ with the image of an idempotent
$e\in\operatorname{End}_A(N_0)$ belonging to the original
Hecke algebra.  Both $N_0$ and $N$ are finite free over the
given local ring $A$.  Nonzero $N$ has a positive rank, so
every field fiber of it is nonzero, even for a nonflat point
map.  The image of $T$ in the endomorphisms of this fiber
generates a finite-dimensional commutative algebra over the
algebraically closed field $\overline{\mathbb Q}_\ell$.
A maximal ideal has residue field equal to that field;
the corresponding local block has a nonzero eigenvector
by the last-power argument in AFS9.  Its character restricts
to a point of $T$ extending $x$.  For any specified extension,
AFS9 applies to the actual faithful $T$-module $N$.

At an arithmetic point ASP11 says that the images under
$\operatorname{sp}$ of every $A$-basis of $N_0$ are
$\mathbb C$-linearly independent.  Therefore its tensor map
\[
 \mathbb C\otimes_{A,\iota x}N_0
     \longrightarrow\operatorname{span}_{\mathbb C}
                                      \operatorname{sp}(N_0)
\]
is an isomorphism.  The inclusion $N\hookrightarrow N_0$
is split by $e$, so it remains injective after arbitrary
scalar extension.  Restricting the preceding tensor map
gives precisely the asserted identification on $N$.
The assertion is about a nonzero specified block; it does
not assert that an arbitrary residual label produces a
nonzero block.
\end{proof}
% END AFS10

% BEGIN AFS11
\paragraph{AFS11. Stable positive-operator images inject into the original Jacquet module.}
\label{claim:AFS11}
Let $F$ be any nonarchimedean local field, let
$G=\mathrm{GL}_n(F)$ with its original $\mathrm{glIwa}$ data,
and let $\pi$ be a smooth complex representation.  Write
$N$ for the compact upper-unipotent subgroup, $N_F$ for the
full upper-unipotent group, $z=\operatorname{diag}
(\varpi^{n-1},\ldots,1)$, and $p:\pi\to J_N(\pi)$ for the
original unnormalized coinvariant map.  For $N$-fixed $v$,
\begin{equation}\label{eq:afs-jac-kernel}
 p(v)=0\quad\Longleftrightarrow\quad
                      U_z^rv=0\text{ for some }r\geq0.
\end{equation}
If $\pi$ is irreducible and the original
$\mathrm{LiteratureRankNHypothesis}$ is supplied, then for
every character $\varphi$ trivial on $T_c(c)$ the map $p$
is injective on
\[
 \operatorname{invPart}(\pi,c,\varphi)
       =\bigcap_{r\geq0}U_z^r\operatorname{eig}(\pi,c,\varphi).
\]
\begin{proof}
The integral double-coset formulas of OST5 give
$U_z^r=U_{z^r}$ on $N$-fixed vectors.  Since $p$ kills
$\pi(u)w-w$ for $u\in N_F$, the original sum over
$N/(aNa^{-1})$ gives
\begin{equation}\label{eq:afs-jac-sum}
 p(U_av)=[N:aNa^{-1}]\,p(\pi(a)v)\qquad(a\in P).
\end{equation}
These are the actual cosets and the integer index; no
normalizing measure is inserted.

If $p(v)=0$, write $v$ as a finite sum of relations
$\pi(u)w-w$ with $u\in N_F$.  Conjugation by $z^r$ carries
all these finitely many $u$ into $N$ for sufficiently large
$r$: their strictly upper entries acquire powers
$\varpi^{r(j-i)}$.  Thus $\pi(z^r)v$ belongs to the span of
$N$-relations.  Average over $N$ on the finite-dimensional
span needed for these relations.  This is a genuine finite
average: smoothness gives an open subgroup fixing the
finitely many vectors, and its normal core has finite index
in compact $N$.  The average kills every $N$-relation by
permutation of cosets.  On $\pi(z^r)v$ it equals the
nonzero scalar $[N:z^rNz^{-r}]^{-1}$ times $U_{z^r}v$,
since this vector is fixed by $z^rNz^{-r}$.  Hence
$U_z^rv=0$.  Conversely, if $U_z^rv=0$, equation
\eqref{eq:afs-jac-sum}, the invertibility of $\pi(z^r)$ on
coinvariants, and the nonzero integer index imply $p(v)=0$.

The eigenspace in the statement is fixed by $I_1(c)$:
factor an element as upper-unipotent, $T_c(c)$, and lower
depth-$c$ factors, and use its three defining eigenlaws.
The admissibility clause of the retained
$h_{\rm Lit}.\mathrm{loc.bz}$ therefore makes it finite-dimensional.
OST5--OST6 show that every $U_a$, $a\in P$, preserves it
and commutes there with $U_z$.  In a finite-dimensional
space the decreasing images of $U_z^r$ stabilize, and $U_z$
is an automorphism on their intersection.  An element of
that intersection killed by a power of $U_z$ is zero.
Equation \eqref{eq:afs-jac-kernel} proves the claimed
injectivity.  The same commutation proves that all $U_a$
preserve this intersection.
\end{proof}
% END AFS11

% BEGIN AFS12
\paragraph{AFS12. Distinct unit characters force the exact principal series and its scalars.}
\label{claim:AFS12}
Retain exactly the hypotheses of
\texttt{Exists\allowbreak Form.\allowbreak exists\_\allowbreak principal\_\allowbreak series\_\allowbreak of\_\allowbreak
distinct\_\allowbreak inv\_\allowbreak ne\_\allowbreak bot}:
$F$ is any nonarchimedean local field, $n\geq0$, $\pi$ is a
smooth irreducible complex representation of
$\mathrm{GL}_n(F)$, $\varphi$ is a character of the original
compact diagonal torus, $\varphi(e_i(y))=1$ for
$y\in U_c(F,c+1)$, its coordinate restrictions are pairwise
distinct, and $\operatorname{invPart}(\pi,c,\varphi)\ne0$.
Assume the original $h_{\rm Lit}$.
There are smooth characters $\omega_i:F^\times\to\mathbb C^\times$
with
\[
 \varphi(e_i(y))=\omega_i(y)^{-1},\qquad
 \pi\simeq\mathrm{PS}(\omega),\qquad
 \dim_{\mathbb C}\operatorname{invPart}(\pi,c,\varphi)=1,
\]
and for every $j\leq n$ and every vector in this same
$\operatorname{invPart}$,
\begin{equation}\label{eq:afs-pseig}
 U_{a_j}v=\mathrm{psEig}(\omega,j)v,
 \qquad
 \mathrm{psEig}(\omega,j)=\prod_{0\leq i<j}
                   \omega_i(\varpi)q^{-i}.
\end{equation}
Here $q=|k_F|$, $a_j=\operatorname{diag}
(\underbrace{\varpi,\ldots,\varpi}_{j},1,\ldots,1)$,
and the source convention is
$\mathrm{PS}(\omega)=\operatorname{Ind}^{\rm norm}_B
(\omega_i|\cdot|^{(n-1)/2})_i$.
\begin{proof}
We spell out the three local literature clauses actually used.
The supplied $h_{\rm Lit}.\mathrm{loc.bz}$ says that smooth
irreducibles are admissible and that Jacquet modules of smooth
finite-length representations have finite length.  Its
$\mathrm{loc.geo}$ clause gives the normalized Jacquet
filtration of an induced representation by the actual Bruhat
cells.  Its $\mathrm{loc.disj}$ clause says that induction of
two irreducibles with totally disjoint cuspidal supports is
irreducible, where the excluded twists have exponents
$-1,0,1$.  These are explicit retained premises of this
source endpoint; the argument below does not assert their
unconditional production.

Let $J$ be the normalized Jacquet representation of $\pi$,
on the same vector space as the coinvariants in AFS11.  It
has finite length over the diagonal torus.  A smooth simple
representation of $(F^\times)^n$ is one-dimensional: an
open compact subgroup fixing one nonzero vector fixes its
entire cyclic representation, because the torus is abelian;
the resulting quotient group is finitely generated abelian.
Its complex group algebra is a finitely generated
commutative algebra, and a simple quotient is $\mathbb C$
by the weak Nullstellensatz.  Thus a finite composition
series makes $J$ finite-dimensional.  This elementary torus
argument applies also in positive residue characteristic.

The nonzero image of $\operatorname{invPart}$ under AFS11
lies in the subspace $E\subset J$ on which the compact
diagonal torus acts by $\varphi^{-1}$; the modulus character
is one on this torus.  This subspace is stable under the
full diagonal torus.  Its commuting endomorphisms have a
common eigencharacter $\mu$, and the corresponding local
block of $J$ has a nonzero quotient functional
$l:J\to\mathbb C$ of character $\mu$ (AFS9's last
finite-dimensional observation).  The character is smooth:
choose a vector on which $l$ is nonzero and an open
stabilizer of that vector.  Write
$\mu(t)=\prod_i\psi_i(t_i)$; then every $\psi_i$ is smooth
and
\[
                 \psi_i(y)=\varphi(e_i(y))^{-1}
                                      \quad(y\in O_F^\times).
\]

For completeness, the needed Frobenius map has the explicit
formula $v\mapsto[g\mapsto l(p(\pi(g)v))]$.
The definition of normalized Jacquet action gives its left
$B$-transformation law $\delta_B^{1/2}\mu$; unipotent factors
act trivially on $p$.  Smoothness follows from a stabilizer
of $v$, and $B\backslash G$ is compact, so this is the
original normalized induced representation.  The map is
nonzero and is equivariant for right translation.  By
irreducibility it embeds $\pi$ into
$\operatorname{Ind}^{\rm norm}_B(\psi)$.

The Jacquet functor is exact here.  To see the only
nontrivial direction, a relation killed in a larger module
uses finitely many upper-unipotent elements.  They lie in
one compact open subgroup of $N_F$; averaging over it as
in AFS11 shows that the same vector is already a sum of
relations in the submodule.  Applying exactness to the
preceding embedding puts $J$ inside the Jacquet module
of this principal series.  The geometric-lemma filtration
has one one-dimensional factor with character
$\prod_i\psi_{w(i)}(t_i)$ for each permutation $w\in S_n$.
Indeed the size matrices for singleton blocks have every
row and column sum one, hence are permutation matrices;
the corresponding zero- and one-dimensional block
inductions simply permute the $\psi_i$, with the normalized
modulus factors canceling.  This verifies the precise
specialization of the supplied geometric-lemma clause.

The functor taking a given smooth compact-torus character
space is exact: twist by its inverse and use the finite
average over a common finite quotient on any proposed
lift.  Its application to that filtration shows that the
$\varphi^{-1}$ subspace has dimension at most the number
of $w$ with
$\psi_{w(i)}|_{O_F^\times}=\psi_i|_{O_F^\times}$ for all $i$.
Distinctness of the given coordinate characters forces
$w=1$.  Therefore $\dim E\leq1$.  It is nonzero, and the
full torus acts on it by $\mu$.  AFS11 now makes
$\operatorname{invPart}$ one-dimensional as well.

Set $\omega_i=\psi_i|\cdot|^{-(n-1)/2}$.  The unit
restrictions remain distinct.  Induction in stages and
$h_{\rm Lit}.\mathrm{loc.disj}$ prove that the principal
series is irreducible: successively add one $\mathrm{GL}_1$
character; no character already present can equal the
new one times $|\cdot|^r$ for $r=-1,0,1$, since these
twists are trivial on units.  The cuspidal supports at
each step are precisely these singleton characters.
The starting $n=0$ representation is the one-dimensional
representation of the trivial group, and $n=1$ is a
character.  Thus the nonzero embedding above is an
isomorphism, proving the asserted principal-series statement.

Finally $[N:a_jNa_j^{-1}]=q^{j(n-j)}$, from the $j(n-j)$
upper entries crossing the two diagonal blocks.  In the
normalized Jacquet module, the torus action gives
\[
 p(\pi(a_j)v)=\delta_B(a_j)^{1/2}
                 \prod_{i<j}\psi_i(\varpi)\,p(v).
\]
Combining with \eqref{eq:afs-jac-sum} yields the scalar
\[
 q^{j(n-j)}\prod_{i<j}
     q^{-(n-1-2i)/2}\,
     \omega_i(\varpi)q^{-(n-1)/2}
   =\prod_{i<j}\omega_i(\varpi)q^{-i}.
\]
The exponent identity is
$j(n-j)-j(n-1)+j(j-1)/2=-j(j-1)/2$.
Both $U_{a_j}v$ and the stated scalar multiple of $v$
remain in $\operatorname{invPart}$ by AFS11, where $p$
is injective.  Their equal Jacquet images therefore give
\eqref{eq:afs-pseig}.  This keeps the original $\varpi$,
operator, normalization and indexing, including $j=0,n$.
\end{proof}
% END AFS12

% BEGIN AFS13
\paragraph{AFS13. Isotypic sums admit separating maps to the original irreducible.}
\label{claim:AFS13}
If $\mathrm{IsoSum}(\rho,Y,\sigma)$ holds in its literal
source sense, then $Y$ is stable under $\rho$, and every
nonzero $F\in Y$ is detected by a complex-linear map
$p:M\to\sigma.V$ satisfying
\[
 p(\rho(g)y)=\sigma.\rho(g)p(y)\quad(y\in Y),
                                  \qquad p(F)\ne0.
\]
\begin{proof}
The source predicate says that $Y$ is spanned by the images
of intertwiners $\sigma\to M$ whose ranges lie in $Y$.
Each nonzero such image is a simple subrepresentation
isomorphic to $\sigma$; this also proves stability of $Y$.
A sum of simple submodules is a direct sum of a suitable
subfamily.  To justify this, choose by the maximal principle
a maximal family whose sum is direct.  Any simple summand
not contained in that sum has zero intersection with it,
by simplicity, and could be added, a contradiction.
Consequently the chosen family spans $Y$.  The nonzero
vector $F$ has a nonzero component in one of its summands.
Projection to that summand followed by an isomorphism to
$\sigma$ detects $F$ and intertwines on $Y$.  Extend this
linear map from $Y$ to $M$ using a vector-space complement.
Only intertwining on $Y$ is asserted, exactly as in the
source's $\mathrm{isoSep}$ statement.
\end{proof}
% END AFS13

% BEGIN AFS14
\paragraph{AFS14. Scalar torus action and surjective positive operators on a local part.}
\label{claim:AFS14}
Let $Y$ be an isotypic sum of a smooth irreducible
$\sigma$ for $\mathrm{GL}_n(F)$ in a complex representation
$\rho$.  Let $0\ne E\subset Y$ be fixed by $N$ and
$N^-(e)$, with
\[
 \rho(e_i(y))F=f_i(y)^{-1}F\quad(F\in E),
 \qquad f_i|_{U_c(F,e+1)}=1,
\]
and suppose the $f_i$ are pairwise distinct as functions
on local units.  Suppose each $U_{a_{j+1}}$ maps $E$ onto
$E$ for $0\leq j<n$.  Under the original $h_{\rm Lit}$,
there are smooth characters $\omega_i$ such that
\[
 \sigma\simeq\mathrm{PS}(\omega),\qquad
 \omega_i(y)f_i(y)=1,
 \qquad U_{a_j}F=\mathrm{psEig}(\omega,j)F
                         \quad(F\in E,\ j\leq n).
\]
\begin{proof}
Every compact diagonal element is a product of the
single-coordinate $e_i(y)$.  Thus it acts as a scalar on
$E$.  Since $E\ne0$, that scalar is unique; it is nonzero
because the group action is invertible.  Comparing its
action on one fixed nonzero vector proves multiplicativity.
The inverse scalars define a character $\varphi$ with
$\varphi(e_i(y))=f_i(y)$.  The specified depth makes
$\varphi$ trivial on $T_c(e)$.

The original identity $z=\prod_{j=0}^{n-1}a_j$ and OST5
show that $U_z$ maps $E$ onto itself.  For any map $p$ in
AFS13, the image of $E$ satisfies the defining eigenlaws
for $\operatorname{eig}(\sigma,e,\varphi)$.  Repeatedly
taking a preimage in $E$ under $U_z$ and using naturality
of the finite sum gives
$p(E)\subset\operatorname{invPart}(\sigma,e,\varphi)$.
AFS13 supplies one $p$ with nonzero image, so AFS12 applies
and gives $\omega$ and its scalar laws on this
$\operatorname{invPart}$.

Fix $F\in E$.  The difference
$U_{a_j}F-\mathrm{psEig}(\omega,j)F$ belongs to $Y$,
since $Y$ is stable under the local group and hence under
the finite sum.  Every AFS13 map kills it, by naturality
and AFS12's scalar identity on $p(E)$.  The separation
in AFS13 forces the difference to be zero.  The assertion
uses all intertwiners only after the single $\omega$ has
been chosen from one nonzero image; their local
$\operatorname{invPart}$ is the same, so no incompatible
character choices occur.
\end{proof}
% END AFS14

% BEGIN AFS15
\paragraph{AFS15. Surjectivity forces integral eigenvalues to be units.}
\label{claim:AFS15}
Let $N$ be a finite module over a commutative ring $A$ and
let $U$ be a surjective endomorphism.  Then
\[
                  1+Uq(U)=0\quad\text{on }N
\]
for some $q\in A[T]$.  If a semilinear map $s$ carries $U$
to a complex operator $T$, its scalar ring map is
$\iota x$, every $|x(a)|\leq1$, and a nonzero eigenvector
of $T$ in $\mathbb C\operatorname{span}s(N)$ has eigenvalue
$\mu$ with $|\iota^{-1}(\mu)|\leq1$, then
$|\iota^{-1}(\mu)|=1$.
\begin{proof}
Choose generators $f_1,\ldots,f_r$ and, by surjectivity,
write $f_i=U\sum_ja_{ij}f_j$.  The adjugate identity for
$I-T(a_{ij})$ gives a polynomial annihilating these
generators when $T=U$.  Its constant coefficient is one,
so it is $1+Tq(T)$.  The identity transfers through $s$
and then to its complex span, because sums, products and
scalar multiplication commute with the semilinear map.
On the nonzero eigenvector it gives
\[
  1+\widetilde\mu\,q_x(\widetilde\mu)=0,
              \qquad\widetilde\mu=\iota^{-1}(\mu).
\]
Every coefficient of $q_x$ has norm at most one, and so
does $\widetilde\mu$.  Hence $|q_x(\widetilde\mu)|\leq1$.
Taking norms in the identity gives
$1\leq|\widetilde\mu|$, proving equality.  This determinant
argument requires neither freeness of $N$ nor flatness
of the coefficient map.
\end{proof}
% END AFS15

% BEGIN AFS16
\paragraph{AFS16. Original Hecke operators preserve every original part.}
\label{claim:AFS16}
For every original $H$, fixed $\iota,\lambda,\xi$, and
$Y:H.\mathrm{PartAltAlt}(\iota,\lambda,\xi)$, each
$H.\operatorname{opC}(\iota,\lambda,d,i)$ preserves $Y.Y$.
On vectors fixed by the product $N$, the positive operator
supported at a place $v$ is exactly the local $U_{a_j}$
at $v$, with the same finite sum and index.
\begin{proof}
AFS13 applied to each $Y.\mathrm{isoL}(v)$ shows that
$Y.Y$ is stable under every single local factor.  Any
element of the finite product is a product of these
single-factor elements, so the full product action
preserves $Y.Y$.  A positive Hecke operator is a finite
sum of these group actions, multiplied by the prescribed
$\operatorname{lowVal}^{-1}$; hence it preserves $Y.Y$.
A spherical generator acts as a scalar on all of $Y.Y$
by the literal $Y.\mathrm{sph}$ field.  This proves the
first assertion for both summands of $H.\mathrm{Op}$.

For the second, take representatives $R$ for the local
quotient $N_v/a_jN_va_j^{-1}$.  Insert each representative
in the $v$th factor and put identity in every other factor.
They represent the product quotient, because the conjugate
subgroup is unchanged at the other factors.  The original
finite sum is thus term by term the local sum on any
product-$N$-fixed vector.  This is precisely the source's
\texttt{u\allowbreak Op\_\allowbreak factor} identity; no normalization or
replacement uniformizer has been introduced.
\end{proof}
% END AFS16

% BEGIN AFS17
\paragraph{AFS17. Full original good-scalar endpoint.}
\label{claim:AFS17}
Retain the complete original hypotheses of
\texttt{good\allowbreak Scal\allowbreak B\_\allowbreak alt}: $h_{\rm Lit}$, $H.\mathrm{Ax}$,
$A$ of class $\mathrm{IsCO}$, $H.\mathrm{IsWt}(\theta)$,
the fixed $\iota$, a given point $x:A\to_O
\overline{\mathbb Q}_\ell$, and any supplied original
$\mathrm{IsSpec}(\theta,\iota,x,\lambda,\xi,c_0,\operatorname{sp})$.
Suppose $H.\mathrm{Good}(\theta,x,\lambda)$, fix an original
part $Y$, level $K$, and $e\geq c_0$, and assume
\[
 E=Y.Y\cap H.\mathrm{OrdAt}(\theta,\iota,x,\lambda,\xi,K,e)
                                                        \ne0.
\]
There are smooth characters $\omega_{v,i}:F_v^\times\to
\mathbb C^\times$ such that
\begin{gather*}
 Y.\sigma_l(v)\simeq\mathrm{PS}(\omega_v),\qquad
 \omega_{v,i}(y)\,
       \iota(\operatorname{finPart}_{v,i}(y))=1,\\
 \mu_{v,j}=\operatorname{lowVal}(\iota,\lambda,v,j+1)^{-1}
                         \mathrm{psEig}(\omega_v,j+1),\qquad
 |\iota^{-1}(\mu_{v,j})|_\ell=1,\\
 H.\operatorname{opC}(\iota,\lambda,d,(v,j))F
                   =\mu_{v,j}F\quad(F\in E, 0\leq j<n).
\end{gather*}
\begin{proof}
Put $S=H.\mathrm{OrdAt}(\theta,\iota,x,\lambda,\xi,K,e)$.
It is finite-dimensional because it lies in the classical
space, whose finiteness is an $\mathrm{IsSpec}$ field.
Its span field identifies $S$ with the complex span of
$\operatorname{sp}(N_0)$ for the original full ordinary
module $N_0$ of AFS10.  The exact ordinary theorem OST25
makes every positive $U_a$ an automorphism of $N_0$.
The original specialization operator identity therefore
makes every normalized local positive operator $T_{v,j}$
map $S$ onto $S$.  In particular it is invertible there.

AFS16 makes $Y.Y$ stable under $T_{v,j}$.  Thus its
restriction to $E=Y.Y\cap S$ is injective.  Since $E$ is
finite-dimensional, it is surjective onto $E$ as well.
The factor $\operatorname{lowVal}(v,j+1)$ is nonzero,
being a product of field images of a uniformizer.
AFS16's product-to-local identity now says that the
unnormalized local $U_{a_{j+1}}$ maps $E$ onto itself.

Membership in the literal classical space gives, at each
$v$, the $N_v$ and $N_v^-(e)$ fixedness and the equations
\[
 \rho_v(e_i(y))F=
             \iota(\operatorname{finPart}_{v,i}(y))^{-1}F.
\]
The supplied specialization is trivial on $U_c(F_v,c_0+1)$;
since $e\geq c_0$, it is also trivial on $U_c(F_v,e+1)$.
The good-point condition says that these $i$-indexed
finite parts are pairwise distinct.  The fixed $\iota$
preserves that distinction.  Apply AFS14 with the literal
$Y.\mathrm{isoL}(v)$ at each place.  It supplies $\omega_v$,
the principal-series isomorphism, and the scalar identity.
Multiplication by the unchanged $\operatorname{lowVal}^{-1}$
gives the displayed $\mu_{v,j}$.

Choose any nonzero $F_0\in E$.  The integrality field of
the given $\mathrm{IsSpec}$ shows
$|\iota^{-1}(\mu_{v,j})|\leq1$, because $F_0$ lies in
the original classical space at the same $K,e$.
Apply AFS15 to $N_0$, its surjective original local
operator, and the same specialization map.  ASP2 gives
$|x(a)|\leq1$ for every $a\in A$.  Thus the norm is one.
The proof has preserved every $K,e\geq c_0$, every
supplied specialization datum, and the source index
$j+1$; it does not impose a domain or flatness condition
on $A$.
\end{proof}
% END AFS17

\subsection{The original action on arithmetic fibers and separation}

% BEGIN AFS18
\paragraph{AFS18. Tensor transport and projection preserve the actual part decomposition.}
\label{claim:AFS18}
Let $N=eN_0$ be an idempotent summand of a finite free
$A$-module, where $e$ belongs to the algebra generated
by specified endomorphisms $U_i$ of $N_0$.  Let a map
$s:N_0\to W$ be semilinear for $\chi:A\to K$, with
$K$ a field, and suppose the images of one $A$-basis of
$N_0$ are $K$-linearly independent.  Suppose
$sU_i=T_is$, where $T_i$ are endomorphisms of $W$.
Set $O'=K\operatorname{span}s(N_0)$ and
$E=K\operatorname{span}s(N)$.  If
\[
 O'=\sum_j(Z_j\cap O')
\]
and every $Z_j$ is stable under all $T_i$, then
\[
 E=\sum_j(Z_j\cap E).
\]
Moreover there is one ring action
\[
 \operatorname{act}:T\longrightarrow\operatorname{End}_K(E),
 \quad T=A[U_i|_N]\subset\operatorname{End}_A(N),
\]
which carries $s(f)$ to $s(tf)$, carries each generator
to $T_i|_E$, and sends $a\in A$ to multiplication by
$\chi(a)$.  Its construction is valid for arbitrary,
possibly nonflat $\chi$.
\begin{proof}
Independence gives the tensor isomorphism
$\Psi_0:K\otimes_A N_0\simeq O'$ by
$k\otimes f\mapsto ks(f)$.  The splitting of $N\subset N_0$
survives every tensor product, so its restriction gives
$\Psi:K\otimes_A N\simeq E$.  Define
\[
 \operatorname{act}(t)=\Psi(1\otimes t)\Psi^{-1}.
\]
It is a ring action, and its stated identities follow
on $s(N)$ and hence on their span.  The same construction
on $N_0$ gives the endomorphism $e_K$ of $O'$.  It fixes
$E$, has image $E$, and is idempotent, by the original
identities for $e$ before base change.

An element of the algebra generated by the $U_i$
transports to a finite polynomial expression in the
$T_i$ and the scalars $\chi(a)$.  It preserves each
$Z_j\cap O'$.  In particular, applying $e_K$ to a finite
expression $F=\sum_jF_j$ with $F_j\in Z_j\cap O'$ and
$F\in E$ gives
$F=\sum_je_KF_j$, with $e_KF_j\in Z_j\cap E$.
This proves the asserted decomposition.  Tensor
injectivity was obtained from the split inclusion,
not inferred from a flatness assumption on $K$.
\end{proof}
% END AFS18

% BEGIN AFS19
\paragraph{AFS19. Full original part-action endpoint on each arithmetic fiber.}
\label{claim:AFS19}
Retain every hypothesis of \texttt{part\allowbreak Act0\_\allowbreak alt}, including
the original $H.\mathrm{AutomorphicHeckeSet}(\iota,P)$,
$\lambda\in H.\mathrm{admitted}$, any supplied
$\mathrm{IsSpec}$ datum at $x$, any residual label $a$,
any level $K$, and the stated commutativity on $H.MK$.
Write $N_0,N,T$ as in AFS10 and
$E=\operatorname{span}_{\mathbb C}\operatorname{sp}(N)$.
Then
\[
 E=\sum_{Y:\,P(\lambda,Y.r)}(Y.Y\cap E),
\]
which in particular proves the source's disjunction of
this equality with $\neg H.\mathrm{Good}(\theta,x,\lambda)$.
There is a single ring action
$\operatorname{act}:T\to\operatorname{End}_{\mathbb C}(E)$
with all the following properties:
\begin{enumerate}
\item For every $t\in T$ and $f\in N$,
$\operatorname{act}(t)\operatorname{sp}(f)=\operatorname{sp}(tf)$.
\item It sends every original generator to the restriction
of the same $H.\operatorname{opC}$, and sends $b\in A$
to multiplication by $\iota(x(b))$.
\item Every original $O$-algebra point $y:T\to
\overline{\mathbb Q}_\ell$ extending $x$ has a nonzero
simultaneous eigenvector $F$ in
$Y.Y\cap H.\mathrm{OrdAt}(\theta,\iota,x,\lambda,\xi,K,c_0)$
for some original part $Y$ with $P(\lambda,Y.r)$.  Its
generator eigenvalues are exactly $\iota(y(\mathrm{gen}_i))$.
\item If $x$ is good and $Y$ is any original part with
$Y.Y\cap E\ne0$, there is one $O$-algebra point
$y:T\to\overline{\mathbb Q}_\ell$ extending $x$ such
that, for every $t\in T$ and every $F\in Y.Y\cap E$,
$\operatorname{act}(t)F=\iota(y(t))F$.
\end{enumerate}
\begin{proof}
OHA6--OHA7 supply the original idempotent $e$ with
$N=eN_0$, its membership in the full ordinary Hecke
algebra, and finite freeness.  ASP's basis statement
(or the corresponding field of the supplied
$\mathrm{IsSpec}$) gives the tensor injectivity required
in AFS18.  Its operator field gives the exact
intertwining of every generator.  Its span field says
\[
 O'=\operatorname{span}_{\mathbb C}\operatorname{sp}(N_0)
         =H.\mathrm{OrdAt}(\theta,\iota,x,\lambda,\xi,K,c_0).
\]
CFD12, applied to the supplied original
$\mathrm{AutomorphicHeckeSet}$, decomposes $O'$ as the
sum of the $Y.Y\cap O'$ for precisely the parts with
$P(\lambda,Y.r)$.  AFS16 gives their generator stability.
AFS18 now gives the displayed decomposition of $E$ and
one action with the first two properties.

For a specified point $y$, AFS9 applies to the faithful
finite $T$-module $N$ and the character $\iota\circ y$.
Transport its fiber through $\Psi$ from AFS18.  The
second assertion of AFS9 uses the just-proved decomposition
to place a nonzero eigenvector in one $Y.Y\cap E$.
Since $E\subset O'$, it lies in the required ordinary
space at depth $c_0$.  This gives property 3, including
the source's existential depth by choosing exactly $c_0$.

For property 4, fix $Y$ with $Z=Y.Y\cap E\ne0$.
Spherical generators act as scalars on all of $Y.Y$.
AFS17, applied to $Y.Y\cap O'\ne0$, makes every local
positive generator scalar on $Y.Y\cap O'$ and thus on
$Z$.  Consequently every element of $T$, being a
finite algebra expression in these generators and
$A$, acts as a scalar on $Z$.  Choose once a nonzero
$F_0\in Z$.  A scalar acting on all $Z$ is unique, since
its action on $F_0$ determines it.  Denote that scalar
by $z(t)$.  The identities for $0,1,+,\cdot$ in the
ring action prove that $z:T\to\mathbb C$ is a ring
homomorphism, and the coefficient identity proves
$z|_A=\iota\circ x$.  Put $y=\iota^{-1}\circ z$.
The original $O$-algebra structures commute because
they do so on $A$.  Thus $y$ is exactly the required
$O$-algebra point, chosen before all $t$ and $F$ in
property 4.  This verifies every clause of the literal
source endpoint, including parts not in the selected
predicate family in its last clause.
\end{proof}
% END AFS19

% BEGIN AFS20
\paragraph{AFS20. Original good arithmetic points separate the original Hecke algebra.}
\label{claim:AFS20}
Under every original hypothesis of
\texttt{TCore\allowbreak Sep.\allowbreak eq\_\allowbreak zero\_\allowbreak of\_\allowbreak vanish\_\allowbreak at\_\allowbreak good\_\allowbreak points\_\allowbreak dense\_\allowbreak weight},
fix $\lambda\in H.\mathrm{admitted}$ with
$H.\mathrm{DenseWt}(\theta,\lambda)$.  For every $t\in
T=H.TK(\theta,a,K)$,
\[
 \left[
 \forall y:T\to_O\overline{\mathbb Q}_\ell,
 \ H.\mathrm{Good}(\theta,y|_A,\lambda)\Rightarrow y(t)=0
 \right]\quad\Longrightarrow\quad t=0.
\]
This retains an arbitrary original $A$ of class
$\mathrm{IsCO}$ and the supplied automorphic-part hypothesis.
It does not presume
reduced fibers or a flat point map.
\begin{proof}
Choose an $A$-basis $(b_i)$ of the original full ordinary
module $N_0$, using OHA7.  To show $t=0$ as its specified
endomorphism of $N=H.MK$, fix $f\in N$ and put $g=tf$,
viewed in $N_0$.  Fix any good point $x:A\to_O
\overline{\mathbb Q}_\ell$.  ASP13 supplies the exact
specialization datum at this $x$ and $\lambda$.  Apply
AFS19 to its space $E$.

Each nonzero piece $Y.Y\cap E$ has by AFS19(4) one
point $y$ extending $x$ on which the entire $T$-action
is scalar.  The hypothesis gives $y(t)=0$, so
$\operatorname{act}(t)$ kills that piece.  It also kills
zero pieces.  The decomposition in AFS19 therefore
shows that it kills all of $E$.  Its compatibility on
$\operatorname{sp}(f)$ gives $\operatorname{sp}(g)=0$.

Write $g=\sum_i a_i b_i$ in the chosen basis.  This is
a finite sum.  Additivity and semilinearity of the same
specialization map give
\[
             0=\sum_i\iota(x(a_i))\operatorname{sp}(b_i).
\]
Its basis-independence clause implies $x(a_i)=0$ for
every $i$.  The point $x$ was arbitrary among all good
points at the fixed $\lambda$.  DenseWt makes every
$a_i$ zero, hence $g=0$.  Since $f$ was arbitrary,
the original endomorphism $t$ is zero.

In particular this proves the exact $H.\mathrm{Out.sep}$
field without requiring the other $\mathrm{Out}$ fields.
For example it also implies that $T$ is reduced under
these hypotheses: a nilpotent element has zero image
under every field point, and the result just proved
annihilates it.  This reducedness is a consequence of
the scalar-part argument, not an input to it.
\end{proof}
% END AFS20

% BEGIN AFS21
\paragraph{AFS21. Separation descends to the very same kernel ideal.}
\label{claim:AFS21}
Let $O,A,B,T,E$ be commutative rings, with their stated
$O$-algebra structures, let $T$ also be an $A$-algebra,
and let
\[
 \lambda:A\to_O B,\qquad f:B\to_O T,
          \qquad f\lambda=\operatorname{struct}_{A,T}.
\]
Let $P$ be a predicate on ring maps $A\to E$ and $G$ a
predicate on $O$-algebra maps $B\to E$.  Assume
$P(x\lambda)\Rightarrow G(x)$ for every such $x$, and
assume the $O$-algebra points $z:T\to E$ with
$P(z|_A)$ separate $T$.  If $b\in B$ vanishes at
every point $x:B\to_O E$ with
$\ker f\subset\ker x$ and $G(x)$, then $b\in\ker f$.
\begin{proof}
For any separating point $z:T\to_O E$, let $x=z\circ f$.
It kills $\ker f$, and
$x\lambda=z\circ\operatorname{struct}_{A,T}$ has property
$P$.  Thus $G(x)$ holds and the hypothesis on $b$ gives
$z(f(b))=x(b)=0$.  Separation in $T$ yields $f(b)=0$.
This is exactly $b\in\ker f$.  Neither surjectivity of
$f$ nor flatness of any map was used.  The conclusion
keeps the literal ideal $\ker f$, without passing to
its radical, a larger support ideal, or a different
quotient.
\end{proof}
% END AFS21

% BEGIN AFS22
\paragraph{AFS22. Exact density on an original IsJ ideal, with the remaining automorphic supplier exposed.}
\label{claim:AFS22}
Fix all the original data of \texttt{TJDense}, including
$h_P:\mathrm{HypothesesH'}$, the prime $\ell$, CM field
$L$, algebraic closure, Hida setting $s$, $n\geq1$,
$V,Q,c,r,\iota$, the supplied original determinant ring
$B$, weight ring $\Lambda$, and
$R:\mathrm{UnitaryHeckeData}(Q,\iota,\Lambda)$.
Retain $V.S'=s.S$ and any original auxiliary level $K$.
The following additional installation input is kept
explicit:
\begin{equation}\label{eq:afs-aut-input}
 H.\mathrm{AutomorphicHeckeSet}
       (\iota,\mathrm{UnitaryPointProperties}(r,\iota)),
             \qquad H=\mathrm{hset}(Q,R.\tau_0).
\end{equation}
Then for every original ideal $J\subset B.B$,
\[
            R.\mathrm{IsJ}(B,K,J)\quad\Longrightarrow\quad
                                      B.\mathrm{Dense}(J).
\]
In full quantified form, for every increasing original
weight $h$ and every $b\in B.B$,
\[
 \bigl[\forall x:B.\mathrm{Pt},\ J\subset\ker x
              \ \land\ x.\mathrm{Good}(h)\Rightarrow x(b)=0\bigr]
                                      \quad\Longrightarrow\quad b\in J.
\]
\begin{proof}
The literal $R.\mathrm{IsJ}$ has two cases.  In its
zero-module case $J=\top$, so the conclusion is immediate.
In its other case it supplies the original commutativity
and coefficient structures, a surjective $O$-algebra map
\begin{gather*}
 f:B.B\longrightarrow T=H.TK(R.\theta,R.a,K),\qquad J=\ker f,\\
 f\circ\Lambda.\mathrm{toPair}(B.u)
                          =\operatorname{struct}_{\Lambda,T},
\end{gather*}
together with the stated universal determinant and
character compatibility.  These maps and this ideal
are retained exactly; the density argument needs only
the displayed compatibility from that supplied record.

Fix $h$ with $h.\mathrm{Incr}$.  AFS6--AFS8 give one
$\lambda$ with $\mathrm{WtOf}(h,\lambda)$ and the exact
equivalence between $x.\mathrm{Good}(h)$ and goodness
of its original weight restriction.  AFS5 gives
$H.\mathrm{DenseWt}(R.\theta,\lambda)$.
FNL16 supplies $H.\mathrm{Ax}$, and the supplied $R$
already has its literal $\mathrm{IsWt}$ field.
Apply AFS20 using \eqref{eq:afs-aut-input}.  It gives
separation of this exact $T$ by the good points at
$\lambda$.  Now apply AFS21 with
$A=\Lambda.\Lambda$, $\lambda=\Lambda.\mathrm{toPair}(B.u)$,
$P$ the HSet good-point predicate, and $G$ the original
$x.\mathrm{Good}(h)$.  AFS8 verifies the required
predicate implication.  The conclusion is
$b\in\ker f=J$, as required, for this same $J$ and
every $h,b$ in the source's $\mathrm{Dense}$ definition.

The source proves \eqref{eq:afs-aut-input} from the
type field of $R$ through its actual automorphic-part
and Galois-attachment constructions.  The present
chapter does not replace those producers by the
existence of a formal proof.  Thus AFS22 proves the
complete same-ideal density implication with this
exact remaining supplier, while the unconditional
installation of the bare \texttt{TJDense} endpoint
remains open until that supplier is reconstructed.
It neither asserts the other $\mathrm{Out.pt}$ fields
nor closes the complete Hida construction.
\end{proof}
% END AFS22


\section{The remaining finite local laws on the original labels}
\label{sec:finite-local-completion-v43}

\status{The first part proves the full original root-injectivity
collision estimate, with its constant $23$, for every finite field
of cardinality at least $2^{12}$ and every degree.  The second part
works with the original finite projectors and normalized polynomial
labels.  Every use of Green's published cuspidal theorem is stated
explicitly.  The finite local statements do not themselves construct
the adelic rows or the arithmetic cylinder counts.}

\subsection{The fifth local error term}

Let $K$ be a finite field of cardinality $q$, and let $L$ be an
algebraically closed extension of $K$.  Write $\mathcal S_N$ for the
monic squarefree polynomials $F\in K[X]$ of degree $N$ with
$F(0)\ne0$, and put
\[
 w(F)=\prod_{P\mid F}(q^{\deg P}-1)^{-1},\qquad
 \mu_N(A)=\sum_{F\in\mathcal S_N\cap A}w(F).
 \tag{FLC.1}
\]
The factors $P$ are monic irreducibles.  This is exactly the
source's \texttt{sq\allowbreak Weight}, with the coefficient-vector
parametrization \texttt{monic\allowbreak Of} suppressed.  RRT5 proves
$\mu_N(\{F:H\mid F\})\le w(H)$ for monic squarefree $H$ with
nonzero constant term and $\deg H\le N$.  The same event is empty
if $\deg H>N$ or $H$ has a repeated irreducible factor.
For a prime $\ell$ and a nonzero finite-order element $z$, write
$z_{\ell'}$ for its prime-to-$\ell$ part.  Thus in the cyclic group
of order $\ell^a m$, with $(\ell,m)=1$, this is the projection
onto its subgroup of order $m$.  It commutes with every power map.
The original exceptional set is
\[
 \mathcal B(q,N)=\{q\}\cup
 \{\ell\text{ prime}:\exists\,1\le d\le N,
          q^d<\ell^{4v_\ell(q^d-1)}\}.
 \tag{FLC.2}
\]
Its finite-range implementation is identical: a prime satisfying
the displayed strict inequality divides $q^d-1$ and is at most
$q^N$.  No primality assumption on $q$ is needed in this subsection.

% BEGIN FLC1
\paragraph{FLC1. Frobenius orbits and the exact root-fiber bound.}
\label{claim:FLC1}
Every monic irreducible $P\in K[X]$ of degree $d$ with $P(0)\ne0$
has exactly $d$ distinct roots in $L$.  These roots form one orbit
under $z\mapsto z^q$, and each satisfies $z^{q^d-1}=1$.
For any integer $n\ge1$, prime $\ell$, and $a\in L^\times$,
\[
 \#\{z\in L:z^n=1,\ z_{\ell'}=a\}
                      \le\ell^{v_\ell(n)}.
 \tag{FLC.3}
\]
\begin{proof}
A root $z$ generates a field of $q^d$ elements over $K$.
The Frobenius orbit polynomial is fixed by Frobenius, hence has
coefficients in $K$, so its length is at least the degree of the
minimal polynomial.  Frobenius to the $d$th power is the identity
on this field.  Thus the orbit length is exactly $d$ and exhausts
the roots.  Finite fields are perfect: the $p$th-power map is an
injection and therefore a surjection on a finite field.  This also
proves separability.  Its nonzero elements have order dividing
$q^d-1$ by Lagrange's theorem.

If the set in (FLC.3) is nonempty, choose $z_0$ in it.
For every other member, $z/z_0$ has $\ell$-power order because the
prime-to-$\ell$ projections agree; its order also divides $n$.
It is therefore a root of $X^{\ell^{v_\ell(n)}}-1$.
The injective map $z\mapsto z/z_0$ and the polynomial root bound
prove (FLC.3), including the case where $\ell$ is the characteristic.
\end{proof}
% END FLC1

% BEGIN FLC2
\paragraph{FLC2. The full same-factor count.}
\label{claim:FLC2}
Let $d\ge2$, let $\ell$ be prime, and assume
$\ell^{4v_\ell(q^d-1)}\le q^d$.
The number of monic irreducible degree-$d$ polynomials for which
$z\mapsto z_{\ell'}$ is not injective on their roots is at most
$q^{3d/4}$.  This is the exact original \texttt{Same\allowbreak Factor\allowbreak Count}.
\begin{proof}
Put $a=\ell^{v_\ell(q^d-1)}\le q^{d/4}$.
An irreducible of degree at least two has nonzero constant term.
For a bad irreducible choose a root $z$ and a distinct conjugate
$z^{q^j}$ with equal projections.  The Frobenius orbit of
$z_{\ell'}$ has a proper length $e\mid d$, so $1\le e\le d/2$.
Consequently
\[
 (z_{\ell'})^{q^e-1}=1,\qquad z^{a(q^e-1)}=1.
 \tag{FLC.4}
\]
The second equation holds since the $\ell$-primary part of $z$
has order dividing $a$.  All $d$ roots of this irreducible
satisfy the same equation.  Different monic irreducibles have
disjoint root sets.  The union, over all proper positive divisors
$e$ of $d$, of the roots in (FLC.4) has cardinality at most
\[
 \sum_{\substack{e\mid d\\1\le e<d}} a(q^e-1)
                   \le d\,q^{d/4}q^{d/2}.
\]
Divide by $d$ to obtain the assertion.  The argument counts all
bad factors, without selecting a distinguished root of each one.
\end{proof}
% END FLC2

% BEGIN FLC3
\paragraph{FLC3. The full cross-factor count.}
\label{claim:FLC3}
Fix any irreducible $P'\in K[X]$ with $P'(0)\ne0$.
Let $d\ge1$ and $\ell$ be prime with
$\ell^{4v_\ell(q^d-1)}\le q^d$.
At most $q^{d/4}$ monic irreducible degree-$d$ polynomials $P$
with $P(0)\ne0$ have roots $z\in P^{-1}(0)$ and
$y\in(P')^{-1}(0)$ satisfying $z_{\ell'}=y_{\ell'}$.
This is the exact original \texttt{Cross\allowbreak Factor\allowbreak Count}.
\begin{proof}
Fix a root $y_0$ of $P'$.  By FLC1, for each such pair $(z,y)$
some common Frobenius power carries $y$ to $y_0$ and carries $z$
to another root of the same $P$.  This root lies in
\[
 T=\{x:x^{q^d-1}=1,\ x_{\ell'}=(y_0)_{\ell'}\}.
\]
Each relevant $P$ has a root in $T$, and roots belonging to
different monic irreducibles cannot coincide.  Therefore their
number is at most $|T|$.  FLC1 gives
$|T|\le\ell^{v_\ell(q^d-1)}\le q^{d/4}$.
No factor $\deg P'$ is lost: the common Frobenius transport is
what permits the single fixed fiber $T$.
\end{proof}
% END FLC3

% BEGIN FLC4
\paragraph{FLC4. Weighted coverage of both collision mechanisms.}
\label{claim:FLC4}
Let $A_N$ be the event that, for some prime
$\ell\notin\mathcal B(q,N)$, two distinct roots of $F$ have equal
prime-to-$\ell$ projections.  Then
\begin{align}
 \mu_N(A_N)&\le A+B,\notag\\
 A&=\sum_{d=2}^{N}\omega(q^d-1)
                \frac{q^{3d/4}}{q^d-1},\notag\\
 B&=\sum_{d=1}^{N}\sum_{e=1}^{d}
       \bigl(\omega(q^d-1)+\omega(q^e-1)\bigr)
                \frac{q^{d/4}}{q^d-1}.
 \tag{FLC.5}
\end{align}
Here $\omega(n)$ is the number of distinct prime factors of $n$.
\begin{proof}
Factor $F$ into its distinct monic irreducibles.  If both roots
belong to one factor of degree $d$, then $d\ge2$.  Also
$\ell\mid q^d-1$, since otherwise the projection is the identity
on all of its roots.  Being outside $\mathcal B(q,N)$ gives
the hypothesis of FLC2 for this $d$.  For each bad factor $P$,
RRT5 bounds the total weight of its multiples in $\mathcal S_N$
by $(q^d-1)^{-1}$.  Union bounds over factors and primes give $A$.

Otherwise the roots belong to distinct factors $P,P'$ of
respective degrees $d\ge e\ge1$.  Their collision forces
$\ell\mid(q^d-1)(q^e-1)$, since otherwise both projections are
identities.  Again the hypothesis of FLC3 holds for the larger
degree $d$.  For fixed $P'$ and $\ell$, FLC3 gives at most
$q^{d/4}$ possibilities for $P$.  The divisor bound assigns
weight at most $((q^d-1)(q^e-1))^{-1}$ to each distinct pair.
The number $I_e$ of monic irreducibles of degree $e$ with
nonzero constant term satisfies
$eI_e\le q^e-1$, by the disjoint-root count in FLC1; in
particular $I_e\le q^e-1$.  Summing first over $P'$, and then
over the at most $\omega(q^d-1)+\omega(q^e-1)$ possible primes,
gives the corresponding term of $B$.

Pairs with $P=P'$ or with $d+e>N$ contribute the empty event.
They can be omitted before applying the squarefree divisor bound.
Allowing them when bounding the number of pairs only increases
its nonnegative upper bound.  Thus no divisor estimate for a
nonsquarefree polynomial is being assumed.
\end{proof}
% END FLC4

% BEGIN FLC5
\paragraph{FLC5. The original constant $23$ and all original quantifiers.}
\label{claim:FLC5}
For every finite field $K$ with $q=|K|\ge2^{12}$, every
algebraically closed extension $L/K$, and every $N\ge0$,
\[
 \sum_{\substack{F\in\mathcal S_N\\
       \exists\ell\notin\mathcal B(q,N)\text{ prime}:\\
       z\mapsto z_{\ell'}\text{ not injective on roots of }F}}
       w(F)\le\frac{23\log q}{\sqrt q}.
 \tag{FLC.6}
\]
This proves the full original \texttt{Local\allowbreak Estimate\allowbreak II}.
\begin{proof}
For $d\ge1$, the product of the distinct prime factors of
$q^d-1$ is at least $2^{\omega(q^d-1)}$ and divides $q^d-1$.
Hence
\[
 \omega(q^d-1)\le\frac{d\log q}{\log2}
                         \le\frac32d\log q.
\]
The numerical inequality $\log2>2/3$ follows, for example,
by integrating the strictly convex function $1/t$ on $[1,2]$:
its integral strictly exceeds its midpoint value $2/3$.
Put $x=q^{-1/4}\le1/8$.  Since $q^d-1\ge q^d/2$,
FLC4 gives
\[
 A\le3\log q\sum_{d=2}^N d x^d,\qquad
 B\le6\log q\sum_{d=1}^N d^2(x^3)^d.
 \tag{FLC.7}
\]
For $0\le x\le1/8$ and $0\le y\le1/512$, the elementary
inequalities $d\le2^d$ and $d^2\le4^d$ give
\[
 \sum_{d=2}^N d x^d\le\frac{4x^2}{1-2x}
                          \le\frac{16}{3}x^2,\qquad
 \sum_{d=1}^N d^2y^d\le\frac{4y}{1-4y}\le5y.
\]
Taking $y=x^3$ yields
\[
 A+B\le(16x^2+30x^3)\log q
       \le\frac{79}{4}x^2\log q
       \le23x^2\log q.
\]
Since $x^2=q^{-1/2}$, this is (FLC.6).  Empty degree ranges
cover $N=0$ and $N=1$ without an extra lower-rank hypothesis.
The definition of $\mathcal B$ and the canonical projection in
(FLC.2) are the original definitions, so this is the original
endpoint on the original polynomial space, not a substitute
collision event.
\end{proof}
% END FLC5

\subsection{Finite representation calculations with the original characters}

For the rest of the chapter, $q$ is prime and
$G_n=\operatorname{GL}_n(\mathbb F_q)$.  All representations are
finite dimensional over $\mathbb C$.  Fix a nontrivial additive
character $\psi$ of $\mathbb F_q$, and write
\[
 U_n=\{\text{upper unitriangular matrices}\},\qquad
 \lambda_n(u)=\psi\left(\sum_{a=0}^{n-2}u_{a,a+1}\right),\qquad
 \operatorname{Wh}(V)=V^{U_n,\lambda_n}.
 \tag{FLC.8}
\]
The superscript means the subspace on which $u$ acts as
$\lambda_n(u)$.  Its projector is exactly the $e_U$ of RRT1.
The rank-zero group and its character are trivial.  In particular
its Whittaker functor is the identity, with no exceptional sign.

% BEGIN FLC6
\paragraph{FLC6. The finite induction identities used below.}
\label{claim:FLC6}
Let $G$ be finite and let $H,K$ be subgroups.  For complex
representations $A$ of $H$ and $B$ of $K$,
\[
 \operatorname{Hom}_G(\operatorname{Ind}_H^G A,
                          \operatorname{Ind}_K^G B)
 \simeq\bigoplus_{x\in H\backslash G/K}
 \operatorname{Hom}_{H\cap xKx^{-1}}(A,{}^xB).
 \tag{FLC.9}
\]
This includes Frobenius reciprocity when $H=G$.
All representations are semisimple, and for a character
$\lambda:H\to\mathbb C^\times$,
\[
 \dim V^{H,\lambda}
       =|H|^{-1}\sum_{h\in H}\lambda(h)^{-1}\chi_V(h).
 \tag{FLC.10}
\]
\begin{proof}
A representation admits an invariant positive Hermitian form:
average any such form over $G$.  Orthogonal complements then
prove semisimplicity.  Realize an induced module as functions
$f:G\to A$ with $f(hg)=h f(g)$, with the right translation
action.  On restriction to $K$, separate these functions
according to their supports in the double cosets $HxK$.
Evaluation at $x$ identifies each summand with induction from
$K\cap x^{-1}Hx$ of the transported representation.  A map
from an induced module is uniquely determined by its values
on functions supported in $H$; conversely a map commuting with
$H$ extends to every coset by translation.  This proves
reciprocity, and applying it to the displayed support
decomposition proves (FLC.9), after replacing representatives
by their inverses if necessary.

The operator $|H|^{-1}\sum_h\lambda(h)^{-1}h$ is an idempotent
with image $V^{H,\lambda}$: multiplying by any $h_0\in H$
multiplies it by $\lambda(h_0)$, and it is the identity on that
subspace.  Taking its trace proves (FLC.10).
\end{proof}
% END FLC6

% BEGIN FLC7
\paragraph{FLC7. The exact Gelfand--Graev norm.}
\label{claim:FLC7}
For $n\ge1$, put $\Gamma_n=\operatorname{Ind}_{U_n}^{G_n}\lambda_n$.
Then
\[
 \dim\operatorname{Wh}(\Gamma_n)
       =\dim\operatorname{End}_{G_n}(\Gamma_n)
       =q^n-q^{n-1}.
 \tag{FLC.11}
\]
\begin{proof}
Gaussian elimination gives the Bruhat decomposition
$G_n=\coprod_{w\in S_n}U_n T_n w U_n$, and the diagonal factor
$t\in T_n$ is unique in the double quotient by $U_n$.
In detail, the pivot positions determine $w$; elimination using
unitriangular row and column operations then leaves the nonzero
pivot values $t$, which those operations cannot change.
Apply FLC6 with $H=K=U_n$.  The contribution of $tw$ is $1$
precisely when the two characters agree on
$U_n\cap twU_n(tw)^{-1}$, and is $0$ otherwise.

This intersection is generated by the elementary matrices
$1+aE_{ij}$ with $i<j$ and $w^{-1}(i)<w^{-1}(j)$.
The character is nontrivial on such a root subgroup exactly
when its two indices are adjacent.  Equality of the two
characters therefore requires
\[
 \begin{split}
 w(j)<w(j+1)&\ \Longrightarrow\ w(j+1)=w(j)+1,\\
 w^{-1}(i)<w^{-1}(i+1)&\ \Longrightarrow\
                  w^{-1}(i+1)=w^{-1}(i)+1.
 \end{split}
 \tag{FLC.12}
\]
The permutations satisfying these conditions are precisely
concatenations of increasing consecutive intervals, with the
intervals themselves arranged in decreasing order.
Here is an elementary verification.  Cut $w(1),\ldots,w(n)$
at each descent.  The first condition makes every resulting
run a consecutive interval.  If two such intervals occurred
in increasing value order, choose two successive value
intervals whose positional order increases; their boundary
values are consecutive, but occupy nonadjacent positions.
The second condition forbids this.  Thus the intervals occur
in decreasing order.  The converse follows immediately by
checking roots within and between intervals.

If there are $r$ intervals, character compatibility further
requires $t_i=t_{i+1}$ within each interval and imposes no
condition between intervals.  Indeed equality
$\psi(ca)=\psi(a)$ for every $a\in\mathbb F_q$ implies $c=1$:
otherwise $(c-1)a$ runs through the field, contradicting
nontriviality of $\psi$.  Hence this permutation contributes
$(q-1)^r$ compatible diagonal matrices.  Choosing the $r-1$
cuts among $n-1$ positions gives
\[
 \sum_{r=1}^n\binom{n-1}{r-1}(q-1)^r
                         =(q-1)q^{n-1}.
\]
Frobenius reciprocity identifies the endomorphism dimension
with $\dim\operatorname{Wh}(\Gamma_n)$, proving (FLC.11).
\end{proof}
% END FLC7

% BEGIN FLC8
\paragraph{FLC8. The generic mirabolic constituent and a sharp lower bound.}
\label{claim:FLC8}
Let
\[
 P_n=\left\{\begin{pmatrix}h&v\\0&1\end{pmatrix}:
              h\in G_{n-1},\ v\in\mathbb F_q^{n-1}\right\}
               \subset G_n.
\]
Then $R_n=\operatorname{Ind}_{U_n}^{P_n}\lambda_n$ is
irreducible.  Its dimension is
$\prod_{a=1}^{n-1}(q^a-1)$, and
$R_n|_{G_{n-1}}\simeq\Gamma_{n-1}$.
For every representation $W$ of $P_{k+1}$, with $k\ge1$,
\[
 \dim W^{U_k,\lambda_k}\ge
          (q^k-q^{k-1})\dim W^{U_{k+1},\lambda_{k+1}}.
 \tag{FLC.13}
\]
\begin{proof}
Write $V_n$ for the last-column normal subgroup of $P_n$.
Its characters are $v\mapsto\psi(av)$, with
$a\in(\mathbb F_q^{n-1})^\vee$; nontriviality of $\psi$
makes this parametrization injective and hence bijective.
The group $G_{n-1}$ has two orbits on these characters,
namely zero and nonzero.  The stabilizer of
$v\mapsto\psi(v_{n-2})$ is precisely $P_{n-1}$.

For completeness, the representation theory needed for this
observation is elementary.  The commuting finite-order
operators from $V_n$ are simultaneously diagonalizable.
Their eigenspaces are permuted by $G_{n-1}$.  An irreducible
$P_n$-module has a single orbit of nonzero eigenspaces.
Its chosen eigenspace is irreducible for the stabilizer;
otherwise translates of a proper invariant subspace would
produce a proper $P_n$-submodule.  Conversely, inducing an
irreducible stabilizer module, extended by the chosen scalar
character of $V_n$, is irreducible because any submodule is
the direct sum of its eigenspace intersections.  This proves
the orbit construction without invoking a classification.

Now $R_1$ is the trivial one-dimensional module.  Induction
in stages identifies $R_n$ with the nonzero-orbit induction
of $R_{n-1}$, since
$U_n=U_{n-1}\ltimes V_n$ and the last-column character is the
one just chosen.  The orbit construction proves irreducibility
by induction.  Counting the successive nonzero orbits gives
its stated dimension.  Moreover $P_n=G_{n-1}U_n$ and
$G_{n-1}\cap U_n=U_{n-1}$, so the one-double-coset restriction
in FLC6 gives $R_n|_{G_{n-1}}\simeq\Gamma_{n-1}$.

By reciprocity and irreducibility, $R_{k+1}$ is the unique
irreducible $P_{k+1}$-module with a nonzero
$(U_{k+1},\lambda_{k+1})$ space, and that space has dimension
one.  Thus a semisimple $W$ contains exactly
$m=\dim W^{U_{k+1},\lambda_{k+1}}$ copies of $R_{k+1}$.
Upon restriction to $G_k$, these contribute $m$ copies of
$\Gamma_k$.  All other summands contribute nonnegative
Whittaker dimensions.  FLC7 proves (FLC.13).
\end{proof}
% END FLC8

% BEGIN FLC9
\paragraph{FLC9. A finite character-exchange identity.}
\label{claim:FLC9}
Let $H_1,H_2\le H\le G$ be finite groups with
$|H_1|=|H_2|$, and let $\alpha_i:H_i\to\mathbb C^\times$
agree on $H_1\cap H_2$.
Suppose that, for $i=1,2$ and every $h\in H\setminus H_i$,
there are $a,a'\in H_i$ such that
$a'=h^{-1}ah$ and $\alpha_i(a')\ne\alpha_i(a)$.
Then for every class function $\chi$ on $G$,
\[
 |H_1|^{-1}\sum_{a\in H_1}\alpha_1(a)^{-1}\chi(a)
       =|H_2|^{-1}\sum_{b\in H_2}\alpha_2(b)^{-1}\chi(b).
 \tag{FLC.14}
\]
\begin{proof}
For $h\notin H_i$, the sum
$F(h)=\sum_{c\in H_i}\alpha_i(c)^{-1}\chi(ch)$ vanishes.
Indeed replace $c$ by $(a')^{-1}ca$ and use
$a'=h^{-1}ah$ and conjugation invariance; this gives
$F(h)=\alpha_i(a')\alpha_i(a)^{-1}F(h)$ with a scalar
different from one.
Now calculate
\[
 S=\sum_{a\in H_1}\sum_{b\in H_2}
        \alpha_1(a)^{-1}\alpha_2(b)^{-1}\chi(ab).
\]
First fix $b$.  Its contribution vanishes unless
$b\in H_1\cap H_2$.  In that case translation inside $H_1$
and agreement of the characters give the unshifted sum over
$H_1$.  Thus $S$ is $|H_1\cap H_2|$ times that sum.
Reversing the order and using $\chi(ab)=\chi(ba)$ gives the
corresponding expression for $H_2$.  Cancel the positive
intersection cardinality, then divide by the common group
order.  This proves the identity, with its exact normalized
averages.
\end{proof}
% END FLC9

% BEGIN FLC10
\paragraph{FLC10. The literal $e_J$ exchanges to the column-vanishing projector.}
\label{claim:FLC10}
Let $N=k+2s+1$, $k\ge1$, $s\ge0$, and use the original
$e_J=e_J(q,N,k,s,\psi)$ from RRT1.  Define
\[
 J'=\{u\in U_N:u_{a,k}=0\ (0\le a<k)\},\qquad
 e_{J'}=|J'|^{-1}\sum_{u\in J'}\lambda_N(u)^{-1}[u].
\]
Then $\operatorname{tr}(e_J\mid V)
=\operatorname{tr}(e_{J'}\mid V)$ for every representation of
$G_N$.  In fact the normalized averages agree for every class
function.  This is the full original pattern-exchange endpoint.
\begin{proof}
Use two ordered chains
\[
 W=(s,s+1,\ldots,s+k-1),\qquad
 C=(0,\ldots,s-1,s+k,s+k+1,\ldots,N-1),
\]
with entries $W_a$ and $C_b$ numbered from zero.  For
$0\le i\le s$, order the basis as
\[
 C_0,\ldots,C_{i-1},W_0,\ldots,W_{k-1},C_i,C_{i+1},\ldots,C_{2s}.
\]
Let $H_i$ consist of the matrices unitriangular in this order
whose $k$ entries from $W$ into column $C_i$ vanish.
Its permitted entries are all forward entries within $W$ and
within $C$, all $C_b\to W$ with $b<i$, and all $W\to C_b$
with $b>i$.  This is precisely the source pattern
\texttt{ph\allowbreak I k s i}.  It is closed under multiplication: the
forbidden entry has no two-step path through permitted entries.
It has $q^{N(N-1)/2-k}$ elements.  Let $\alpha_i$ sum the
entries along the adjacent edges of the two fixed chains and
apply $\psi$.  No character edge has a two-step path through
permitted entries, so this sum is additive under matrix
multiplication.  These are exactly \texttt{ee\allowbreak I} and its
pattern character.

For $1\le i\le s$, put $p=C_{i-1}$ and $t=C_i$.
The patterns of $H_i$ and $H_{i-1}$ differ only as follows:
\[
 \begin{array}{c|c|c}
       &H_i\setminus H_{i-1}\text{ entries}
       &H_{i-1}\setminus H_i\text{ entries}\\ \hline
       &p\longrightarrow W&W\longrightarrow t.
 \end{array}
 \tag{FLC.15}
\]
The union pattern is a group $H$: insert $p,W,t$ in this
order and allow both displayed families.  The entry $p\to t$
is a character edge.  The two characters agree on the
intersection.

If $h\in H\setminus H_i$, some $h_{w,t}\ne0$ with $w\in W$.
Take $a=1+uE_{p,w}$.  Then $a\in H_i$, its character is one,
and
\[
 h^{-1}ah=1+u(h^{-1}e_p)(e_w^{\mathsf t}h)\in H_i,
 \qquad \alpha_i(h^{-1}ah)=\psi(u h_{w,t}).
\]
To check both assertions, $h^{-1}e_p$ is supported on the
$C$ entries up to $p$, whereas $e_w^{\mathsf t}h$ is supported
on $W$ from $w$ onward and on $C$ from $t$ onward.  All their
products are permitted in $H_i$; the only character edge
among them is $p\to t$, and $(h^{-1})_{p,p}=1$.
Choose $u$ so that the displayed character is not one.
Likewise, if $h\notin H_{i-1}$, then $h^{-1}\notin H_{i-1}$
and some $(h^{-1})_{p,w}\ne0$.  Use $a=1+uE_{w,t}$.
Now its conjugate lies in $H_{i-1}$ and has character
$\psi(u(h^{-1})_{p,w})$, again different from one for some $u$.
The support check is the same with the row and column roles
interchanged.  FLC9 therefore equates the two normalized
character averages.

Iterating from $i=s$ to $i=0$ proves the claim.  At $i=s$ the
order is the original order and $H_s=J$, with exactly
$\phi_J$ from RRT1.  At $i=0$, reordering the basis as $W,C$
identifies $H_0$ with $J'$ and its character with
$\lambda_N|_{J'}$.  Conjugating by this permutation matrix
does not change a class function.  For $s=0$ the single
permutation identification already proves the result.
\end{proof}
% END FLC10

% BEGIN FLC11
\paragraph{FLC11. The exact lower comparison on every original series.}
\label{claim:FLC11}
For every representation $V$ of $G_N$, with $N=k+2s+1$ and
$k\ge1$,
\[
 \dim e_JV\ge(q^k-q^{k-1})\dim e_UV.
 \tag{FLC.16}
\]
Consequently, for every central series projector $E_F$ that
is a sum of primitive central projectors,
\[
           (1-q^{-1})w_{e_U}(F)\le w_{e_J}(F).
 \tag{FLC.17}
\]
In particular this proves the full original
\texttt{lower\allowbreak Wt\allowbreak Comparison}, on its actual \texttt{e\allowbreak Ser\allowbreak C}
projectors, for every label, without Green's theorem.
\begin{proof}
Set $m=k+1$.  Let $T_m\subset U_N$ have identity upper-left
$m\times m$ block, arbitrary upper-right block, and upper
unitriangular lower-right block.  Define its character by
\[
 \beta(t)=\psi\left(\sum_{a=m-1}^{N-2}t_{a,a+1}\right).
\]
For $m=N$ this group and character are trivial.
The group $P_m$, embedded in the upper-left block of $G_N$,
normalizes $T_m$ and preserves $\beta$: its last row is fixed,
so the cross-block coordinate $t_{m-1,m}$ is unchanged.
Thus $W=V^{T_m,\beta}$ is a representation of $P_m$.
The actual semidirect products give
\[
 V^{U_N,\lambda_N}=W^{U_m,\lambda_m},\qquad
 V^{J',\lambda_N|_{J'}}=W^{U_k,\lambda_k}.
\]
Both equalities concern the original subspaces: $U_N=U_mT_m$,
$J'=U_kT_m$, and the indicated characters multiply to the
original adjacent-entry sums.  FLC8 and FLC10 prove (FLC.16).

For (FLC.17), apply (FLC.16) separately to each irreducible
in $E_F$.  Multiply by its positive degree, sum, and use RRT3's
coefficient identity together with
$[1]e_J/[1]e_U=q^k$ from RRT1.  The ratio of multiplicity
factors is $(q^k-q^{k-1})/q^k=1-q^{-1}$.
The argument works for any such central projector, so it
also applies directly to the original normalized projectors
without renaming their labels.
\end{proof}
% END FLC11

\subsection{Squarefree series and the unchanged normalized labels}

For a standard block parabolic $P=MU\subset G_n$, write
$\operatorname{HC}_P(\rho)=\operatorname{Ind}_P^{G_n}\rho$,
where $U$ acts trivially.  This is ordinary finite induction;
there is no modulus character.  An irreducible representation
of $G_d$ is cuspidal when its invariants under every proper
parabolic unipotent radical vanish.  It is enough to require
this for the two-block parabolics: a longer flag radical
contains a conjugate two-block radical.  The source's
\texttt{res a (d-a) chi 1 1=0} says precisely that the
averaging idempotent onto these invariants has trace zero.
Thus this definition is the original \texttt{Is\allowbreak Cusp}.

% BEGIN FLC12
\paragraph{FLC12. Cuspidal degree determines the Whittaker multiplicity.}
\label{claim:FLC12}
For every cuspidal representation $V$ of $G_d$,
\[
 \dim V=\left(\prod_{a=1}^{d-1}(q^a-1)\right)
                         \dim\operatorname{Wh}(V).
 \tag{FLC.18}
\]
In particular, a cuspidal of degree
$\prod_{a=1}^{d-1}(q^a-1)$ has Whittaker multiplicity one.
\begin{proof}
For $1\le m\le d$, use exactly the tail subgroup $T_m$ and
character $\beta_m$ from FLC11, with ambient rank $d$,
and put $W_m=V^{T_m,\beta_m}$.
Then $W_d=V$, $W_1=\operatorname{Wh}(V)$, and $W_m$ is a
$P_m$-module.  For $m\ge2$, let $V_m$ be the last-column
normal subgroup of $P_m$.
The zero-character eigenspace $W_m^{V_m}$ vanishes: its
vectors are fixed by the entire unipotent radical of the
parabolic with blocks $(m-1,d-m+1)$.
Indeed the roots from the first $m-1$ rows into column $m-1$
belong to $V_m$, and those into later columns belong to $T_m$
and have trivial $\beta_m$; together they are that radical.
Cuspidality kills its invariant subspace.

All nonzero $V_m$-character eigenspaces in $W_m$ have equal
dimension, because $G_{m-1}$ acts transitively on their
characters.  The eigenspace for the last-coordinate character
is exactly $W_{m-1}$: adjoining $V_m$ to $T_m$ gives
$T_{m-1}$, and adds the single adjacent-entry term to its
character.  There are $q^{m-1}-1$ nonzero characters.  Hence
$\dim W_m=(q^{m-1}-1)\dim W_{m-1}$.
Multiplying these identities proves (FLC.18), including
$d=1$, when the product is empty.
\end{proof}
% END FLC12

% BEGIN FLC13
\paragraph{FLC13. Cuspidal support and the squarefree induction.}
\label{claim:FLC13}
Let $(\rho_i)_{i=1}^r$ and $(\rho'_j)_{j=1}^{r'}$ be lists
of cuspidal irreducibles, of positive ranks whose sums equal
$n$.  Their Harish--Chandra inductions satisfy
\[
 \dim\operatorname{Hom}_{G_n}
       (\operatorname{HC}(\rho_1,\ldots,\rho_r),
        \operatorname{HC}(\rho'_1,\ldots,\rho'_{r'}))
   =\#\{\sigma:\{1,\ldots,r\}\xrightarrow{\sim}
                       \{1,\ldots,r'\}:\rho_i\simeq\rho'_{\sigma(i)}\}.
 \tag{FLC.19}
\]
An isomorphism here includes equality of ranks.
Every irreducible of $G_n$ occurs in such an induction, and
the unordered list of cuspidal factors is uniquely determined.
If the factors are pairwise nonisomorphic, their induction
is irreducible.
\begin{proof}
For block sizes $(d_i)$ and $(e_j)$, the double cosets of the
two block parabolics are indexed by matrices of nonnegative
integers $(a_{ij})$ with row sums $d_i$ and column sums $e_j$.
One sees this directly from two flags: $a_{ij}$ is the dimension
of the successive quotient of their intersections.
Choosing bases for these quotients gives a simultaneous
adapted basis, and hence a representative; two pairs with
the same dimensions are carried to one another by a linear
isomorphism.  This proves both existence and completeness of
the double-coset indexing.

Apply FLC6.  If row $i$ has two or more nonzero entries, the
intersection subgroup on the $i$th Levi block includes the
unipotent radical of the proper flag with block sizes the
nonzero $a_{ij}$.  This subgroup acts trivially on the other
inducing representation, because there it lies in the
parabolic radical; its matrices are the identity on every
other block of the first Levi.  An intertwiner must consequently factor
through the invariants of $\rho_i$ under this proper
radical: average over that subgroup to see this explicitly.
Those invariants vanish by cuspidality.  The same argument,
with the flags interchanged, applies to a split column.
Only matrices having one nonzero entry in each row and each
column survive.  They are exactly rank-matching bijections
$\sigma$, and their intertwining spaces are
$\bigotimes_i\operatorname{Hom}_{G_{d_i}}(\rho_i,\rho'_{\sigma(i)})$.
Schur's lemma makes their dimensions zero or one, proving
(FLC.19).  Schur's lemma itself follows because the kernel
and image of an intertwiner are invariant, and any complex
endomorphism has an eigenvalue.

For existence of a cuspidal support, start with an irreducible
$V$.  If it is not cuspidal, some proper two-block radical
has nonzero invariants.  Choose an irreducible constituent
of this Levi module.  Irreducibles of a product are tensor
products, as follows by decomposing under one factor and
using its commuting multiplicity spaces.  Reciprocity embeds
$V$ as a constituent of the corresponding induction.
Repeat in smaller positive ranks until all factors are
cuspidal.  Induction in stages, proved by composing the
function models, gives one induction from the resulting list.
If two such inductions share $V$, semisimplicity gives a
nonzero map between them, so (FLC.19) identifies their lists.
Finally a list of pairwise nonisomorphic factors has just
the identity self-matching; its induced module has
endomorphism dimension one and therefore is irreducible.
The empty list gives the trivial rank-zero representation.
\end{proof}
% END FLC13

% BEGIN FLC14
\paragraph{FLC14. Whittaker multiplicity is multiplicative in finite induction.}
\label{claim:FLC14}
For representations $V_i$ of $G_{d_i}$,
\[
 \dim\operatorname{Wh}(\operatorname{HC}(V_1,\ldots,V_r))
                    =\prod_{i=1}^r\dim\operatorname{Wh}(V_i).
 \tag{FLC.20}
\]
\begin{proof}
The double cosets $U_n\backslash G_n/P$ are represented by
shuffles of the ordered Levi blocks.  To verify this, Gaussian
elimination gives $G_n=\coprod_w B_n wP$, with $w$ preserving
the order within each block.  Since diagonal matrices can
be moved through $w$ into $P$, the same representatives give
$U_n\backslash G_n/P$.
Let $b_a$ be the block occupying output coordinate $a$ of
such a shuffle.  If $b_a<b_{a+1}$, the elementary root subgroup
$1+tE_{a,a+1}$ lies in the conjugate parabolic radical.
It acts trivially on the inducing module but nontrivially
through $\lambda_n$.  Its term in reciprocity is therefore
zero.  The only shuffle without such an increase has the
blocks in decreasing order, with each block contiguous.
For this shuffle, the intersection with $U_n$ has the
within-block groups $U_{d_i}$, with their original adjacent
characters; every remaining radical coordinate has trivial
character.  Its intertwining space is the tensor product
of the $\operatorname{Wh}(V_i)$.  FLC6 proves (FLC.20).
This also proves the formula for character traces and all
parenthesizations of \texttt{hc\allowbreak L}, by induction in stages.
\end{proof}
% END FLC14

% BEGIN FLC15
\paragraph{FLC15. Elliptic character values uniquely determine the prime label.}
\label{claim:FLC15}
Fix $q\ge5$ and the original \texttt{Label\allowbreak Datum} $L_0$:
its injective multiplicative Teichmuller map is $\iota_0$,
and its degree-$d$ generator is $\gamma_d$, of order $q^d-1$.
For a root $x$ of a prime label of degree $d$, let
\[
 \theta_x(\gamma_d^a)=\iota_0(x)^a.
 \tag{FLC.21}
\]
If the same cuspidal character satisfies the original
\texttt{Is\allowbreak Cusp\allowbreak Of} formula for two prime labels $P,P'$ of
degree $d$, then $P=P'$.
\begin{proof}
The expression (FLC.21) is independent of $a$ modulo
$q^d-1$ and is exactly the source's \texttt{theta\allowbreak Of}.
Choose roots $x$ and $x'$ of the two labels.  On
$\mathbb F_{q^d}^{\times}$ the difference
\[
 f(y)=\sum_{i=0}^{d-1}\theta_x(y^{q^i})
              -\sum_{i=0}^{d-1}\theta_{x'}(y^{q^i})
\]
vanishes whenever $y$ has degree $d$ over $\mathbb F_q$.
Indeed multiplication by such a $y$ on
$\mathbb F_{q^d}$ is a matrix with irreducible characteristic
polynomial, and its elliptic character value is the same
under the two given formulas.  The common sign
$(-1)^{d-1}$ is nonzero and cancels.

Let $B$ be the number of nonzero elements of degree less than
$d$ in $\mathbb F_{q^d}$.  Then
\[
                         2dB<q^d-1.                 \tag{FLC.22}
\]
For $d=1$, $B=0$.  For $d=2$, $B=q-1$, so the assertion is
$4<q+1$.  For $d\ge3$ put $h=\lfloor d/2\rfloor$.
Every proper subfield has degree a proper divisor of $d$,
so $B<q^{h+1}/(q-1)$ by a geometric sum.
The inequality $q^{\lceil d/2\rceil-1}\ge d$ follows for
$d=3,4$ from $q\ge5$, and then follows for $d+2$ by
multiplying by $q\ge5$.  Therefore
$2dB<2q^d/(q-1)\le q^d/2<q^d-1$, proving (FLC.22).

A nonzero function on a finite abelian group of order $M$
whose Fourier expansion uses at most $r$ characters has
support of cardinality at least $M/r$.  To prove this,
write $f=\sum c_\eta\eta$.  Orthogonality gives
$\sum_y|f(y)|^2=M\sum_\eta|c_\eta|^2$, whereas
$|f(y)|^2\le r\sum_\eta|c_\eta|^2$ by Cauchy--Schwarz.
Summing only over the support and canceling proves the claim.
Here $M=q^d-1$ and $r\le2d$, so (FLC.22) forces $f=0$.
Character orthogonality now identifies its two multisets of
characters.  Thus $\theta_{x'}=\theta_{x^{q^i}}$ for some $i$.
Evaluating at $\gamma_d$ and using injectivity of $\iota_0$
on nonzero elements gives $x'=x^{q^i}$.  Their monic minimal
polynomials are equal, proving $P=P'$.
\end{proof}
% END FLC15

\paragraph{The one published finite-group input.}
We use Green's cuspidal theorem in the following precise form.
For every finite field $\mathbb F_q$ and $d\ge1$, the regular
characters $\theta:\mathbb F_{q^d}^{\times}\to\mathbb C^\times$,
modulo Frobenius, parametrize the cuspidal irreducibles of
$G_d$.  Regular means that the Frobenius orbit has length $d$.
Their degrees and elliptic values are
\[
 \dim\rho_\theta=\prod_{a=1}^{d-1}(q^a-1),\qquad
 \chi_{\rho_\theta}(g)=(-1)^{d-1}
                     \sum_{i=0}^{d-1}\theta(y^{q^i})
 \tag{FLC.23}
\]
when $g$ has irreducible characteristic polynomial with root
$y\in\mathbb F_{q^d}$.
The published input is
J.~A.~Green, \emph{The characters of the finite general linear
groups}, Trans. Amer. Math. Soc. \textbf{80} (1955), 402--447:
Theorems 12--13, Lemma 7.4, and Theorem 9.
\footnote{\url{https://www.ams.org/tran/1955-080-02/S0002-9947-1955-0072878-2/S0002-9947-1955-0072878-2.pdf}.}
Its exact dictionary is as follows.  Writing $\theta=\theta_0^a$
for a primitive character, the regular orbit is Green's
degree-$d$ simplex $g=\{a,aq,\ldots,aq^{d-1}\}$ modulo $q^d-1$.
Theorem 12 gives
\[
 \rho_\theta=(g^{(1)})=(-1)^{d-1}I_d^a[(1)].
\]
The paper's $\circ$-product is our finite Harish--Chandra induction.
Its Section 7 Schur-function homomorphism preserves simplex
support, with $s_{(1)}^v=\sum_{\lambda\vdash v}f^\lambda s_\lambda$
and $f^\lambda>0$.  Thus $(g^{(1)})$ cannot occur in proper
induction; every other symbol does, by factoring distinct
supports or this identity.  Reciprocity and Theorem 13 identify
exactly the cuspidals.  Definition 3.1 and (12), page 417,
with Theorem 9, give
\[
 I_d^a[(1)](g_y)=T_{d,1}(a;y)
                    =\sum_{i=0}^{d-1}\theta_0^{a q^i}(y).
\]
Lemma 7.4 gives
$I_d^a[(1)](1)=\prod_{j=1}^{d-1}(1-q^j)$.
Both formulas, with the displayed sign, are (FLC.23).
This applies Green's classification; it does not reprove it.

% BEGIN FLC16
\paragraph{FLC16. Green's theorem gives the literal \texttt{Green} record.}
\label{claim:FLC16}
For every prime $q$, every original \texttt{Label\allowbreak Datum} $L_0$,
and every $d\ge1$, there is a bijection
\[
 \{P\text{ monic irreducible}:\deg P=d,\ P(0)\ne0\}
       \longleftrightarrow
       \{\text{cuspidal irreducible characters of }G_d\}
 \tag{FLC.24}
\]
whose character $\tau_P$ satisfies \texttt{Is\allowbreak Cusp\allowbreak Of} for that
same $L_0$ and the degree formula in (FLC.23).
This supplies every field of the original \texttt{Green}
predicate using the stated published theorem.
\begin{proof}
Work in the unique degree-$d$ subfield of the fixed algebraic
closure.  The restriction of $\iota_0$ to its multiplicative
group is a bijection onto the complex roots of unity of
order dividing $q^d-1$: it is injective and both finite sets
have $q^d-1$ elements.  Since $\gamma_d$ is a generator,
(FLC.21) therefore identifies elements $x$ of this group
with all its complex characters.  It intertwines Frobenius
on $x$ with Frobenius on characters.  Characters with orbit
length $d$ correspond exactly to elements of degree $d$,
and their Frobenius orbits correspond exactly to the monic
irreducible polynomials of degree $d$ with nonzero constant
term.  Apply (FLC.23) to these characters.

Choosing any root of $P$ gives the same cuspidal character,
because the published parameter is a Frobenius orbit.
Every root of $P$ is a Frobenius conjugate, so the elliptic
formula holds for each root, exactly as required by
\texttt{Is\allowbreak Cusp\allowbreak Of}.  Roots of an irreducible characteristic
polynomial of a rank-$d$ matrix also lie in this same finite
subfield, so its quantified elliptic arguments are unchanged.
Formula (FLC.21) is the literal \texttt{theta\allowbreak Of}; independence
of the chosen exponent was proved in FLC15's first sentence
and does not require $q\ge5$.
The published injectivity and surjectivity on Frobenius
orbits give the remaining two fields of \texttt{Green}.
The supplied $\gamma_d$, $\iota_0$, and their norm compatibility
are retained.  No replacement compatible generator system
or new embedding is selected.
\end{proof}
% END FLC16

% BEGIN FLC17
\paragraph{FLC17. The exact squarefree $e_U$ law in character labels.}
\label{claim:FLC17}
Let $q\ge5$ be prime and use any original \texttt{Label\allowbreak Datum}.
For a monic squarefree degree-$N$ polynomial $F$ with
$F(0)\ne0$, the original \texttt{e\allowbreak Ser\allowbreak C} series consists of
one irreducible $\tau_F$, and
\[
 \dim\operatorname{Wh}(\tau_F)=1,\qquad
 \dim\tau_F=
       \frac{\prod_{a=1}^{N}(q^a-1)}
                         {\prod_{P\mid F}(q^{\deg P}-1)},\qquad
 w_{e_U}(F)=\prod_{P\mid F}(q^{\deg P}-1)^{-1}.
 \tag{FLC.25}
\]
This is the full original squarefree Gelfand--Graev weight law.
\begin{proof}
FLC16 assigns a distinct cuspidal $\rho_P$ to every distinct
factor $P\mid F$.  FLC15 says that the source's
\texttt{Is\allowbreak Cusp\allowbreak Of} relation has no other prime label for a
cuspidal.  Thus FLC13's support uniqueness identifies the
original \texttt{Has\allowbreak Label} series of $F$ with the constituents
of $\operatorname{HC}((\rho_P)_{P\mid F})$.
The factors are pairwise nonisomorphic, so FLC13 makes this
induction irreducible, proving the singleton assertion.
Its Whittaker multiplicity is one by FLC12, FLC14, and the
published cuspidal degree supplied in FLC16.

For completeness, its degree uses only the index of the
actual block parabolic.  A parabolic with blocks of sizes
$d_1,\ldots,d_r$ has order
$q^{\sum_{i<j}d_id_j}\prod_i|G_{d_i}|$.
Counting ordered bases gives
$|G_n|=q^{n(n-1)/2}\prod_{a=1}^n(q^a-1)$.
Multiplying the parabolic index by
$\prod_i\prod_{a=1}^{d_i-1}(q^a-1)$ cancels the powers of
$q$ and gives the displayed degree.  Finally RRT3 gives
\[
 w_{e_U}(F)=\frac{\dim\tau_F}{|G_N|[1]e_U}
          =\frac{\dim\tau_F}{\prod_{a=1}^{N}(q^a-1)},
\]
since $\dim e_U\tau_F=1$ and $[1]e_U=q^{-N(N-1)/2}$.
This proves the exact normalization in (FLC.25).
For $N=0$, the empty support and all empty products equal one,
which is also the original rank-zero statement.
\end{proof}
% END FLC17

In the original datum $D$, the stored type-prime congruence is
\[
D.\mathrm{host}.g(k)\equiv1\pmod{D.\mathrm{host}.QS\,
D.\mathrm{host}.q\,D.\mathrm{host}.d}.
\]
The stored tower divisibility makes this modulus divisible by
$4$.  Thus each original type prime is $1$ modulo $4$ and,
being prime, is at least $5$.  The preceding character-label
hypotheses therefore apply at both of the original type primes.

% BEGIN FLC18
\paragraph{FLC18. The original normalized-label comparison.}
\label{claim:FLC18}
\mbox{}\par\noindent
Fix an actual \texttt{Normalized\allowbreak Types\allowbreak Compatibility} $T$ for
the original datum $D$ and its specified embedding $j_E$.
At either type prime $q=D.\mathrm{host}.g(k)$, let $O=T.O(k)$
be its original \texttt{Label\allowbreak Map}.  Let $L_0$ use exactly
$D.\mathrm{ntypes}.\mathrm{teichC}(k)$ and $T.\zeta(k)$.
There is a degree-preserving permutation $\sigma$ of prime
polynomial labels such that, for every valid degree-$N$ label,
\[
 \widehat\sigma(F)=\prod_{P\mid F}\sigma(P)^{v_P(F)},\qquad
 E_F^{\mathrm{original}}=e\mathrm{SerC}_{L_0}(\widehat\sigma(F)).
 \tag{FLC.26}
\]
It preserves squarefreeness and the multiset of factor degrees.
Consequently the exact squarefree law (FLC.25) and lower
comparison (FLC.17) hold on the original labels themselves.
\begin{proof}
Fix $d\ge1$.  For each Green prime label $P$ choose a
representation $V_P$ of its cuspidal character from FLC16.
Send $P$ to $O.\mathrm{lab}(d,V_P)$.
The original \texttt{lab\_\allowbreak spec} makes this another prime
label of degree $d$.  This map is injective: equal $O$-labels
give isomorphic representations by the original
\texttt{lab\_\allowbreak inj}, hence equal Green characters and equal
$P$.  The finite set of degree-$d$ prime labels maps
injectively to itself, so this map is a permutation $\pi_d$.
Take $\sigma_d=\pi_d^{-1}$.

For any cuspidal representation $V$, FLC16 gives a $P$ with
its character.  Equal characters imply isomorphic complex
representations: averaging on $\operatorname{Hom}(V,V_P)$
computes its invariant dimension as their character pairing,
which is one by Schur's lemma and semisimplicity.
Thus \texttt{lab\_\allowbreak iso} gives
$O.\mathrm{lab}(d,V)=\pi_d(P)$ and
$\sigma_d(O.\mathrm{lab}(d,V))=P$.
This constructs the same degree-preserving prime-label
permutation required in the source comparison.
Its extension to products preserves degrees and
multiplicities; unique factorization gives
$\widehat\pi\widehat\sigma(F)=F$ and preservation of
squarefreeness.

The original \texttt{Has\allowbreak Label\allowbreak O} says exactly that an
irreducible is a constituent of a block-parabolic induction
of cuspidals and that $F$ is the product of their $O$-labels.
Transporting the original block indices to consecutive
indices conjugates the groups and representations; finite
induction is carried by the same change of basis, so its
character is the source's \texttt{hc\allowbreak L}.  No induction
normalization changes, since the group is finite.
FLC13's unique cuspidal support and the preceding prime-label
identity therefore give, for each irreducible $V$,
\[
 \mathrm{HasLabelO}(O,V,F)
 \quad\Longleftrightarrow\quad
 \mathrm{HasLabel}(L_0,\chi_V,\widehat\sigma(F)).
\]
FLC15 ensures that the right-hand label relation is the
same one used in the source.  Summing the same primitive
central idempotents proves (FLC.26), using the supplied
literal equality $T.\mathrm{eSer\_\allowbreak def}$.
In particular, this is an equality of the original group
algebra elements, not just of their ranks or eigenvalues.

The squarefree weight in (FLC.25) depends only on the factor
degrees, which $\widehat\sigma$ preserves.  The lower law
holds for every series and transfers through (FLC.26).
These comparisons are made at the original polynomial $F$;
no root-collision predicate is transported by an arbitrary
permutation of labels.  Subsequent polynomial estimates
therefore still use the actual roots of $F$.
\end{proof}
% END FLC18

% BEGIN FLC19
\paragraph{FLC19. The retained local weight producer on the original row object.}
\label{claim:FLC19}
Fix the original datum $D$, the original
\texttt{Normalized\allowbreak Types\allowbreak Compatibility} $T$ of FLC18, and an
actual \texttt{All\allowbreak Rows} object $A$ with
\texttt{All\allowbreak Join\allowbreak A A}.  Retain the original
\texttt{Tie\allowbreak Adm\allowbreak G} and \texttt{Range\allowbreak W} hypotheses.
Write $f=D.R.f$, $m=D.R.m$, $N=D.R.r$, and let $N_0$ be the
absolute threshold in RRT9--RRT10.  If $N\ge N_0$, then the
actual assertion \texttt{Weight\allowbreak All\allowbreak Rows A} holds, with its
unchanged constant
\begin{align}
 W_g={}&\frac{3f}{q}+\frac{23f\log q}{\sqrt q}
       +2fm\left(\frac1{4\min(q_1,d_0)}
                     +\frac{1+\log N}{q-1}\right)\notag\\
     &+\frac N2(7/8)^f+\exp\{-f/(5\log N)\},
       \qquad q=D.\mathrm{host}.g(k).
 \tag{FLC.27}
\end{align}
Here $(q_1,d_0)=(D.\mathrm{host}.q,D.\mathrm{host}.d)$.
The conclusion contains all local nonnegativity assertions
and the weighted bound for every finite set of bad tuples,
on the original $A.eVec$ and original \texttt{good\allowbreak At}.
For the original selected eventual family, the additional
$N_0$ guard is absorbed by its actual rank growth, so the
original eventual retained-defect application has no
remaining finite local quantitative input beyond the stated
published theorem and the original supplied normalized-type
record.
\begin{proof}
The exact guards supplied by the two original records include
$q\ge2^{12}$, distinct odd primes $q_1,d_0$ with
$\max(q_1,d_0)\ge16$, and $16,q_1,d_0\mid q-1$.
The original \texttt{All\allowbreak Join\allowbreak A.\allowbreak level} gives, at each type
prime, one nontrivial additive character and the literal
$e_J$ at every place.  More generally the argument permits
the exact \texttt{Tie\allowbreak Vec\allowbreak W} alternatives $e_U$ or $e_J$.
RRT1--RRT2 give their nonnegative normalized probability
laws on the actual degree-$N$ labels, with weight zero on
invalid labels.  FLC18 supplies the squarefree $e_U$ formula
on those original labels, and FLC11 supplies the required
lower $e_J$ comparison.

Apply RRT11, replacing its last collision input by FLC5.
Its five bad-event mechanisms now all have proofs:
RRT6 handles nonsquarefree factors; FLC5 handles failure of
root injectivity; RRT16 handles each of the $2fm$ root-to-twist
collisions; RRT8 handles a common proper divisor degree;
and RRT10 handles absence of the long prime-degree witness.
RRT4 transfers the resulting bound to the actual local
projectors at a cost at most $f/q$.  This is exactly (FLC.27).
The original factor count $|D.\mathrm{nd}.\mathrm{PlG}(k)|=f$
and twist count $m$ give its coefficients.
Any finite set of bad tuples is contained in the full bad
event and has no larger weight.  Valid polynomials are
identified injectively by their coefficients; invalid ones
have zero projector.  These facts give precisely the two
clauses of \texttt{Weight\allowbreak All\allowbreak Rows}, even when different global
rows have the same local label.

The original consumer \path{ConcNWeight.locShare_le_Wg}
uses only the first three \texttt{adm\allowbreak Lq} clauses.
We have proved those clauses directly for the actual
projectors.  The fourth clause, the upper $e_J$ comparison,
is unused in this argument.

For the actual selected family, the original rank record
has $f=2\operatorname{ord}\le k\le r=N$, and its size
choice has $x^2\le f$.  Thus $N\ge N_0$ whenever
$x\ge\sqrt{N_0}$, uniformly in every later local choice.
This is the same legitimate rank comparison written in
RRT11.  It does not follow from the abstract upper/lower
bounds on $N$ in \texttt{Sizes\allowbreak At} alone.
\end{proof}
% END FLC19

\subsection{Exact scope of the finite completion}

FLC1--FLC5 close the full original \texttt{Same\allowbreak Factor\allowbreak Count},
\texttt{Cross\allowbreak Factor\allowbreak Count}, and \texttt{Local\allowbreak Estimate\allowbreak II}.
FLC6--FLC11 close the actual pattern exchange and the full
lower local character comparison.  FLC12--FLC18 give the
squarefree series law and the degree-preserving normalized
label comparison, using the explicitly stated published
cuspidal theorem.  Every polynomial collision condition is
then tested on the original label, not on a relabeled root
set.  FLC19 supplies the actual retained local weight clause
for the selected eventual family with the original $W_g$.

This chapter retains the original supplied normalized-type
compatibility record; it does not construct that record from
bare abstract projectors.  It does not produce the original
adelic \texttt{Hida\allowbreak Setting}, the actual reduction and Haar
normalization in \texttt{Et\allowbreak One}, the original
\texttt{Rows\allowbreak Add} and \texttt{Rows\allowbreak Good} identities, the common
\texttt{All\allowbreak Rows} object, or any of the three sufficient
arithmetic cylinder counts in GQT5.  The full fourth
\texttt{adm\allowbreak Lq} clause, an upper comparison unused by this
consumer, is not claimed here.  The long-prime-degree term
still uses RRT10's eventual theorem, with its explicit
actual-rank justification.  These scope statements prevent
a local completion from being mistaken for completion of
\texttt{GLOBAL\_\allowbreak RETAINED\_\allowbreak DEFECT}, RH, or GRH.

\endgroup
\end{document}
